\subsection{Problema manufacturado bidimensional}
\label{sec:es-results-2d}
Las tablas de esta subsección se generan directamente desde el árbol de
resultados \path{/home/pablo/OpenFOAM/pablo-v2312/run/tutorials_electro/highOrderManufacturedFDAImplicitPETSc/results/example0}. Las raíces por método son \path{/home/pablo/OpenFOAM/pablo-v2312/run/tutorials_electro/highOrderManufacturedFDAImplicitPETSc/results/example0/2D_CN_cons_Picard_RKF45}, \path{/home/pablo/OpenFOAM/pablo-v2312/run/tutorials_electro/highOrderManufacturedFDAImplicitPETSc/results/example0/2D_CN_cons_diagonalIion_RKF45}, \path{/home/pablo/OpenFOAM/pablo-v2312/run/tutorials_electro/highOrderManufacturedFDAImplicitPETSc/results/example0/2D_CN_cons_JFNK_RKF45}.
Todas las entradas usan \code{crankNicolson}, la matriz de masa
\code{consistent} y \code{RKF45} para las ODEs de estado.
La primera columna reporta \code{Vm-states-Iion}: el grado de reconstrucción
de $V_m$, el tratamiento de estados (\code{CCp*}, \code{GP} o \code{na}) y el
ajuste de cuadratura/reconstrucción de $\Iion$.
Los ajustes usan $N=10$--$100$ para el rango amplio de mallas y
$N=70$--$100$ para el rango de mallas finas.
Las dos últimas columnas auditan el solver no lineal: \code{RB} cuenta los
pasos revertidos (rollback) y $n_{\mathrm{nl}}$ resume las iteraciones no
lineales por paso (mediana/máximo).

\subsubsection{\code{Picard}}
\paragraph{$\Delta t=10^{-2}$}
\tiny
\setlength{\tabcolsep}{1.4pt}
\begin{longtable}{lllrrrrrrrrrrrrrrrr}
\caption{Resultados manufacturados 2D para \code{Picard}, \code{crankNicolson}, masa \code{consistent}, estados \code{RKF45} y $\Delta t=10^{-2}$. El primer bloque de convergencia se ajusta sobre $N=10$--$100$; el segundo sobre $N=70$--$100$. La primera columna registra la terna \code{Vm-states-Iion} efectivamente usada en la corrida. \textbf{RB} es la cantidad de pasos de tiempo que agotaron el presupuesto de iteraciones no lineales, acumulada sobre los 91 tamaños de malla de la fila (455 pasos de tiempo en total por fila). \code{Picard} acepta el último iterado con advertencia en vez de hacer rollback, de modo que para este método la columna cuenta pasos aceptados sin converger. \textbf{$n_{\mathrm{nl}}$} reporta la mediana y el máximo de iteraciones no lineales por paso sobre esos mismos pasos; el presupuesto configurado es \code{nonlinearIterations}$=2000$.}\label{tab:es-results-2d-picard-1p0em02}\\
\toprule
& & & \multicolumn{6}{c}{Ajuste $N=10$--$100$} & \multicolumn{6}{c}{Ajuste $N=70$--$100$} & & & & \\
\cmidrule(lr){4-9}\cmidrule(lr){10-15}
\code{Vm-states-Iion} & Tipo de malla & $\alpha$ & $V_m$ & $V_m^G$ & $u_1$ & $u_1^G$ & $u_2$ & $u_2^G$ & $V_m$ & $V_m^G$ & $u_1$ & $u_1^G$ & $u_2$ & $u_2^G$ & $t_{\Sigma}$ [s] & RSS$_{\Sigma}$ [MB] & RB & $n_{\mathrm{nl}}$ \\
\midrule
\endfirsthead
\toprule
& & & \multicolumn{6}{c}{Ajuste $N=10$--$100$} & \multicolumn{6}{c}{Ajuste $N=70$--$100$} & & & & \\
\cmidrule(lr){4-9}\cmidrule(lr){10-15}
\code{Vm-states-Iion} & Tipo de malla & $\alpha$ & $V_m$ & $V_m^G$ & $u_1$ & $u_1^G$ & $u_2$ & $u_2^G$ & $V_m$ & $V_m^G$ & $u_1$ & $u_1^G$ & $u_2$ & $u_2^G$ & $t_{\Sigma}$ [s] & RSS$_{\Sigma}$ [MB] & RB & $n_{\mathrm{nl}}$ \\
\midrule
\endhead
\code{NO--na--NO} & hexa & 0.0 & 1.992 & 1.992 & 2.135 & 2.135 & 2.134 & 2.134 & 1.985 & 1.985 & 2.500 & 2.500 & 2.497 & 2.497 & 5.143 & 16406.2 & 0 & 4.00/4 \\
\code{NO--na--NO} & hexa & 0.1 & 1.849 & 1.849 & 1.612 & 1.612 & 1.618 & 1.618 & 2.278 & 2.278 & 2.159 & 2.159 & 2.163 & 2.163 & 18.689 & 18355.6 & 0 & 4.00/4 \\
\code{NO--na--NO} & triang. no estr. & 0.0 & 0.730 & 0.730 & 0.778 & 0.778 & 0.784 & 0.784 & 0.862 & 0.862 & 1.037 & 1.037 & 1.107 & 1.107 & 66.651 & 17158.0 & 0 & 4.00/4 \\
\code{NO--na--NO} & triang. no estr. & 0.1 & 0.722 & 0.722 & 0.772 & 0.772 & 0.777 & 0.777 & 0.871 & 0.871 & 1.027 & 1.027 & 1.097 & 1.097 & 118.52 & 20786.3 & 0 & 4.00/4 \\
\hline
\code{NO--CCp1--p1} & hexa & 0.0 & 1.993 & 1.993 & 2.136 & 2.354 & 2.134 & 2.333 & 1.985 & 1.985 & 2.505 & 2.289 & 2.497 & 2.242 & 17.995 & 16908.0 & 0 & 4.00/4 \\
\code{NO--CCp1--p1} & hexa & 0.1 & 1.849 & 1.849 & 1.611 & 2.354 & 1.618 & 2.310 & 2.278 & 2.278 & 2.159 & 2.289 & 2.163 & 2.225 & 37.501 & 18456.9 & 0 & 4.00/4 \\
\code{NO--CCp1--p1} & triang. no estr. & 0.0 & 0.730 & 0.730 & 0.778 & 0.778 & 0.784 & 0.784 & 0.862 & 0.862 & 1.037 & 1.037 & 1.107 & 1.107 & 126.46 & 17892.6 & 0 & 4.00/4 \\
\code{NO--CCp1--p1} & triang. no estr. & 0.1 & 0.722 & 0.722 & 0.772 & 0.772 & 0.777 & 0.777 & 0.871 & 0.871 & 1.027 & 1.027 & 1.097 & 1.097 & 158.68 & 20895.8 & 0 & 4.00/4 \\
\hline
\code{NO--CCp1--p2} & hexa & 0.0 & 1.993 & 1.993 & 2.140 & 2.229 & 2.134 & 2.184 & 1.985 & 1.985 & 2.513 & 2.150 & 2.495 & 2.094 & 21.941 & 17118.3 & 0 & 4.00/4 \\
\code{NO--CCp1--p2} & hexa & 0.1 & 1.849 & 1.849 & 1.609 & 2.229 & 1.619 & 2.181 & 2.278 & 2.278 & 2.159 & 2.149 & 2.163 & 2.091 & 32.920 & 18586.4 & 0 & 4.00/4 \\
\code{NO--CCp1--p2} & triang. no estr. & 0.0 & 0.730 & 0.730 & 0.778 & 2.111 & 0.784 & 1.164 & 0.862 & 0.862 & 1.037 & 2.101 & 1.107 & 1.150 & 133.37 & 18167.6 & 0 & 4.00/4 \\
\code{NO--CCp1--p2} & triang. no estr. & 0.1 & 0.722 & 0.722 & 0.772 & 2.112 & 0.777 & 1.172 & 0.871 & 0.871 & 1.027 & 2.105 & 1.097 & 1.143 & 153.34 & 21019.5 & 0 & 4.00/4 \\
\hline
\code{NO--CCp1--p3} & hexa & 0.0 & 1.993 & 1.993 & 2.140 & 2.247 & 2.134 & 2.211 & 1.985 & 1.985 & 2.513 & 2.167 & 2.495 & 2.112 & 20.564 & 17226.5 & 0 & 4.00/4 \\
\code{NO--CCp1--p3} & hexa & 0.1 & 1.849 & 1.849 & 1.609 & 2.247 & 1.619 & 2.206 & 2.278 & 2.278 & 2.159 & 2.166 & 2.163 & 2.109 & 25.935 & 18641.4 & 0 & 4.00/4 \\
\code{NO--CCp1--p3} & triang. no estr. & 0.0 & 0.730 & 0.730 & 0.778 & 2.139 & 0.784 & 1.132 & 0.862 & 0.862 & 1.037 & 2.131 & 1.107 & 1.144 & 115.55 & 18286.4 & 0 & 4.00/4 \\
\code{NO--CCp1--p3} & triang. no estr. & 0.1 & 0.722 & 0.722 & 0.772 & 2.140 & 0.777 & 1.139 & 0.871 & 0.871 & 1.027 & 2.135 & 1.097 & 1.136 & 153.47 & 21092.0 & 0 & 4.00/4 \\
\hline
\code{NO--CCp2--p1} & hexa & 0.0 & 1.992 & 1.992 & 2.141 & 3.042 & 2.132 & 2.624 & 1.985 & 1.985 & 2.518 & 2.999 & 2.491 & 2.604 & 25.070 & 17413.8 & 0 & 4.00/4 \\
\code{NO--CCp2--p1} & hexa & 0.1 & 1.849 & 1.849 & 1.607 & 3.038 & 1.618 & 1.948 & 2.278 & 2.278 & 2.158 & 2.987 & 2.162 & 2.187 & 42.260 & 18557.9 & 0 & 4.00/4 \\
\code{NO--CCp2--p1} & triang. no estr. & 0.0 & 0.730 & 0.730 & 0.778 & 0.778 & 0.784 & 0.784 & 0.862 & 0.862 & 1.037 & 1.037 & 1.107 & 1.107 & 120.75 & 19093.8 & 0 & 4.00/4 \\
\code{NO--CCp2--p1} & triang. no estr. & 0.1 & 0.722 & 0.722 & 0.772 & 0.772 & 0.777 & 0.777 & 0.871 & 0.871 & 1.027 & 1.027 & 1.097 & 1.097 & 180.69 & 20978.0 & 0 & 4.00/4 \\
\hline
\code{NO--CCp2--p2} & hexa & 0.0 & 1.991 & 1.991 & 2.145 & 3.063 & 2.121 & 2.765 & 1.986 & 1.986 & 2.519 & 3.002 & 2.444 & 2.630 & 29.600 & 17632.8 & 0 & 4.00/4 \\
\code{NO--CCp2--p2} & hexa & 0.1 & 1.851 & 1.851 & 1.600 & 3.061 & 1.624 & 2.068 & 2.274 & 2.274 & 2.155 & 2.996 & 2.162 & 2.210 & 30.224 & 18673.4 & 0 & 4.00/4 \\
\code{NO--CCp2--p2} & triang. no estr. & 0.0 & 0.730 & 0.730 & 0.778 & 1.941 & 0.784 & 0.776 & 0.862 & 0.862 & 1.037 & 1.095 & 1.107 & 1.097 & 117.04 & 19335.1 & 0 & 4.00/4 \\
\code{NO--CCp2--p2} & triang. no estr. & 0.1 & 0.722 & 0.722 & 0.772 & 1.957 & 0.777 & 0.773 & 0.871 & 0.871 & 1.027 & 1.090 & 1.097 & 1.087 & 177.40 & 21128.6 & 0 & 4.00/4 \\
\hline
\code{NO--CCp2--p3} & hexa & 0.0 & 1.991 & 1.991 & 2.145 & 3.058 & 2.121 & 2.763 & 1.986 & 1.986 & 2.519 & 3.001 & 2.444 & 2.632 & 28.256 & 17729.6 & 0 & 4.00/4 \\
\code{NO--CCp2--p3} & hexa & 0.1 & 1.851 & 1.851 & 1.600 & 3.056 & 1.624 & 2.069 & 2.274 & 2.274 & 2.155 & 2.995 & 2.162 & 2.211 & 42.963 & 18800.3 & 0 & 4.00/4 \\
\code{NO--CCp2--p3} & triang. no estr. & 0.0 & 0.730 & 0.730 & 0.778 & 1.952 & 0.784 & 0.776 & 0.862 & 0.862 & 1.037 & 1.098 & 1.107 & 1.096 & 112.55 & 19472.0 & 0 & 4.00/4 \\
\code{NO--CCp2--p3} & triang. no estr. & 0.1 & 0.722 & 0.722 & 0.772 & 1.968 & 0.777 & 0.773 & 0.871 & 0.871 & 1.027 & 1.093 & 1.097 & 1.087 & 173.29 & 21226.5 & 0 & 4.00/4 \\
\hline
\code{NO--CCp3--p1} & hexa & 0.0 & 1.992 & 1.992 & 2.139 & 4.063 & 2.132 & 2.165 & 1.985 & 1.985 & 2.518 & 3.183 & 2.491 & 2.491 & 40.748 & 18769.0 & 0 & 4.00/4 \\
\code{NO--CCp3--p1} & hexa & 0.1 & 1.849 & 1.849 & 1.605 & 3.118 & 1.618 & 1.661 & 2.278 & 2.278 & 2.158 & 2.340 & 2.163 & 2.189 & 52.484 & 19843.9 & 0 & 4.00/4 \\
\code{NO--CCp3--p1} & triang. no estr. & 0.0 & 0.730 & 0.730 & 0.778 & 0.778 & 0.784 & 0.784 & 0.862 & 0.862 & 1.037 & 1.037 & 1.107 & 1.107 & 131.85 & 21560.9 & 0 & 4.00/4 \\
\code{NO--CCp3--p1} & triang. no estr. & 0.1 & 0.722 & 0.722 & 0.772 & 0.772 & 0.777 & 0.777 & 0.871 & 0.871 & 1.027 & 1.027 & 1.097 & 1.097 & 207.53 & 23322.9 & 0 & 4.00/4 \\
\hline
\code{NO--CCp3--p2} & hexa & 0.0 & 1.992 & 1.992 & 2.141 & 4.277 & 2.120 & 2.217 & 1.985 & 1.985 & 2.519 & 4.233 & 2.444 & 2.444 & 34.667 & 18981.8 & 0 & 4.00/4 \\
\code{NO--CCp3--p2} & hexa & 0.1 & 1.851 & 1.851 & 1.594 & 3.421 & 1.622 & 1.722 & 2.274 & 2.274 & 2.155 & 2.589 & 2.162 & 2.221 & 54.473 & 20039.5 & 0 & 4.00/4 \\
\code{NO--CCp3--p2} & triang. no estr. & 0.0 & 0.730 & 0.730 & 0.778 & 1.071 & 0.784 & 0.796 & 0.862 & 0.862 & 1.037 & 1.037 & 1.107 & 1.109 & 165.34 & 21845.6 & 0 & 4.00/4 \\
\code{NO--CCp3--p2} & triang. no estr. & 0.1 & 0.722 & 0.722 & 0.772 & 1.081 & 0.777 & 0.798 & 0.871 & 0.871 & 1.027 & 1.028 & 1.097 & 1.100 & 208.25 & 23561.6 & 0 & 4.00/4 \\
\hline
\code{NO--CCp3--p3} & hexa & 0.0 & 1.992 & 1.992 & 2.141 & 4.293 & 2.120 & 2.214 & 1.985 & 1.985 & 2.519 & 4.254 & 2.444 & 2.444 & 47.414 & 19082.2 & 0 & 4.00/4 \\
\code{NO--CCp3--p3} & hexa & 0.1 & 1.851 & 1.851 & 1.594 & 3.409 & 1.622 & 1.721 & 2.274 & 2.274 & 2.155 & 2.578 & 2.162 & 2.221 & 56.993 & 20145.0 & 0 & 4.00/4 \\
\code{NO--CCp3--p3} & triang. no estr. & 0.0 & 0.730 & 0.730 & 0.778 & 1.054 & 0.784 & 0.796 & 0.862 & 0.862 & 1.037 & 1.037 & 1.107 & 1.109 & 180.29 & 21975.3 & 0 & 4.00/4 \\
\code{NO--CCp3--p3} & triang. no estr. & 0.1 & 0.722 & 0.722 & 0.772 & 1.064 & 0.777 & 0.797 & 0.871 & 0.871 & 1.027 & 1.028 & 1.097 & 1.100 & 205.65 & 23684.7 & 0 & 4.00/4 \\
\hline
\code{NO--GP--p1} & hexa & 0.0 & 1.992 & 1.992 & 2.000 & 2.057 & 2.005 & 2.061 & 1.985 & 1.985 & 2.002 & 2.031 & 2.023 & 2.031 & 6.823 & 16540.8 & 0 & 4.00/4 \\
\code{NO--GP--p1} & hexa & 0.1 & 1.849 & 1.849 & 2.000 & 2.049 & 1.997 & 2.053 & 2.278 & 2.278 & 2.002 & 2.039 & 2.030 & 2.039 & 21.520 & 18420.0 & 0 & 4.00/4 \\
\code{NO--GP--p1} & triang. no estr. & 0.0 & 0.730 & 0.730 & 0.778 & 0.778 & 0.784 & 0.784 & 0.862 & 0.862 & 1.037 & 1.037 & 1.107 & 1.107 & 72.359 & 17303.3 & 0 & 4.00/4 \\
\code{NO--GP--p1} & triang. no estr. & 0.1 & 0.722 & 0.722 & 0.772 & 0.772 & 0.777 & 0.777 & 0.871 & 0.871 & 1.027 & 1.027 & 1.097 & 1.097 & 128.31 & 20871.0 & 0 & 4.00/4 \\
\hline
\code{NO--GP--p2} & hexa & 0.0 & 1.992 & 1.992 & 1.999 & 2.015 & 2.002 & 2.017 & 1.986 & 1.986 & 2.001 & 2.010 & 2.011 & 2.010 & 14.302 & 16794.5 & 0 & 4.00/4 \\
\code{NO--GP--p2} & hexa & 0.1 & 1.851 & 1.851 & 1.999 & 2.015 & 2.001 & 2.017 & 2.274 & 2.274 & 2.001 & 2.011 & 2.013 & 2.011 & 26.791 & 18563.5 & 0 & 4.00/4 \\
\code{NO--GP--p2} & triang. no estr. & 0.0 & 0.730 & 0.730 & 1.878 & 1.418 & 1.102 & 1.411 & 0.862 & 0.862 & 1.678 & 1.219 & 1.112 & 1.269 & 71.948 & 17559.7 & 0 & 4.00/4 \\
\code{NO--GP--p2} & triang. no estr. & 0.1 & 0.722 & 0.722 & 1.880 & 1.427 & 1.107 & 1.420 & 0.871 & 0.871 & 1.682 & 1.220 & 1.104 & 1.270 & 132.64 & 21009.7 & 0 & 4.00/4 \\
\hline
\code{NO--GP--p3} & hexa & 0.0 & 1.992 & 1.992 & 1.999 & 2.019 & 2.002 & 2.021 & 1.986 & 1.986 & 2.001 & 2.011 & 2.011 & 2.011 & 34.842 & 16938.5 & 0 & 4.00/4 \\
\code{NO--GP--p3} & hexa & 0.1 & 1.851 & 1.851 & 1.999 & 2.019 & 2.001 & 2.021 & 2.274 & 2.274 & 2.001 & 2.013 & 2.013 & 2.013 & 36.836 & 18613.1 & 0 & 4.00/4 \\
\code{NO--GP--p3} & triang. no estr. & 0.0 & 0.730 & 0.730 & 1.878 & 1.389 & 1.102 & 1.384 & 0.862 & 0.862 & 1.678 & 1.206 & 1.112 & 1.258 & 78.186 & 17682.8 & 0 & 4.00/4 \\
\code{NO--GP--p3} & triang. no estr. & 0.1 & 0.722 & 0.722 & 1.881 & 1.398 & 1.107 & 1.392 & 0.871 & 0.871 & 1.682 & 1.207 & 1.104 & 1.259 & 129.33 & 21069.0 & 0 & 4.00/4 \\
\hline
\code{p1--na--NO} & hexa & 0.0 & 1.911 & 2.146 & 1.616 & 1.616 & 1.626 & 1.626 & 2.612 & 2.168 & 2.381 & 2.381 & 2.387 & 2.387 & 116.43 & 17750.5 & 0 & 4.00/4 \\
\code{p1--na--NO} & hexa & 0.1 & 1.971 & 2.163 & 1.677 & 1.677 & 1.687 & 1.687 & 2.614 & 2.141 & 2.416 & 2.416 & 2.422 & 2.422 & 93.581 & 19067.1 & 0 & 4.00/4 \\
\code{p1--na--NO} & triang. no estr. & 0.0 & 2.271 & 2.126 & 2.004 & 2.004 & 2.002 & 2.002 & 2.384 & 2.115 & 2.552 & 2.552 & 2.549 & 2.549 & 566.18 & 19578.9 & 0 & 4.00/4 \\
\code{p1--na--NO} & triang. no estr. & 0.1 & 2.301 & 2.106 & 2.092 & 2.092 & 2.090 & 2.090 & 2.222 & 2.054 & 2.524 & 2.524 & 2.524 & 2.524 & 636.61 & 21518.4 & 0 & 4.00/4 \\
\hline
\code{p1--CCp1--p1} & hexa & 0.0 & 1.911 & 2.146 & 1.616 & 2.349 & 1.626 & 1.955 & 2.612 & 2.168 & 2.381 & 2.294 & 2.387 & 2.369 & 114.17 & 18231.1 & 0 & 4.00/4 \\
\code{p1--CCp1--p1} & hexa & 0.1 & 1.971 & 2.163 & 1.677 & 2.351 & 1.687 & 2.045 & 2.614 & 2.141 & 2.416 & 2.293 & 2.422 & 2.383 & 93.147 & 19257.6 & 0 & 4.00/4 \\
\code{p1--CCp1--p1} & triang. no estr. & 0.0 & 2.271 & 2.126 & 2.004 & 2.004 & 2.002 & 2.002 & 2.384 & 2.115 & 2.552 & 2.552 & 2.549 & 2.549 & 591.55 & 20179.6 & 0 & 4.00/4 \\
\code{p1--CCp1--p1} & triang. no estr. & 0.1 & 2.301 & 2.106 & 2.092 & 2.092 & 2.090 & 2.090 & 2.222 & 2.054 & 2.524 & 2.524 & 2.524 & 2.524 & 609.37 & 21594.2 & 0 & 4.00/4 \\
\hline
\code{p1--CCp1--p2} & hexa & 0.0 & 1.911 & 2.146 & 1.616 & 2.228 & 1.626 & 2.057 & 2.612 & 2.168 & 2.381 & 2.153 & 2.387 & 2.232 & 116.12 & 18427.4 & 0 & 4.00/4 \\
\code{p1--CCp1--p2} & hexa & 0.1 & 1.971 & 2.163 & 1.676 & 2.228 & 1.687 & 2.104 & 2.614 & 2.141 & 2.416 & 2.152 & 2.422 & 2.208 & 104.61 & 19367.0 & 0 & 4.00/4 \\
\code{p1--CCp1--p2} & triang. no estr. & 0.0 & 2.271 & 2.126 & 2.004 & 2.102 & 2.002 & 2.036 & 2.384 & 2.115 & 2.552 & 2.063 & 2.549 & 2.285 & 668.38 & 20456.4 & 0 & 4.00/4 \\
\code{p1--CCp1--p2} & triang. no estr. & 0.1 & 2.302 & 2.106 & 2.092 & 2.102 & 2.090 & 2.066 & 2.222 & 2.054 & 2.524 & 2.061 & 2.524 & 2.190 & 625.36 & 21732.8 & 0 & 4.00/4 \\
\hline
\code{p1--CCp1--p3} & hexa & 0.0 & 1.911 & 2.146 & 1.616 & 2.246 & 1.626 & 2.063 & 2.612 & 2.168 & 2.381 & 2.170 & 2.387 & 2.255 & 83.912 & 18517.0 & 0 & 4.00/4 \\
\code{p1--CCp1--p3} & hexa & 0.1 & 1.971 & 2.163 & 1.676 & 2.246 & 1.687 & 2.115 & 2.614 & 2.141 & 2.416 & 2.169 & 2.422 & 2.234 & 97.852 & 19438.2 & 0 & 4.00/4 \\
\code{p1--CCp1--p3} & triang. no estr. & 0.0 & 2.271 & 2.126 & 2.004 & 2.127 & 2.002 & 2.040 & 2.384 & 2.115 & 2.552 & 2.080 & 2.549 & 2.315 & 725.98 & 20581.6 & 0 & 4.00/4 \\
\code{p1--CCp1--p3} & triang. no estr. & 0.1 & 2.302 & 2.106 & 2.092 & 2.127 & 2.090 & 2.075 & 2.222 & 2.054 & 2.524 & 2.078 & 2.524 & 2.217 & 606.33 & 21809.8 & 0 & 4.00/4 \\
\hline
\code{p1--CCp2--p1} & hexa & 0.0 & 1.911 & 2.146 & 1.616 & 2.904 & 1.626 & 1.665 & 2.612 & 2.168 & 2.381 & 2.783 & 2.387 & 2.395 & 91.324 & 18731.6 & 0 & 4.00/4 \\
\code{p1--CCp2--p1} & hexa & 0.1 & 1.971 & 2.163 & 1.676 & 2.954 & 1.687 & 1.736 & 2.614 & 2.141 & 2.416 & 2.859 & 2.422 & 2.431 & 109.98 & 19338.5 & 0 & 4.00/4 \\
\code{p1--CCp2--p1} & triang. no estr. & 0.0 & 2.271 & 2.126 & 2.004 & 2.004 & 2.002 & 2.002 & 2.384 & 2.115 & 2.552 & 2.552 & 2.549 & 2.549 & 797.98 & 21404.6 & 0 & 4.00/4 \\
\code{p1--CCp2--p1} & triang. no estr. & 0.1 & 2.301 & 2.106 & 2.092 & 2.092 & 2.090 & 2.090 & 2.222 & 2.054 & 2.524 & 2.524 & 2.524 & 2.524 & 616.74 & 21869.5 & 0 & 4.00/4 \\
\hline
\code{p1--CCp2--p2} & hexa & 0.0 & 1.912 & 2.146 & 1.616 & 2.990 & 1.627 & 1.688 & 2.612 & 2.168 & 2.381 & 2.882 & 2.387 & 2.400 & 98.916 & 18950.6 & 0 & 4.00/4 \\
\code{p1--CCp2--p2} & hexa & 0.1 & 1.971 & 2.163 & 1.676 & 3.018 & 1.687 & 1.765 & 2.613 & 2.141 & 2.416 & 2.931 & 2.422 & 2.437 & 105.47 & 19476.9 & 0 & 4.00/4 \\
\code{p1--CCp2--p2} & triang. no estr. & 0.0 & 2.271 & 2.126 & 2.004 & 2.677 & 2.002 & 1.977 & 2.384 & 2.115 & 2.552 & 2.669 & 2.549 & 2.528 & 768.73 & 21638.0 & 0 & 4.00/4 \\
\code{p1--CCp2--p2} & triang. no estr. & 0.1 & 2.301 & 2.106 & 2.092 & 2.799 & 2.090 & 2.062 & 2.222 & 2.054 & 2.524 & 2.735 & 2.523 & 2.511 & 678.26 & 22074.3 & 0 & 4.00/4 \\
\hline
\code{p1--CCp2--p3} & hexa & 0.0 & 1.912 & 2.146 & 1.616 & 2.987 & 1.627 & 1.688 & 2.612 & 2.168 & 2.381 & 2.884 & 2.387 & 2.400 & 102.86 & 19056.1 & 0 & 4.00/4 \\
\code{p1--CCp2--p3} & hexa & 0.1 & 1.971 & 2.163 & 1.676 & 3.015 & 1.687 & 1.765 & 2.613 & 2.141 & 2.416 & 2.932 & 2.422 & 2.437 & 100.90 & 19565.9 & 0 & 4.00/4 \\
\code{p1--CCp2--p3} & triang. no estr. & 0.0 & 2.271 & 2.126 & 2.004 & 2.681 & 2.002 & 1.977 & 2.384 & 2.115 & 2.552 & 2.672 & 2.549 & 2.527 & 746.78 & 21749.9 & 0 & 4.00/4 \\
\code{p1--CCp2--p3} & triang. no estr. & 0.1 & 2.301 & 2.106 & 2.092 & 2.802 & 2.090 & 2.062 & 2.222 & 2.054 & 2.524 & 2.740 & 2.523 & 2.510 & 743.82 & 22186.1 & 0 & 4.00/4 \\
\hline
\code{p1--CCp3--p1} & hexa & 0.0 & 1.911 & 2.146 & 1.616 & 2.034 & 1.626 & 1.630 & 2.612 & 2.168 & 2.381 & 2.378 & 2.387 & 2.407 & 125.49 & 20075.4 & 0 & 4.00/4 \\
\code{p1--CCp3--p1} & hexa & 0.1 & 1.971 & 2.163 & 1.677 & 2.165 & 1.687 & 1.692 & 2.614 & 2.141 & 2.416 & 2.409 & 2.422 & 2.443 & 124.88 & 20548.1 & 0 & 4.00/4 \\
\code{p1--CCp3--p1} & triang. no estr. & 0.0 & 2.271 & 2.126 & 2.004 & 2.004 & 2.002 & 2.002 & 2.384 & 2.115 & 2.552 & 2.552 & 2.549 & 2.549 & 680.89 & 23892.8 & 0 & 4.00/4 \\
\code{p1--CCp3--p1} & triang. no estr. & 0.1 & 2.301 & 2.106 & 2.092 & 2.092 & 2.090 & 2.090 & 2.222 & 2.054 & 2.524 & 2.524 & 2.524 & 2.524 & 773.90 & 24267.6 & 0 & 4.00/4 \\
\hline
\code{p1--CCp3--p2} & hexa & 0.0 & 1.912 & 2.146 & 1.617 & 2.239 & 1.627 & 1.634 & 2.612 & 2.168 & 2.381 & 2.377 & 2.387 & 2.426 & 131.82 & 20280.0 & 0 & 4.00/4 \\
\code{p1--CCp3--p2} & hexa & 0.1 & 1.971 & 2.163 & 1.677 & 2.394 & 1.688 & 1.698 & 2.613 & 2.141 & 2.416 & 2.402 & 2.422 & 2.464 & 138.65 & 20735.8 & 0 & 4.00/4 \\
\code{p1--CCp3--p2} & triang. no estr. & 0.0 & 2.271 & 2.126 & 2.004 & 2.031 & 2.002 & 1.964 & 2.384 & 2.115 & 2.552 & 2.527 & 2.549 & 2.528 & 642.44 & 24134.8 & 0 & 4.00/4 \\
\code{p1--CCp3--p2} & triang. no estr. & 0.1 & 2.301 & 2.106 & 2.092 & 2.141 & 2.090 & 2.052 & 2.222 & 2.054 & 2.524 & 2.511 & 2.523 & 2.516 & 762.00 & 24520.7 & 0 & 4.00/4 \\
\hline
\code{p1--CCp3--p3} & hexa & 0.0 & 1.912 & 2.146 & 1.617 & 2.225 & 1.627 & 1.634 & 2.612 & 2.168 & 2.381 & 2.376 & 2.387 & 2.426 & 131.02 & 20374.9 & 0 & 4.00/4 \\
\code{p1--CCp3--p3} & hexa & 0.1 & 1.971 & 2.163 & 1.677 & 2.379 & 1.688 & 1.697 & 2.613 & 2.141 & 2.416 & 2.402 & 2.422 & 2.463 & 141.29 & 20835.1 & 0 & 4.00/4 \\
\code{p1--CCp3--p3} & triang. no estr. & 0.0 & 2.271 & 2.126 & 2.004 & 2.027 & 2.002 & 1.964 & 2.384 & 2.115 & 2.552 & 2.527 & 2.549 & 2.528 & 639.63 & 24254.8 & 0 & 4.00/4 \\
\code{p1--CCp3--p3} & triang. no estr. & 0.1 & 2.301 & 2.106 & 2.092 & 2.136 & 2.090 & 2.052 & 2.222 & 2.054 & 2.524 & 2.512 & 2.523 & 2.516 & 799.58 & 24633.1 & 0 & 4.00/4 \\
\hline
\code{p1--GP--p1} & hexa & 0.0 & 1.911 & 2.146 & 1.998 & 2.154 & 1.838 & 2.160 & 2.612 & 2.168 & 2.007 & 2.316 & 2.258 & 2.316 & 81.360 & 17778.0 & 0 & 4.00/4 \\
\code{p1--GP--p1} & hexa & 0.1 & 1.971 & 2.163 & 1.999 & 2.221 & 1.898 & 2.227 & 2.614 & 2.141 & 2.005 & 2.319 & 2.234 & 2.318 & 86.674 & 19080.5 & 0 & 4.00/4 \\
\code{p1--GP--p1} & triang. no estr. & 0.0 & 2.271 & 2.126 & 2.004 & 2.004 & 2.002 & 2.002 & 2.384 & 2.115 & 2.552 & 2.552 & 2.549 & 2.549 & 713.46 & 19611.5 & 0 & 4.00/4 \\
\code{p1--GP--p1} & triang. no estr. & 0.1 & 2.301 & 2.106 & 2.092 & 2.092 & 2.090 & 2.090 & 2.222 & 2.054 & 2.524 & 2.524 & 2.524 & 2.524 & 644.65 & 21555.8 & 0 & 4.00/4 \\
\hline
\code{p1--GP--p2} & hexa & 0.0 & 1.912 & 2.146 & 1.999 & 2.182 & 1.951 & 2.181 & 2.612 & 2.168 & 2.002 & 2.161 & 2.105 & 2.157 & 87.844 & 17999.5 & 0 & 4.00/4 \\
\code{p1--GP--p2} & hexa & 0.1 & 1.971 & 2.163 & 1.999 & 2.201 & 1.975 & 2.199 & 2.613 & 2.141 & 2.001 & 2.156 & 2.080 & 2.151 & 94.921 & 19190.6 & 0 & 4.00/4 \\
\code{p1--GP--p2} & triang. no estr. & 0.0 & 2.271 & 2.126 & 1.992 & 2.069 & 2.012 & 2.068 & 2.384 & 2.115 & 2.004 & 2.184 & 2.284 & 2.176 & 648.18 & 19866.6 & 0 & 4.00/4 \\
\code{p1--GP--p2} & triang. no estr. & 0.1 & 2.301 & 2.106 & 1.992 & 2.089 & 2.038 & 2.087 & 2.222 & 2.054 & 2.002 & 2.117 & 2.186 & 2.113 & 717.87 & 21700.1 & 0 & 4.00/4 \\
\hline
\code{p1--GP--p3} & hexa & 0.0 & 1.912 & 2.146 & 1.999 & 2.196 & 1.951 & 2.195 & 2.612 & 2.168 & 2.002 & 2.177 & 2.105 & 2.173 & 90.087 & 18115.0 & 0 & 4.00/4 \\
\code{p1--GP--p3} & hexa & 0.1 & 1.971 & 2.163 & 1.999 & 2.217 & 1.975 & 2.215 & 2.613 & 2.141 & 2.001 & 2.172 & 2.080 & 2.167 & 100.92 & 19245.1 & 0 & 4.00/4 \\
\code{p1--GP--p3} & triang. no estr. & 0.0 & 2.271 & 2.126 & 1.992 & 2.082 & 2.012 & 2.081 & 2.384 & 2.115 & 2.004 & 2.219 & 2.284 & 2.210 & 575.95 & 19986.2 & 0 & 4.00/4 \\
\code{p1--GP--p3} & triang. no estr. & 0.1 & 2.301 & 2.106 & 1.992 & 2.109 & 2.038 & 2.106 & 2.222 & 2.054 & 2.002 & 2.144 & 2.186 & 2.139 & 677.12 & 21763.5 & 0 & 4.00/4 \\
\hline
\code{p2--na--NO} & hexa & 0.0 & 2.127 & 2.567 & 2.176 & 2.176 & 2.176 & 2.176 & 2.016 & 2.179 & 2.139 & 2.139 & 2.139 & 2.139 & 85.517 & 18507.1 & 0 & 4.00/4 \\
\code{p2--na--NO} & hexa & 0.1 & 2.128 & 2.607 & 2.180 & 2.180 & 2.181 & 2.181 & 2.018 & 2.209 & 2.153 & 2.153 & 2.153 & 2.153 & 87.941 & 20275.2 & 0 & 4.00/4 \\
\code{p2--na--NO} & triang. no estr. & 0.0 & 2.056 & 2.374 & 2.166 & 2.166 & 2.162 & 2.162 & 1.982 & 2.068 & 2.261 & 2.261 & 2.263 & 2.263 & 593.35 & 21544.9 & 0 & 4.00/4 \\
\code{p2--na--NO} & triang. no estr. & 0.1 & 2.043 & 2.402 & 2.162 & 2.162 & 2.158 & 2.158 & 1.980 & 2.081 & 2.286 & 2.286 & 2.285 & 2.285 & 591.36 & 24379.3 & 0 & 4.00/4 \\
\hline
\code{p2--CCp1--p1} & hexa & 0.0 & 2.149 & 2.608 & 2.201 & 2.356 & 2.203 & 2.337 & 2.019 & 2.209 & 2.154 & 2.292 & 2.155 & 2.244 & 119.60 & 19444.1 & 0 & 4.00/4 \\
\code{p2--CCp1--p1} & hexa & 0.1 & 2.151 & 2.651 & 2.208 & 2.356 & 2.211 & 2.338 & 2.021 & 2.247 & 2.171 & 2.291 & 2.172 & 2.246 & 129.42 & 20862.0 & 0 & 4.00/4 \\
\code{p2--CCp1--p1} & triang. no estr. & 0.0 & 2.056 & 2.374 & 2.166 & 2.166 & 2.162 & 2.162 & 1.982 & 2.068 & 2.261 & 2.261 & 2.263 & 2.263 & 628.11 & 22471.6 & 0 & 4.00/4 \\
\code{p2--CCp1--p1} & triang. no estr. & 0.1 & 2.043 & 2.402 & 2.162 & 2.162 & 2.158 & 2.158 & 1.980 & 2.081 & 2.286 & 2.286 & 2.285 & 2.285 & 644.03 & 24808.0 & 0 & 4.00/4 \\
\hline
\code{p2--CCp1--p2} & hexa & 0.0 & 2.212 & 2.708 & 2.272 & 2.230 & 2.283 & 2.188 & 2.026 & 2.303 & 2.199 & 2.150 & 2.206 & 2.095 & 100.91 & 19646.5 & 0 & 4.00/4 \\
\code{p2--CCp1--p2} & hexa & 0.1 & 2.223 & 2.759 & 2.291 & 2.229 & 2.304 & 2.188 & 2.030 & 2.364 & 2.229 & 2.150 & 2.238 & 2.095 & 112.47 & 20964.3 & 0 & 4.00/4 \\
\code{p2--CCp1--p2} & triang. no estr. & 0.0 & 2.118 & 2.532 & 2.273 & 2.101 & 2.279 & 2.040 & 1.977 & 2.149 & 2.394 & 2.057 & 2.419 & 2.018 & 643.37 & 22685.6 & 0 & 4.00/4 \\
\code{p2--CCp1--p2} & triang. no estr. & 0.1 & 2.104 & 2.578 & 2.283 & 2.101 & 2.292 & 2.039 & 1.972 & 2.185 & 2.455 & 2.057 & 2.480 & 2.018 & 583.69 & 24950.0 & 0 & 4.00/4 \\
\hline
\code{p2--CCp1--p3} & hexa & 0.0 & 2.212 & 2.708 & 2.272 & 2.248 & 2.283 & 2.215 & 2.026 & 2.303 & 2.199 & 2.167 & 2.206 & 2.114 & 110.96 & 19740.6 & 0 & 4.00/4 \\
\code{p2--CCp1--p3} & hexa & 0.1 & 2.223 & 2.759 & 2.291 & 2.248 & 2.304 & 2.214 & 2.030 & 2.364 & 2.229 & 2.167 & 2.238 & 2.114 & 104.35 & 21027.4 & 0 & 4.00/4 \\
\code{p2--CCp1--p3} & triang. no estr. & 0.0 & 2.118 & 2.532 & 2.273 & 2.126 & 2.279 & 2.050 & 1.977 & 2.149 & 2.394 & 2.074 & 2.419 & 2.023 & 641.34 & 22793.8 & 0 & 4.00/4 \\
\code{p2--CCp1--p3} & triang. no estr. & 0.1 & 2.104 & 2.578 & 2.283 & 2.126 & 2.292 & 2.050 & 1.972 & 2.185 & 2.455 & 2.074 & 2.480 & 2.023 & 584.37 & 25015.4 & 0 & 4.00/4 \\
\hline
\code{p2--CCp2--p1} & hexa & 0.0 & 2.149 & 2.608 & 2.203 & 3.039 & 2.203 & 2.289 & 2.019 & 2.210 & 2.155 & 2.989 & 2.155 & 2.169 & 125.93 & 19924.8 & 0 & 4.00/4 \\
\code{p2--CCp2--p1} & hexa & 0.1 & 2.152 & 2.651 & 2.211 & 3.040 & 2.211 & 2.310 & 2.021 & 2.247 & 2.173 & 2.991 & 2.172 & 2.190 & 99.702 & 20884.2 & 0 & 4.00/4 \\
\code{p2--CCp2--p1} & triang. no estr. & 0.0 & 2.056 & 2.374 & 2.166 & 2.166 & 2.162 & 2.162 & 1.982 & 2.068 & 2.261 & 2.261 & 2.263 & 2.263 & 658.27 & 23562.2 & 0 & 4.00/4 \\
\code{p2--CCp2--p1} & triang. no estr. & 0.1 & 2.043 & 2.402 & 2.162 & 2.162 & 2.158 & 2.158 & 1.980 & 2.081 & 2.286 & 2.286 & 2.285 & 2.285 & 597.89 & 24901.6 & 0 & 4.00/4 \\
\hline
\code{p2--CCp2--p2} & hexa & 0.0 & 2.212 & 2.708 & 2.284 & 3.063 & 2.283 & 2.480 & 2.026 & 2.303 & 2.207 & 3.000 & 2.206 & 2.250 & 104.23 & 20117.0 & 0 & 4.00/4 \\
\code{p2--CCp2--p2} & hexa & 0.1 & 2.224 & 2.760 & 2.305 & 3.063 & 2.305 & 2.528 & 2.030 & 2.365 & 2.239 & 3.001 & 2.239 & 2.294 & 102.48 & 21009.3 & 0 & 4.00/4 \\
\code{p2--CCp2--p2} & triang. no estr. & 0.0 & 2.118 & 2.532 & 2.286 & 3.012 & 2.280 & 2.428 & 1.977 & 2.149 & 2.413 & 2.995 & 2.420 & 2.443 & 661.07 & 23788.5 & 0 & 4.00/4 \\
\code{p2--CCp2--p2} & triang. no estr. & 0.1 & 2.104 & 2.579 & 2.300 & 3.012 & 2.293 & 2.467 & 1.972 & 2.186 & 2.481 & 2.996 & 2.483 & 2.512 & 607.64 & 25036.1 & 0 & 4.00/4 \\
\hline
\code{p2--CCp2--p3} & hexa & 0.0 & 2.212 & 2.708 & 2.284 & 3.058 & 2.283 & 2.480 & 2.026 & 2.303 & 2.207 & 2.999 & 2.206 & 2.251 & 108.61 & 20220.7 & 0 & 4.00/4 \\
\code{p2--CCp2--p3} & hexa & 0.1 & 2.224 & 2.760 & 2.305 & 3.058 & 2.305 & 2.527 & 2.030 & 2.365 & 2.239 & 3.000 & 2.239 & 2.296 & 106.87 & 21070.0 & 0 & 4.00/4 \\
\code{p2--CCp2--p3} & triang. no estr. & 0.0 & 2.118 & 2.532 & 2.286 & 3.010 & 2.280 & 2.425 & 1.977 & 2.149 & 2.413 & 2.995 & 2.420 & 2.444 & 657.26 & 23893.4 & 0 & 4.00/4 \\
\code{p2--CCp2--p3} & triang. no estr. & 0.1 & 2.104 & 2.579 & 2.300 & 3.010 & 2.293 & 2.465 & 1.972 & 2.186 & 2.481 & 2.996 & 2.483 & 2.512 & 604.31 & 25115.3 & 0 & 4.00/4 \\
\hline
\code{p2--CCp3--p1} & hexa & 0.0 & 2.149 & 2.608 & 2.203 & 3.175 & 2.203 & 2.220 & 2.019 & 2.210 & 2.155 & 2.123 & 2.155 & 2.154 & 123.23 & 21147.6 & 0 & 4.00/4 \\
\code{p2--CCp3--p1} & hexa & 0.1 & 2.152 & 2.651 & 2.211 & 3.250 & 2.211 & 2.227 & 2.021 & 2.247 & 2.173 & 2.151 & 2.172 & 2.171 & 123.06 & 21550.1 & 0 & 4.00/4 \\
\code{p2--CCp3--p1} & triang. no estr. & 0.0 & 2.056 & 2.374 & 2.166 & 2.166 & 2.162 & 2.162 & 1.982 & 2.068 & 2.261 & 2.261 & 2.263 & 2.263 & 701.21 & 26008.8 & 0 & 4.00/4 \\
\code{p2--CCp3--p1} & triang. no estr. & 0.1 & 2.043 & 2.402 & 2.162 & 2.162 & 2.158 & 2.158 & 1.980 & 2.081 & 2.286 & 2.286 & 2.285 & 2.285 & 636.67 & 26829.5 & 0 & 4.00/4 \\
\hline
\code{p2--CCp3--p2} & hexa & 0.0 & 2.212 & 2.708 & 2.284 & 3.846 & 2.283 & 2.347 & 2.026 & 2.303 & 2.207 & 2.397 & 2.206 & 2.204 & 130.67 & 21307.4 & 0 & 4.00/4 \\
\code{p2--CCp3--p2} & hexa & 0.1 & 2.224 & 2.760 & 2.305 & 3.936 & 2.304 & 2.369 & 2.030 & 2.365 & 2.239 & 2.584 & 2.239 & 2.237 & 126.33 & 21721.4 & 0 & 4.00/4 \\
\code{p2--CCp3--p2} & triang. no estr. & 0.0 & 2.118 & 2.532 & 2.286 & 3.422 & 2.279 & 2.345 & 1.977 & 2.149 & 2.413 & 2.437 & 2.420 & 2.419 & 706.14 & 26226.4 & 0 & 4.00/4 \\
\code{p2--CCp3--p2} & triang. no estr. & 0.1 & 2.104 & 2.579 & 2.300 & 3.528 & 2.293 & 2.358 & 1.972 & 2.186 & 2.481 & 2.550 & 2.483 & 2.482 & 647.27 & 27058.1 & 0 & 4.00/4 \\
\hline
\code{p2--CCp3--p3} & hexa & 0.0 & 2.212 & 2.708 & 2.284 & 3.839 & 2.283 & 2.346 & 2.026 & 2.303 & 2.207 & 2.351 & 2.206 & 2.204 & 133.61 & 21423.4 & 0 & 4.00/4 \\
\code{p2--CCp3--p3} & hexa & 0.1 & 2.224 & 2.760 & 2.305 & 3.933 & 2.304 & 2.368 & 2.030 & 2.365 & 2.239 & 2.531 & 2.239 & 2.237 & 131.80 & 21828.1 & 0 & 4.00/4 \\
\code{p2--CCp3--p3} & triang. no estr. & 0.0 & 2.118 & 2.532 & 2.286 & 3.389 & 2.279 & 2.343 & 1.977 & 2.149 & 2.413 & 2.404 & 2.420 & 2.419 & 687.20 & 26345.2 & 0 & 4.00/4 \\
\code{p2--CCp3--p3} & triang. no estr. & 0.1 & 2.104 & 2.579 & 2.300 & 3.497 & 2.293 & 2.356 & 1.972 & 2.186 & 2.481 & 2.509 & 2.483 & 2.482 & 651.90 & 27158.1 & 0 & 4.00/4 \\
\hline
\code{p2--GP--p1} & hexa & 0.0 & 2.149 & 2.608 & 2.004 & 2.716 & 2.109 & 2.720 & 2.019 & 2.210 & 2.002 & 2.416 & 2.034 & 2.422 & 113.49 & 18606.1 & 0 & 4.00/4 \\
\code{p2--GP--p1} & hexa & 0.1 & 2.152 & 2.651 & 2.004 & 2.757 & 2.107 & 2.760 & 2.021 & 2.247 & 2.002 & 2.470 & 2.033 & 2.476 & 95.066 & 20320.7 & 0 & 4.00/4 \\
\code{p2--GP--p1} & triang. no estr. & 0.0 & 2.056 & 2.374 & 2.166 & 2.166 & 2.162 & 2.162 & 1.982 & 2.068 & 2.261 & 2.261 & 2.263 & 2.263 & 613.49 & 21589.3 & 0 & 4.00/4 \\
\code{p2--GP--p1} & triang. no estr. & 0.1 & 2.043 & 2.402 & 2.162 & 2.162 & 2.158 & 2.158 & 1.980 & 2.081 & 2.286 & 2.286 & 2.285 & 2.285 & 601.41 & 24401.3 & 0 & 4.00/4 \\
\hline
\code{p2--GP--p2} & hexa & 0.0 & 2.212 & 2.708 & 2.004 & 2.942 & 2.118 & 2.941 & 2.026 & 2.303 & 2.001 & 2.735 & 2.020 & 2.738 & 101.01 & 18797.0 & 0 & 4.00/4 \\
\code{p2--GP--p2} & hexa & 0.1 & 2.224 & 2.760 & 2.004 & 2.969 & 2.116 & 2.968 & 2.030 & 2.365 & 2.001 & 2.794 & 2.019 & 2.798 & 97.950 & 20436.8 & 0 & 4.00/4 \\
\code{p2--GP--p2} & triang. no estr. & 0.0 & 2.118 & 2.532 & 1.993 & 2.860 & 2.044 & 2.862 & 1.977 & 2.149 & 2.000 & 2.751 & 2.026 & 2.758 & 607.86 & 21808.2 & 0 & 4.00/4 \\
\code{p2--GP--p2} & triang. no estr. & 0.1 & 2.104 & 2.579 & 1.993 & 2.895 & 2.042 & 2.896 & 1.972 & 2.186 & 2.000 & 2.820 & 2.025 & 2.823 & 595.60 & 24545.1 & 0 & 4.00/4 \\
\hline
\code{p2--GP--p3} & hexa & 0.0 & 2.212 & 2.708 & 2.004 & 2.940 & 2.118 & 2.940 & 2.026 & 2.303 & 2.001 & 2.739 & 2.020 & 2.743 & 115.81 & 18890.9 & 0 & 4.00/4 \\
\code{p2--GP--p3} & hexa & 0.1 & 2.224 & 2.760 & 2.004 & 2.967 & 2.117 & 2.966 & 2.030 & 2.365 & 2.001 & 2.798 & 2.019 & 2.801 & 105.75 & 20496.0 & 0 & 4.00/4 \\
\code{p2--GP--p3} & triang. no estr. & 0.0 & 2.118 & 2.532 & 1.993 & 2.863 & 2.044 & 2.865 & 1.977 & 2.149 & 2.000 & 2.758 & 2.026 & 2.764 & 600.35 & 21930.4 & 0 & 4.00/4 \\
\code{p2--GP--p3} & triang. no estr. & 0.1 & 2.104 & 2.579 & 1.993 & 2.897 & 2.042 & 2.898 & 1.972 & 2.186 & 2.000 & 2.825 & 2.025 & 2.829 & 571.33 & 24606.0 & 0 & 4.00/4 \\
\hline
\code{p3--na--NO} & hexa & 0.0 & 2.121 & 2.481 & 2.256 & 2.256 & 2.258 & 2.258 & 1.999 & 2.037 & 2.472 & 2.472 & 2.469 & 2.469 & 139.64 & 21801.9 & 0 & 4.00/4 \\
\code{p3--na--NO} & hexa & 0.1 & 2.104 & 2.476 & 2.238 & 2.238 & 2.240 & 2.240 & 1.996 & 2.034 & 2.468 & 2.468 & 2.465 & 2.465 & 131.88 & 24501.4 & 0 & 4.00/4 \\
\code{p3--na--NO} & triang. no estr. & 0.0 & 2.033 & 2.201 & 2.348 & 2.348 & 2.346 & 2.346 & 1.972 & 1.988 & 3.183 & 3.183 & 3.173 & 3.173 & 702.42 & 27521.8 & 0 & 4.00/4 \\
\code{p3--na--NO} & triang. no estr. & 0.1 & 2.021 & 2.196 & 2.333 & 2.333 & 2.331 & 2.331 & 1.972 & 1.987 & 3.183 & 3.183 & 3.173 & 3.173 & 668.22 & 31503.9 & 0 & 4.00/4 \\
\hline
\code{p3--CCp1--p1} & hexa & 0.0 & 2.180 & 2.629 & 2.347 & 2.353 & 2.360 & 2.339 & 2.003 & 2.055 & 2.608 & 2.289 & 2.617 & 2.247 & 148.33 & 22964.6 & 0 & 4.00/4 \\
\code{p3--CCp1--p1} & hexa & 0.1 & 2.155 & 2.624 & 2.323 & 2.353 & 2.335 & 2.339 & 1.998 & 2.051 & 2.600 & 2.289 & 2.609 & 2.247 & 150.54 & 25237.0 & 0 & 4.00/4 \\
\code{p3--CCp1--p1} & triang. no estr. & 0.0 & 2.033 & 2.201 & 2.348 & 2.348 & 2.346 & 2.346 & 1.972 & 1.988 & 3.183 & 3.183 & 3.173 & 3.173 & 731.36 & 28009.0 & 0 & 4.00/4 \\
\code{p3--CCp1--p1} & triang. no estr. & 0.1 & 2.021 & 2.196 & 2.333 & 2.333 & 2.331 & 2.331 & 1.972 & 1.987 & 3.183 & 3.183 & 3.173 & 3.173 & 731.11 & 31964.1 & 0 & 4.00/4 \\
\hline
\code{p3--CCp1--p2} & hexa & 0.0 & 2.672 & 3.361 & 2.672 & 2.229 & 2.497 & 2.184 & 2.142 & 2.316 & 1.404 & 2.150 & 1.266 & 2.093 & 149.27 & 23080.2 & 0 & 4.00/4 \\
\code{p3--CCp1--p2} & hexa & 0.1 & 2.615 & 3.360 & 2.623 & 2.229 & 2.442 & 2.184 & 2.103 & 2.281 & 1.305 & 2.150 & 1.208 & 2.093 & 163.93 & 25346.8 & 0 & 4.00/4 \\
\code{p3--CCp1--p2} & triang. no estr. & 0.0 & 2.439 & 2.843 & 2.196 & 2.101 & 2.045 & 2.038 & 1.999 & 2.026 & 0.132 & 2.057 & 0.463 & 2.016 & 761.76 & 28151.6 & 0 & 4.00/4 \\
\code{p3--CCp1--p2} & triang. no estr. & 0.1 & 2.369 & 2.829 & 2.121 & 2.101 & 1.975 & 2.038 & 1.995 & 2.023 & 0.125 & 2.057 & 0.459 & 2.016 & 716.74 & 32105.2 & 0 & 4.00/4 \\
\hline
\code{p3--CCp1--p3} & hexa & 0.0 & 2.672 & 3.361 & 2.672 & 2.247 & 2.497 & 2.210 & 2.142 & 2.316 & 1.404 & 2.166 & 1.266 & 2.111 & 142.46 & 23133.2 & 0 & 4.00/4 \\
\code{p3--CCp1--p3} & hexa & 0.1 & 2.615 & 3.360 & 2.623 & 2.247 & 2.442 & 2.210 & 2.103 & 2.281 & 1.305 & 2.166 & 1.208 & 2.111 & 148.27 & 25409.8 & 0 & 4.00/4 \\
\code{p3--CCp1--p3} & triang. no estr. & 0.0 & 2.439 & 2.843 & 2.196 & 2.126 & 2.045 & 2.048 & 1.999 & 2.026 & 0.132 & 2.073 & 0.463 & 2.021 & 736.86 & 28242.9 & 0 & 4.00/4 \\
\code{p3--CCp1--p3} & triang. no estr. & 0.1 & 2.369 & 2.829 & 2.120 & 2.126 & 1.975 & 2.048 & 1.995 & 2.023 & 0.125 & 2.073 & 0.459 & 2.021 & 707.29 & 32170.5 & 0 & 4.00/4 \\
\hline
\code{p3--CCp2--p1} & hexa & 0.0 & 2.181 & 2.631 & 2.362 & 3.042 & 2.362 & 2.762 & 2.003 & 2.056 & 2.629 & 2.999 & 2.620 & 2.746 & 146.87 & 23022.2 & 0 & 4.00/4 \\
\code{p3--CCp2--p1} & hexa & 0.1 & 2.156 & 2.626 & 2.337 & 3.042 & 2.337 & 2.758 & 1.998 & 2.051 & 2.621 & 2.999 & 2.611 & 2.740 & 142.23 & 25237.2 & 0 & 4.00/4 \\
\code{p3--CCp2--p1} & triang. no estr. & 0.0 & 2.033 & 2.201 & 2.348 & 2.348 & 2.346 & 2.346 & 1.972 & 1.988 & 3.183 & 3.183 & 3.173 & 3.173 & 743.90 & 28372.6 & 0 & 4.00/4 \\
\code{p3--CCp2--p1} & triang. no estr. & 0.1 & 2.021 & 2.196 & 2.333 & 2.333 & 2.331 & 2.331 & 1.972 & 1.987 & 3.183 & 3.183 & 3.173 & 3.173 & 753.95 & 31971.9 & 0 & 4.00/4 \\
\hline
\code{p3--CCp2--p2} & hexa & 0.0 & 3.520 & 3.971 & 2.698 & 3.063 & 2.707 & 3.071 & 4.171 & 4.124 & 0.669 & 3.002 & 0.640 & 2.724 & 147.94 & 23158.8 & 0 & 4.00/4 \\
\code{p3--CCp2--p2} & hexa & 0.1 & 3.547 & 4.002 & 2.648 & 3.063 & 2.656 & 3.069 & 4.105 & 4.100 & 0.514 & 3.002 & 0.488 & 2.716 & 151.97 & 25350.9 & 0 & 4.00/4 \\
\code{p3--CCp2--p2} & triang. no estr. & 0.0 & 3.749 & 3.924 & 2.081 & 3.013 & 2.060 & 2.855 & 2.166 & 2.959 & 0.011 & 2.997 & 0.007 & 1.543 & 744.77 & 28535.6 & 0 & 4.00/4 \\
\code{p3--CCp2--p2} & triang. no estr. & 0.1 & 3.693 & 3.923 & 1.982 & 3.013 & 1.961 & 2.852 & 1.803 & 2.888 & 0.004 & 2.997 & 0.001 & 1.542 & 693.97 & 32108.5 & 0 & 4.00/4 \\
\hline
\code{p3--CCp2--p3} & hexa & 0.0 & 3.520 & 3.971 & 2.698 & 3.058 & 2.707 & 3.064 & 4.171 & 4.124 & 0.669 & 3.001 & 0.640 & 2.728 & 155.77 & 23227.6 & 0 & 4.00/4 \\
\code{p3--CCp2--p3} & hexa & 0.1 & 3.547 & 4.002 & 2.648 & 3.058 & 2.656 & 3.062 & 4.105 & 4.100 & 0.514 & 3.001 & 0.488 & 2.721 & 148.70 & 25410.3 & 0 & 4.00/4 \\
\code{p3--CCp2--p3} & triang. no estr. & 0.0 & 3.748 & 3.924 & 2.081 & 3.011 & 2.059 & 2.848 & 2.167 & 2.959 & 0.011 & 2.996 & 0.007 & 1.566 & 753.06 & 28621.7 & 0 & 4.00/4 \\
\code{p3--CCp2--p3} & triang. no estr. & 0.1 & 3.692 & 3.923 & 1.982 & 3.011 & 1.960 & 2.845 & 1.803 & 2.888 & 0.004 & 2.996 & 0.001 & 1.564 & 686.54 & 32165.2 & 0 & 4.00/4 \\
\hline
\code{p3--CCp3--p1} & hexa & 0.0 & 2.181 & 2.631 & 2.362 & 4.082 & 2.362 & 2.385 & 2.003 & 2.056 & 2.629 & 3.391 & 2.620 & 2.620 & 168.18 & 23706.2 & 0 & 4.00/4 \\
\code{p3--CCp3--p1} & hexa & 0.1 & 2.156 & 2.626 & 2.337 & 4.082 & 2.337 & 2.361 & 1.998 & 2.051 & 2.621 & 3.389 & 2.611 & 2.611 & 160.28 & 25271.4 & 0 & 4.00/4 \\
\code{p3--CCp3--p1} & triang. no estr. & 0.0 & 2.033 & 2.201 & 2.348 & 2.348 & 2.346 & 2.346 & 1.972 & 1.988 & 3.183 & 3.183 & 3.173 & 3.173 & 798.28 & 30639.9 & 0 & 4.00/4 \\
\code{p3--CCp3--p1} & triang. no estr. & 0.1 & 2.021 & 2.196 & 2.333 & 2.333 & 2.331 & 2.331 & 1.972 & 1.987 & 3.183 & 3.183 & 3.173 & 3.173 & 715.06 & 32100.9 & 0 & 4.00/4 \\
\hline
\code{p3--CCp3--p2} & hexa & 0.0 & 3.520 & 3.971 & 2.705 & 4.177 & 2.708 & 2.869 & 4.171 & 4.124 & 0.670 & 3.889 & 0.642 & 0.729 & 173.31 & 23875.3 & 0 & 4.00/4 \\
\code{p3--CCp3--p2} & hexa & 0.1 & 3.547 & 4.002 & 2.656 & 4.177 & 2.658 & 2.839 & 4.105 & 4.100 & 0.515 & 3.888 & 0.489 & 0.577 & 154.17 & 25374.4 & 0 & 4.00/4 \\
\code{p3--CCp3--p2} & triang. no estr. & 0.0 & 3.749 & 3.924 & 2.085 & 3.906 & 2.061 & 2.111 & 2.166 & 2.959 & 0.011 & 2.694 & 0.007 & 0.004 & 784.14 & 30867.0 & 0 & 4.00/4 \\
\code{p3--CCp3--p2} & triang. no estr. & 0.1 & 3.693 & 3.924 & 1.986 & 3.906 & 1.962 & 2.027 & 1.803 & 2.888 & 0.004 & 2.693 & 0.001 & -0.001 & 721.74 & 32251.5 & 0 & 4.00/4 \\
\hline
\code{p3--CCp3--p3} & hexa & 0.0 & 3.520 & 3.971 & 2.705 & 4.182 & 2.708 & 2.865 & 4.171 & 4.124 & 0.670 & 3.874 & 0.642 & 0.728 & 184.56 & 23970.5 & 0 & 4.00/4 \\
\code{p3--CCp3--p3} & hexa & 0.1 & 3.547 & 4.002 & 2.656 & 4.182 & 2.658 & 2.835 & 4.105 & 4.100 & 0.515 & 3.873 & 0.489 & 0.577 & 160.92 & 25445.4 & 0 & 4.00/4 \\
\code{p3--CCp3--p3} & triang. no estr. & 0.0 & 3.749 & 3.924 & 2.084 & 3.889 & 2.061 & 2.107 & 2.167 & 2.959 & 0.011 & 2.591 & 0.007 & 0.004 & 784.89 & 30979.2 & 0 & 4.00/4 \\
\code{p3--CCp3--p3} & triang. no estr. & 0.1 & 3.693 & 3.924 & 1.986 & 3.889 & 1.962 & 2.022 & 1.803 & 2.888 & 0.004 & 2.591 & 0.001 & -0.001 & 719.74 & 32318.2 & 0 & 4.00/4 \\
\hline
\code{p3--GP--p1} & hexa & 0.0 & 2.181 & 2.631 & 2.002 & 2.762 & 1.992 & 2.758 & 2.003 & 2.056 & 2.002 & 2.661 & 2.008 & 2.651 & 139.74 & 21832.7 & 0 & 4.00/4 \\
\code{p3--GP--p1} & hexa & 0.1 & 2.156 & 2.626 & 2.002 & 2.756 & 1.991 & 2.752 & 1.998 & 2.051 & 2.002 & 2.653 & 2.008 & 2.642 & 137.21 & 24543.6 & 0 & 4.00/4 \\
\code{p3--GP--p1} & triang. no estr. & 0.0 & 2.033 & 2.201 & 2.348 & 2.348 & 2.346 & 2.346 & 1.972 & 1.988 & 3.183 & 3.183 & 3.173 & 3.173 & 723.30 & 27552.4 & 0 & 4.00/4 \\
\code{p3--GP--p1} & triang. no estr. & 0.1 & 2.021 & 2.196 & 2.333 & 2.333 & 2.331 & 2.331 & 1.972 & 1.987 & 3.183 & 3.183 & 3.173 & 3.173 & 697.04 & 31547.6 & 0 & 4.00/4 \\
\hline
\code{p3--GP--p2} & hexa & 0.0 & 3.520 & 3.971 & 2.000 & 3.744 & 1.999 & 3.748 & 4.171 & 4.124 & 2.001 & 1.457 & 2.006 & 1.395 & 127.20 & 21964.0 & 0 & 4.00/4 \\
\code{p3--GP--p2} & hexa & 0.1 & 3.547 & 4.002 & 1.999 & 3.757 & 1.998 & 3.761 & 4.105 & 4.100 & 2.001 & 1.347 & 2.006 & 1.283 & 134.32 & 24618.0 & 0 & 4.00/4 \\
\code{p3--GP--p2} & triang. no estr. & 0.0 & 3.749 & 3.924 & 1.992 & 2.844 & 2.005 & 2.843 & 2.166 & 2.959 & 2.000 & -0.092 & 2.012 & -0.097 & 721.62 & 27688.9 & 0 & 4.00/4 \\
\code{p3--GP--p2} & triang. no estr. & 0.1 & 3.693 & 3.924 & 1.992 & 2.836 & 2.005 & 2.836 & 1.803 & 2.888 & 2.000 & -0.098 & 2.012 & -0.103 & 725.45 & 31682.7 & 0 & 4.00/4 \\
\hline
\code{p3--GP--p3} & hexa & 0.0 & 3.520 & 3.971 & 2.000 & 3.738 & 1.999 & 3.742 & 4.171 & 4.124 & 2.001 & 1.419 & 2.006 & 1.354 & 135.58 & 22012.0 & 0 & 4.00/4 \\
\code{p3--GP--p3} & hexa & 0.1 & 3.547 & 4.002 & 2.000 & 3.751 & 1.999 & 3.755 & 4.105 & 4.100 & 2.001 & 1.304 & 2.006 & 1.237 & 142.03 & 24707.6 & 0 & 4.00/4 \\
\code{p3--GP--p3} & triang. no estr. & 0.0 & 3.749 & 3.924 & 1.992 & 2.812 & 2.005 & 2.811 & 2.166 & 2.959 & 2.000 & -0.108 & 2.012 & -0.114 & 732.26 & 27747.5 & 0 & 4.00/4 \\
\code{p3--GP--p3} & triang. no estr. & 0.1 & 3.693 & 3.924 & 1.992 & 2.803 & 2.005 & 2.803 & 1.803 & 2.888 & 2.000 & -0.115 & 2.012 & -0.120 & 691.27 & 31746.0 & 0 & 4.00/4 \\
\bottomrule
\end{longtable}
\normalsize

\paragraph{$\Delta t=10^{-3}$}
\tiny
\setlength{\tabcolsep}{1.4pt}
\begin{longtable}{lllrrrrrrrrrrrrrrrr}
\caption{Resultados manufacturados 2D para \code{Picard}, \code{crankNicolson}, masa \code{consistent}, estados \code{RKF45} y $\Delta t=10^{-3}$. El primer bloque de convergencia se ajusta sobre $N=10$--$100$; el segundo sobre $N=70$--$100$. La primera columna registra la terna \code{Vm-states-Iion} efectivamente usada en la corrida. \textbf{RB} es la cantidad de pasos de tiempo que agotaron el presupuesto de iteraciones no lineales, acumulada sobre los 91 tamaños de malla de la fila (4550 pasos de tiempo en total por fila). \code{Picard} acepta el último iterado con advertencia en vez de hacer rollback, de modo que para este método la columna cuenta pasos aceptados sin converger. \textbf{$n_{\mathrm{nl}}$} reporta la mediana y el máximo de iteraciones no lineales por paso sobre esos mismos pasos; el presupuesto configurado es \code{nonlinearIterations}$=2000$.}\label{tab:es-results-2d-picard-1p0em03}\\
\toprule
& & & \multicolumn{6}{c}{Ajuste $N=10$--$100$} & \multicolumn{6}{c}{Ajuste $N=70$--$100$} & & & & \\
\cmidrule(lr){4-9}\cmidrule(lr){10-15}
\code{Vm-states-Iion} & Tipo de malla & $\alpha$ & $V_m$ & $V_m^G$ & $u_1$ & $u_1^G$ & $u_2$ & $u_2^G$ & $V_m$ & $V_m^G$ & $u_1$ & $u_1^G$ & $u_2$ & $u_2^G$ & $t_{\Sigma}$ [s] & RSS$_{\Sigma}$ [MB] & RB & $n_{\mathrm{nl}}$ \\
\midrule
\endfirsthead
\toprule
& & & \multicolumn{6}{c}{Ajuste $N=10$--$100$} & \multicolumn{6}{c}{Ajuste $N=70$--$100$} & & & & \\
\cmidrule(lr){4-9}\cmidrule(lr){10-15}
\code{Vm-states-Iion} & Tipo de malla & $\alpha$ & $V_m$ & $V_m^G$ & $u_1$ & $u_1^G$ & $u_2$ & $u_2^G$ & $V_m$ & $V_m^G$ & $u_1$ & $u_1^G$ & $u_2$ & $u_2^G$ & $t_{\Sigma}$ [s] & RSS$_{\Sigma}$ [MB] & RB & $n_{\mathrm{nl}}$ \\
\midrule
\endhead
\code{NO--na--NO} & hexa & 0.0 & 1.997 & 1.997 & 1.998 & 1.998 & 1.998 & 1.998 & 2.000 & 2.000 & 2.004 & 2.004 & 2.004 & 2.004 & 37.411 & 16434.0 & 0 & 3.00/3 \\
\code{NO--na--NO} & hexa & 0.1 & 1.851 & 1.851 & 1.607 & 1.607 & 1.613 & 1.613 & 2.294 & 2.294 & 2.140 & 2.140 & 2.142 & 2.142 & 60.356 & 18355.3 & 0 & 3.00/3 \\
\code{NO--na--NO} & triang. no estr. & 0.0 & 0.798 & 0.798 & 0.777 & 0.777 & 0.783 & 0.783 & 0.963 & 0.963 & 1.038 & 1.038 & 1.107 & 1.107 & 356.55 & 17094.1 & 0 & 3.00/3 \\
\code{NO--na--NO} & triang. no estr. & 0.1 & 0.785 & 0.785 & 0.771 & 0.771 & 0.776 & 0.776 & 0.953 & 0.953 & 1.027 & 1.027 & 1.097 & 1.097 & 409.76 & 20788.7 & 0 & 3.00/3 \\
\hline
\code{NO--CCp1--p1} & hexa & 0.0 & 1.997 & 1.997 & 1.999 & 2.354 & 1.998 & 2.333 & 2.000 & 2.000 & 2.004 & 2.290 & 2.004 & 2.242 & 65.100 & 16893.3 & 0 & 3.00/3 \\
\code{NO--CCp1--p1} & hexa & 0.1 & 1.851 & 1.851 & 1.608 & 2.354 & 1.613 & 2.310 & 2.294 & 2.294 & 2.140 & 2.290 & 2.142 & 2.226 & 103.65 & 18467.3 & 0 & 3.00/3 \\
\code{NO--CCp1--p1} & triang. no estr. & 0.0 & 0.798 & 0.798 & 0.777 & 0.777 & 0.783 & 0.783 & 0.963 & 0.963 & 1.038 & 1.038 & 1.107 & 1.107 & 466.11 & 17834.6 & 0 & 3.00/3 \\
\code{NO--CCp1--p1} & triang. no estr. & 0.1 & 0.785 & 0.785 & 0.771 & 0.771 & 0.776 & 0.776 & 0.953 & 0.953 & 1.027 & 1.027 & 1.097 & 1.097 & 445.29 & 20897.4 & 0 & 3.00/3 \\
\hline
\code{NO--CCp1--p2} & hexa & 0.0 & 1.997 & 1.997 & 2.000 & 2.229 & 1.999 & 2.184 & 2.000 & 2.000 & 2.004 & 2.150 & 2.004 & 2.094 & 77.641 & 17123.0 & 0 & 3.00/3 \\
\code{NO--CCp1--p2} & hexa & 0.1 & 1.851 & 1.851 & 1.608 & 2.229 & 1.614 & 2.181 & 2.295 & 2.295 & 2.140 & 2.150 & 2.142 & 2.092 & 75.766 & 18575.5 & 0 & 3.00/3 \\
\code{NO--CCp1--p2} & triang. no estr. & 0.0 & 0.798 & 0.798 & 0.777 & 2.111 & 0.783 & 1.163 & 0.963 & 0.963 & 1.038 & 2.102 & 1.107 & 1.150 & 488.02 & 18105.3 & 0 & 3.00/3 \\
\code{NO--CCp1--p2} & triang. no estr. & 0.1 & 0.785 & 0.785 & 0.771 & 2.112 & 0.776 & 1.171 & 0.953 & 0.953 & 1.027 & 2.105 & 1.097 & 1.143 & 494.01 & 21026.6 & 0 & 3.00/3 \\
\hline
\code{NO--CCp1--p3} & hexa & 0.0 & 1.997 & 1.997 & 2.000 & 2.247 & 1.999 & 2.211 & 2.000 & 2.000 & 2.004 & 2.167 & 2.004 & 2.113 & 82.027 & 17228.8 & 0 & 3.00/3 \\
\code{NO--CCp1--p3} & hexa & 0.1 & 1.851 & 1.851 & 1.608 & 2.247 & 1.614 & 2.206 & 2.295 & 2.295 & 2.140 & 2.167 & 2.142 & 2.110 & 106.15 & 18649.7 & 0 & 3.00/3 \\
\code{NO--CCp1--p3} & triang. no estr. & 0.0 & 0.798 & 0.798 & 0.777 & 2.139 & 0.783 & 1.131 & 0.963 & 0.963 & 1.038 & 2.131 & 1.107 & 1.144 & 510.96 & 18236.0 & 0 & 3.00/3 \\
\code{NO--CCp1--p3} & triang. no estr. & 0.1 & 0.785 & 0.785 & 0.771 & 2.140 & 0.776 & 1.138 & 0.953 & 0.953 & 1.027 & 2.135 & 1.097 & 1.137 & 489.56 & 21093.5 & 0 & 3.00/3 \\
\hline
\code{NO--CCp2--p1} & hexa & 0.0 & 1.997 & 1.997 & 2.000 & 3.042 & 1.997 & 2.510 & 1.999 & 1.999 & 2.004 & 2.998 & 2.004 & 2.157 & 98.814 & 17419.8 & 0 & 3.00/3 \\
\code{NO--CCp2--p1} & hexa & 0.1 & 1.851 & 1.851 & 1.603 & 3.038 & 1.613 & 1.942 & 2.294 & 2.294 & 2.140 & 2.987 & 2.142 & 2.167 & 132.60 & 18553.2 & 0 & 3.00/3 \\
\code{NO--CCp2--p1} & triang. no estr. & 0.0 & 0.798 & 0.798 & 0.777 & 0.777 & 0.783 & 0.783 & 0.963 & 0.963 & 1.038 & 1.038 & 1.107 & 1.107 & 577.70 & 19034.2 & 0 & 3.00/3 \\
\code{NO--CCp2--p1} & triang. no estr. & 0.1 & 0.785 & 0.785 & 0.771 & 0.771 & 0.776 & 0.776 & 0.953 & 0.953 & 1.027 & 1.027 & 1.097 & 1.097 & 569.23 & 20972.7 & 0 & 3.00/3 \\
\hline
\code{NO--CCp2--p2} & hexa & 0.0 & 1.995 & 1.995 & 2.003 & 3.063 & 1.997 & 2.671 & 1.999 & 1.999 & 2.004 & 3.002 & 2.003 & 2.246 & 122.24 & 17643.1 & 0 & 3.00/3 \\
\code{NO--CCp2--p2} & hexa & 0.1 & 1.853 & 1.853 & 1.597 & 3.061 & 1.619 & 2.063 & 2.289 & 2.289 & 2.140 & 2.996 & 2.141 & 2.192 & 110.37 & 18698.7 & 0 & 3.00/3 \\
\code{NO--CCp2--p2} & triang. no estr. & 0.0 & 0.799 & 0.799 & 0.777 & 1.940 & 0.783 & 0.775 & 0.963 & 0.963 & 1.038 & 1.095 & 1.107 & 1.097 & 591.70 & 19305.3 & 0 & 3.00/3 \\
\code{NO--CCp2--p2} & triang. no estr. & 0.1 & 0.786 & 0.786 & 0.771 & 1.955 & 0.776 & 0.772 & 0.953 & 0.953 & 1.027 & 1.090 & 1.097 & 1.087 & 634.90 & 21133.8 & 0 & 3.00/3 \\
\hline
\code{NO--CCp2--p3} & hexa & 0.0 & 1.995 & 1.995 & 2.003 & 3.058 & 1.997 & 2.670 & 1.999 & 1.999 & 2.004 & 3.001 & 2.003 & 2.251 & 137.62 & 17753.7 & 0 & 3.00/3 \\
\code{NO--CCp2--p3} & hexa & 0.1 & 1.853 & 1.853 & 1.597 & 3.056 & 1.619 & 2.064 & 2.289 & 2.289 & 2.140 & 2.995 & 2.141 & 2.193 & 154.38 & 18785.2 & 0 & 3.00/3 \\
\code{NO--CCp2--p3} & triang. no estr. & 0.0 & 0.799 & 0.799 & 0.777 & 1.950 & 0.783 & 0.775 & 0.963 & 0.963 & 1.038 & 1.098 & 1.107 & 1.097 & 611.98 & 19426.7 & 0 & 3.00/3 \\
\code{NO--CCp2--p3} & triang. no estr. & 0.1 & 0.786 & 0.786 & 0.771 & 1.966 & 0.776 & 0.772 & 0.953 & 0.953 & 1.027 & 1.092 & 1.097 & 1.087 & 607.74 & 21224.9 & 0 & 3.00/3 \\
\hline
\code{NO--CCp3--p1} & hexa & 0.0 & 1.997 & 1.997 & 1.998 & 3.944 & 1.997 & 2.030 & 1.999 & 1.999 & 2.004 & 2.505 & 2.004 & 2.004 & 226.62 & 18786.9 & 0 & 3.00/3 \\
\code{NO--CCp3--p1} & hexa & 0.1 & 1.851 & 1.851 & 1.600 & 3.114 & 1.613 & 1.657 & 2.294 & 2.294 & 2.140 & 2.324 & 2.142 & 2.170 & 182.05 & 19809.8 & 0 & 3.00/3 \\
\code{NO--CCp3--p1} & triang. no estr. & 0.0 & 0.798 & 0.798 & 0.777 & 0.777 & 0.783 & 0.783 & 0.963 & 0.963 & 1.038 & 1.038 & 1.107 & 1.107 & 734.01 & 21534.7 & 0 & 3.00/3 \\
\code{NO--CCp3--p1} & triang. no estr. & 0.1 & 0.785 & 0.785 & 0.771 & 0.771 & 0.776 & 0.776 & 0.953 & 0.953 & 1.027 & 1.027 & 1.097 & 1.097 & 741.32 & 23329.8 & 0 & 3.00/3 \\
\hline
\code{NO--CCp3--p2} & hexa & 0.0 & 1.996 & 1.996 & 1.998 & 4.244 & 1.996 & 2.092 & 1.999 & 1.999 & 2.004 & 3.862 & 2.003 & 2.002 & 226.96 & 19000.2 & 0 & 3.00/3 \\
\code{NO--CCp3--p2} & hexa & 0.1 & 1.853 & 1.853 & 1.590 & 3.420 & 1.617 & 1.718 & 2.289 & 2.289 & 2.140 & 2.581 & 2.141 & 2.203 & 247.54 & 20029.0 & 0 & 3.00/3 \\
\code{NO--CCp3--p2} & triang. no estr. & 0.0 & 0.799 & 0.799 & 0.777 & 1.070 & 0.783 & 0.795 & 0.963 & 0.963 & 1.038 & 1.038 & 1.107 & 1.110 & 748.16 & 21804.4 & 0 & 3.00/3 \\
\code{NO--CCp3--p2} & triang. no estr. & 0.1 & 0.786 & 0.786 & 0.771 & 1.080 & 0.776 & 0.797 & 0.953 & 0.953 & 1.027 & 1.028 & 1.097 & 1.100 & 751.27 & 23567.3 & 0 & 3.00/3 \\
\hline
\code{NO--CCp3--p3} & hexa & 0.0 & 1.996 & 1.996 & 1.998 & 4.257 & 1.996 & 2.090 & 1.999 & 1.999 & 2.004 & 3.849 & 2.003 & 2.002 & 216.69 & 19095.7 & 0 & 3.00/3 \\
\code{NO--CCp3--p3} & hexa & 0.1 & 1.853 & 1.853 & 1.590 & 3.408 & 1.617 & 1.717 & 2.289 & 2.289 & 2.140 & 2.570 & 2.141 & 2.203 & 194.93 & 20116.4 & 0 & 3.00/3 \\
\code{NO--CCp3--p3} & triang. no estr. & 0.0 & 0.799 & 0.799 & 0.777 & 1.053 & 0.783 & 0.795 & 0.963 & 0.963 & 1.038 & 1.038 & 1.107 & 1.110 & 738.56 & 21921.5 & 0 & 3.00/3 \\
\code{NO--CCp3--p3} & triang. no estr. & 0.1 & 0.786 & 0.786 & 0.771 & 1.062 & 0.776 & 0.796 & 0.953 & 0.953 & 1.027 & 1.028 & 1.097 & 1.100 & 782.87 & 23682.0 & 0 & 3.00/3 \\
\hline
\code{NO--GP--p1} & hexa & 0.0 & 1.997 & 1.997 & 1.999 & 2.050 & 1.998 & 2.053 & 2.000 & 2.000 & 2.000 & 2.007 & 2.000 & 2.007 & 50.315 & 16582.5 & 0 & 3.00/3 \\
\code{NO--GP--p1} & hexa & 0.1 & 1.851 & 1.851 & 1.999 & 2.042 & 1.990 & 2.046 & 2.294 & 2.294 & 2.000 & 2.015 & 2.008 & 2.016 & 73.314 & 18444.8 & 0 & 3.00/3 \\
\code{NO--GP--p1} & triang. no estr. & 0.0 & 0.798 & 0.798 & 0.777 & 0.777 & 0.783 & 0.783 & 0.963 & 0.963 & 1.038 & 1.038 & 1.107 & 1.107 & 405.08 & 17260.0 & 0 & 3.00/3 \\
\code{NO--GP--p1} & triang. no estr. & 0.1 & 0.785 & 0.785 & 0.771 & 0.771 & 0.776 & 0.776 & 0.953 & 0.953 & 1.027 & 1.027 & 1.097 & 1.097 & 469.96 & 20883.1 & 0 & 3.00/3 \\
\hline
\code{NO--GP--p2} & hexa & 0.0 & 1.997 & 1.997 & 1.999 & 2.013 & 1.998 & 2.015 & 1.999 & 1.999 & 2.000 & 2.002 & 2.000 & 2.002 & 110.15 & 16824.2 & 0 & 3.00/3 \\
\code{NO--GP--p2} & hexa & 0.1 & 1.853 & 1.853 & 1.999 & 2.013 & 1.998 & 2.014 & 2.290 & 2.290 & 2.000 & 2.003 & 2.001 & 2.003 & 107.80 & 18552.9 & 0 & 3.00/3 \\
\code{NO--GP--p2} & triang. no estr. & 0.0 & 0.799 & 0.799 & 1.878 & 1.417 & 1.101 & 1.410 & 0.963 & 0.963 & 1.677 & 1.218 & 1.112 & 1.269 & 425.12 & 17491.6 & 0 & 3.00/3 \\
\code{NO--GP--p2} & triang. no estr. & 0.1 & 0.786 & 0.786 & 1.880 & 1.426 & 1.105 & 1.419 & 0.953 & 0.953 & 1.680 & 1.219 & 1.104 & 1.270 & 468.71 & 21015.9 & 0 & 3.00/3 \\
\hline
\code{NO--GP--p3} & hexa & 0.0 & 1.997 & 1.997 & 1.999 & 2.016 & 1.999 & 2.018 & 1.999 & 1.999 & 2.000 & 2.002 & 2.000 & 2.002 & 113.09 & 16926.6 & 0 & 3.00/3 \\
\code{NO--GP--p3} & hexa & 0.1 & 1.853 & 1.853 & 1.999 & 2.016 & 1.998 & 2.018 & 2.290 & 2.290 & 2.000 & 2.004 & 2.002 & 2.004 & 139.45 & 18620.6 & 0 & 3.00/3 \\
\code{NO--GP--p3} & triang. no estr. & 0.0 & 0.799 & 0.799 & 1.878 & 1.388 & 1.101 & 1.383 & 0.963 & 0.963 & 1.677 & 1.205 & 1.112 & 1.258 & 469.91 & 17633.3 & 0 & 3.00/3 \\
\code{NO--GP--p3} & triang. no estr. & 0.1 & 0.786 & 0.786 & 1.880 & 1.397 & 1.106 & 1.391 & 0.953 & 0.953 & 1.680 & 1.206 & 1.104 & 1.258 & 503.51 & 21076.0 & 0 & 3.00/3 \\
\hline
\code{p1--na--NO} & hexa & 0.0 & 1.912 & 2.147 & 1.616 & 1.616 & 1.626 & 1.626 & 2.611 & 2.168 & 2.380 & 2.380 & 2.385 & 2.385 & 119.56 & 17760.4 & 0 & 3.00/3 \\
\code{p1--na--NO} & hexa & 0.1 & 1.971 & 2.163 & 1.677 & 1.677 & 1.687 & 1.687 & 2.615 & 2.141 & 2.414 & 2.414 & 2.419 & 2.419 & 138.64 & 19053.5 & 0 & 3.00/3 \\
\code{p1--na--NO} & triang. no estr. & 0.0 & 2.292 & 2.131 & 2.004 & 2.004 & 2.003 & 2.003 & 2.600 & 2.164 & 2.545 & 2.545 & 2.542 & 2.542 & 1039.5 & 19596.0 & 0 & 3.00/3 \\
\code{p1--na--NO} & triang. no estr. & 0.1 & 2.345 & 2.112 & 2.092 & 2.092 & 2.089 & 2.089 & 2.552 & 2.095 & 2.516 & 2.516 & 2.515 & 2.515 & 884.21 & 21530.4 & 0 & 3.00/3 \\
\hline
\code{p1--CCp1--p1} & hexa & 0.0 & 1.912 & 2.147 & 1.616 & 2.350 & 1.626 & 1.955 & 2.611 & 2.168 & 2.380 & 2.295 & 2.385 & 2.369 & 145.59 & 18215.8 & 0 & 3.00/3 \\
\code{p1--CCp1--p1} & hexa & 0.1 & 1.971 & 2.163 & 1.677 & 2.351 & 1.687 & 2.045 & 2.615 & 2.141 & 2.414 & 2.294 & 2.419 & 2.382 & 175.98 & 19252.1 & 0 & 3.00/3 \\
\code{p1--CCp1--p1} & triang. no estr. & 0.0 & 2.292 & 2.131 & 2.004 & 2.004 & 2.003 & 2.003 & 2.600 & 2.164 & 2.545 & 2.545 & 2.542 & 2.542 & 1116.5 & 20203.4 & 0 & 3.00/3 \\
\code{p1--CCp1--p1} & triang. no estr. & 0.1 & 2.345 & 2.112 & 2.092 & 2.092 & 2.089 & 2.089 & 2.552 & 2.095 & 2.516 & 2.516 & 2.515 & 2.515 & 961.19 & 21599.8 & 0 & 3.00/3 \\
\hline
\code{p1--CCp1--p2} & hexa & 0.0 & 1.912 & 2.147 & 1.616 & 2.228 & 1.626 & 2.057 & 2.611 & 2.168 & 2.380 & 2.154 & 2.385 & 2.233 & 191.20 & 18422.0 & 0 & 3.00/3 \\
\code{p1--CCp1--p2} & hexa & 0.1 & 1.971 & 2.163 & 1.676 & 2.228 & 1.686 & 2.104 & 2.615 & 2.141 & 2.414 & 2.153 & 2.419 & 2.209 & 145.60 & 19364.9 & 0 & 3.00/3 \\
\code{p1--CCp1--p2} & triang. no estr. & 0.0 & 2.292 & 2.131 & 2.004 & 2.103 & 2.003 & 2.036 & 2.600 & 2.164 & 2.545 & 2.063 & 2.542 & 2.283 & 1189.9 & 20458.9 & 0 & 3.00/3 \\
\code{p1--CCp1--p2} & triang. no estr. & 0.1 & 2.345 & 2.112 & 2.092 & 2.103 & 2.089 & 2.066 & 2.552 & 2.095 & 2.516 & 2.061 & 2.515 & 2.188 & 956.58 & 21743.2 & 0 & 3.00/3 \\
\hline
\code{p1--CCp1--p3} & hexa & 0.0 & 1.912 & 2.147 & 1.616 & 2.246 & 1.626 & 2.063 & 2.611 & 2.168 & 2.380 & 2.171 & 2.385 & 2.256 & 167.32 & 18527.7 & 0 & 3.00/3 \\
\code{p1--CCp1--p3} & hexa & 0.1 & 1.971 & 2.163 & 1.676 & 2.246 & 1.686 & 2.115 & 2.615 & 2.141 & 2.414 & 2.170 & 2.419 & 2.234 & 167.15 & 19424.6 & 0 & 3.00/3 \\
\code{p1--CCp1--p3} & triang. no estr. & 0.0 & 2.292 & 2.131 & 2.004 & 2.127 & 2.003 & 2.040 & 2.600 & 2.164 & 2.545 & 2.081 & 2.542 & 2.312 & 1277.0 & 20590.8 & 0 & 3.00/3 \\
\code{p1--CCp1--p3} & triang. no estr. & 0.1 & 2.345 & 2.112 & 2.092 & 2.128 & 2.089 & 2.076 & 2.552 & 2.095 & 2.516 & 2.078 & 2.515 & 2.216 & 992.42 & 21805.8 & 0 & 3.00/3 \\
\hline
\code{p1--CCp2--p1} & hexa & 0.0 & 1.912 & 2.147 & 1.616 & 2.903 & 1.626 & 1.665 & 2.611 & 2.168 & 2.380 & 2.782 & 2.385 & 2.394 & 178.87 & 18729.5 & 0 & 3.00/3 \\
\code{p1--CCp2--p1} & hexa & 0.1 & 1.971 & 2.163 & 1.676 & 2.953 & 1.687 & 1.736 & 2.615 & 2.141 & 2.414 & 2.858 & 2.419 & 2.428 & 180.49 & 19321.9 & 0 & 3.00/3 \\
\code{p1--CCp2--p1} & triang. no estr. & 0.0 & 2.292 & 2.131 & 2.004 & 2.004 & 2.003 & 2.003 & 2.600 & 2.164 & 2.545 & 2.545 & 2.542 & 2.542 & 1330.2 & 21406.7 & 0 & 3.00/3 \\
\code{p1--CCp2--p1} & triang. no estr. & 0.1 & 2.345 & 2.112 & 2.092 & 2.092 & 2.089 & 2.089 & 2.552 & 2.095 & 2.516 & 2.516 & 2.515 & 2.515 & 1045.8 & 21863.4 & 0 & 3.00/3 \\
\hline
\code{p1--CCp2--p2} & hexa & 0.0 & 1.912 & 2.147 & 1.616 & 2.990 & 1.627 & 1.689 & 2.611 & 2.168 & 2.380 & 2.882 & 2.385 & 2.400 & 175.61 & 18938.8 & 0 & 3.00/3 \\
\code{p1--CCp2--p2} & hexa & 0.1 & 1.972 & 2.163 & 1.676 & 3.018 & 1.687 & 1.765 & 2.615 & 2.141 & 2.414 & 2.931 & 2.419 & 2.435 & 216.23 & 19466.0 & 0 & 3.00/3 \\
\code{p1--CCp2--p2} & triang. no estr. & 0.0 & 2.292 & 2.131 & 2.004 & 2.676 & 2.003 & 1.978 & 2.600 & 2.164 & 2.546 & 2.665 & 2.542 & 2.522 & 1361.8 & 21636.3 & 0 & 3.00/3 \\
\code{p1--CCp2--p2} & triang. no estr. & 0.1 & 2.345 & 2.112 & 2.092 & 2.799 & 2.089 & 2.063 & 2.552 & 2.095 & 2.516 & 2.731 & 2.515 & 2.504 & 1202.3 & 22085.7 & 0 & 3.00/3 \\
\hline
\code{p1--CCp2--p3} & hexa & 0.0 & 1.912 & 2.147 & 1.616 & 2.987 & 1.627 & 1.688 & 2.611 & 2.168 & 2.380 & 2.884 & 2.385 & 2.400 & 163.85 & 19042.7 & 0 & 3.00/3 \\
\code{p1--CCp2--p3} & hexa & 0.1 & 1.972 & 2.163 & 1.676 & 3.015 & 1.687 & 1.765 & 2.615 & 2.141 & 2.414 & 2.931 & 2.419 & 2.435 & 227.61 & 19539.8 & 0 & 3.00/3 \\
\code{p1--CCp2--p3} & triang. no estr. & 0.0 & 2.292 & 2.131 & 2.004 & 2.681 & 2.003 & 1.978 & 2.600 & 2.164 & 2.546 & 2.668 & 2.542 & 2.522 & 1316.6 & 21754.9 & 0 & 3.00/3 \\
\code{p1--CCp2--p3} & triang. no estr. & 0.1 & 2.345 & 2.112 & 2.092 & 2.802 & 2.089 & 2.062 & 2.552 & 2.095 & 2.516 & 2.736 & 2.515 & 2.504 & 1276.9 & 22203.2 & 0 & 3.00/3 \\
\hline
\code{p1--CCp3--p1} & hexa & 0.0 & 1.912 & 2.147 & 1.617 & 2.034 & 1.626 & 1.630 & 2.611 & 2.168 & 2.380 & 2.378 & 2.385 & 2.406 & 243.76 & 20052.7 & 0 & 3.00/3 \\
\code{p1--CCp3--p1} & hexa & 0.1 & 1.971 & 2.163 & 1.677 & 2.164 & 1.687 & 1.692 & 2.615 & 2.141 & 2.414 & 2.408 & 2.419 & 2.441 & 315.44 & 20522.5 & 0 & 3.00/3 \\
\code{p1--CCp3--p1} & triang. no estr. & 0.0 & 2.292 & 2.131 & 2.004 & 2.004 & 2.003 & 2.003 & 2.600 & 2.164 & 2.545 & 2.545 & 2.542 & 2.542 & 1178.4 & 23901.7 & 0 & 3.00/3 \\
\code{p1--CCp3--p1} & triang. no estr. & 0.1 & 2.345 & 2.112 & 2.092 & 2.092 & 2.089 & 2.089 & 2.552 & 2.095 & 2.516 & 2.516 & 2.515 & 2.515 & 1484.4 & 24289.4 & 0 & 3.00/3 \\
\hline
\code{p1--CCp3--p2} & hexa & 0.0 & 1.913 & 2.147 & 1.617 & 2.239 & 1.627 & 1.635 & 2.611 & 2.168 & 2.380 & 2.377 & 2.385 & 2.426 & 259.12 & 20250.1 & 0 & 3.00/3 \\
\code{p1--CCp3--p2} & hexa & 0.1 & 1.972 & 2.163 & 1.677 & 2.394 & 1.688 & 1.698 & 2.615 & 2.141 & 2.414 & 2.402 & 2.419 & 2.463 & 300.54 & 20715.8 & 0 & 3.00/3 \\
\code{p1--CCp3--p2} & triang. no estr. & 0.0 & 2.292 & 2.131 & 2.004 & 2.031 & 2.003 & 1.965 & 2.600 & 2.164 & 2.546 & 2.521 & 2.542 & 2.523 & 1225.7 & 24140.6 & 0 & 3.00/3 \\
\code{p1--CCp3--p2} & triang. no estr. & 0.1 & 2.345 & 2.112 & 2.092 & 2.141 & 2.089 & 2.052 & 2.552 & 2.095 & 2.516 & 2.506 & 2.515 & 2.510 & 1474.8 & 24516.2 & 0 & 3.00/3 \\
\hline
\code{p1--CCp3--p3} & hexa & 0.0 & 1.913 & 2.147 & 1.617 & 2.224 & 1.627 & 1.634 & 2.611 & 2.168 & 2.380 & 2.376 & 2.385 & 2.426 & 282.39 & 20356.2 & 0 & 3.00/3 \\
\code{p1--CCp3--p3} & hexa & 0.1 & 1.972 & 2.163 & 1.677 & 2.379 & 1.688 & 1.698 & 2.615 & 2.141 & 2.414 & 2.401 & 2.419 & 2.462 & 288.53 & 20809.4 & 0 & 3.00/3 \\
\code{p1--CCp3--p3} & triang. no estr. & 0.0 & 2.292 & 2.131 & 2.004 & 2.028 & 2.003 & 1.965 & 2.600 & 2.164 & 2.546 & 2.522 & 2.542 & 2.523 & 1239.2 & 24244.2 & 0 & 3.00/3 \\
\code{p1--CCp3--p3} & triang. no estr. & 0.1 & 2.345 & 2.112 & 2.092 & 2.136 & 2.089 & 2.052 & 2.552 & 2.095 & 2.516 & 2.506 & 2.515 & 2.510 & 1459.8 & 24637.2 & 0 & 3.00/3 \\
\hline
\code{p1--GP--p1} & hexa & 0.0 & 1.912 & 2.147 & 1.997 & 2.152 & 1.836 & 2.158 & 2.611 & 2.168 & 2.005 & 2.309 & 2.249 & 2.308 & 123.49 & 17807.9 & 0 & 3.00/3 \\
\code{p1--GP--p1} & hexa & 0.1 & 1.971 & 2.163 & 1.999 & 2.219 & 1.895 & 2.225 & 2.615 & 2.141 & 2.003 & 2.309 & 2.222 & 2.308 & 138.59 & 19090.8 & 0 & 3.00/3 \\
\code{p1--GP--p1} & triang. no estr. & 0.0 & 2.292 & 2.131 & 2.004 & 2.004 & 2.003 & 2.003 & 2.600 & 2.164 & 2.545 & 2.545 & 2.542 & 2.542 & 1204.8 & 19634.4 & 0 & 3.00/3 \\
\code{p1--GP--p1} & triang. no estr. & 0.1 & 2.345 & 2.112 & 2.092 & 2.092 & 2.089 & 2.089 & 2.552 & 2.095 & 2.516 & 2.516 & 2.515 & 2.515 & 1079.4 & 21558.4 & 0 & 3.00/3 \\
\hline
\code{p1--GP--p2} & hexa & 0.0 & 1.913 & 2.147 & 1.998 & 2.181 & 1.949 & 2.179 & 2.611 & 2.168 & 2.001 & 2.155 & 2.096 & 2.151 & 180.40 & 18008.2 & 0 & 3.00/3 \\
\code{p1--GP--p2} & hexa & 0.1 & 1.972 & 2.163 & 1.999 & 2.199 & 1.972 & 2.197 & 2.615 & 2.141 & 2.001 & 2.150 & 2.071 & 2.145 & 209.63 & 19212.9 & 0 & 3.00/3 \\
\code{p1--GP--p2} & triang. no estr. & 0.0 & 2.292 & 2.131 & 1.991 & 2.067 & 2.009 & 2.066 & 2.600 & 2.164 & 2.002 & 2.176 & 2.268 & 2.168 & 849.42 & 19867.6 & 0 & 3.00/3 \\
\code{p1--GP--p2} & triang. no estr. & 0.1 & 2.345 & 2.112 & 1.992 & 2.087 & 2.034 & 2.085 & 2.552 & 2.095 & 2.000 & 2.109 & 2.167 & 2.105 & 1197.1 & 21695.8 & 0 & 3.00/3 \\
\hline
\code{p1--GP--p3} & hexa & 0.0 & 1.913 & 2.147 & 1.999 & 2.194 & 1.949 & 2.193 & 2.611 & 2.168 & 2.001 & 2.170 & 2.096 & 2.166 & 249.46 & 18130.7 & 0 & 3.00/3 \\
\code{p1--GP--p3} & hexa & 0.1 & 1.972 & 2.163 & 1.999 & 2.215 & 1.973 & 2.213 & 2.615 & 2.141 & 2.001 & 2.165 & 2.071 & 2.160 & 247.86 & 19286.1 & 0 & 3.00/3 \\
\code{p1--GP--p3} & triang. no estr. & 0.0 & 2.292 & 2.131 & 1.992 & 2.080 & 2.009 & 2.079 & 2.600 & 2.164 & 2.002 & 2.209 & 2.268 & 2.200 & 895.96 & 19987.1 & 0 & 3.00/3 \\
\code{p1--GP--p3} & triang. no estr. & 0.1 & 2.345 & 2.112 & 1.992 & 2.106 & 2.034 & 2.104 & 2.552 & 2.095 & 2.000 & 2.134 & 2.167 & 2.129 & 1063.9 & 21763.2 & 0 & 3.00/3 \\
\hline
\code{p2--na--NO} & hexa & 0.0 & 2.128 & 2.568 & 2.142 & 2.142 & 2.143 & 2.143 & 2.020 & 2.183 & 2.025 & 2.025 & 2.026 & 2.026 & 159.22 & 18524.9 & 0 & 3.00/3 \\
\code{p2--na--NO} & hexa & 0.1 & 2.129 & 2.608 & 2.143 & 2.143 & 2.144 & 2.144 & 2.022 & 2.213 & 2.027 & 2.027 & 2.028 & 2.028 & 156.57 & 20279.8 & 0 & 3.00/3 \\
\code{p2--na--NO} & triang. no estr. & 0.0 & 2.058 & 2.376 & 2.090 & 2.090 & 2.086 & 2.086 & 1.991 & 2.077 & 1.998 & 1.998 & 2.002 & 2.002 & 832.28 & 21545.1 & 0 & 3.00/3 \\
\code{p2--na--NO} & triang. no estr. & 0.1 & 2.045 & 2.405 & 2.079 & 2.079 & 2.075 & 2.075 & 1.989 & 2.091 & 1.996 & 1.996 & 1.997 & 1.997 & 894.20 & 24372.9 & 0 & 3.00/3 \\
\hline
\code{p2--CCp1--p1} & hexa & 0.0 & 2.150 & 2.609 & 2.164 & 2.356 & 2.166 & 2.334 & 2.023 & 2.214 & 2.029 & 2.293 & 2.030 & 2.232 & 185.32 & 19424.4 & 0 & 3.00/3 \\
\code{p2--CCp1--p1} & hexa & 0.1 & 2.153 & 2.652 & 2.167 & 2.356 & 2.170 & 2.335 & 2.026 & 2.251 & 2.031 & 2.293 & 2.032 & 2.235 & 174.20 & 20861.9 & 0 & 3.00/3 \\
\code{p2--CCp1--p1} & triang. no estr. & 0.0 & 2.058 & 2.376 & 2.090 & 2.090 & 2.086 & 2.086 & 1.991 & 2.077 & 1.998 & 1.998 & 2.002 & 2.002 & 1017.2 & 22478.0 & 0 & 3.00/3 \\
\code{p2--CCp1--p1} & triang. no estr. & 0.1 & 2.045 & 2.405 & 2.079 & 2.079 & 2.075 & 2.075 & 1.989 & 2.091 & 1.996 & 1.996 & 1.997 & 1.997 & 1104.4 & 24819.2 & 0 & 3.00/3 \\
\hline
\code{p2--CCp1--p2} & hexa & 0.0 & 2.213 & 2.709 & 2.224 & 2.230 & 2.235 & 2.188 & 2.031 & 2.308 & 2.039 & 2.151 & 2.041 & 2.094 & 197.88 & 19638.1 & 0 & 3.00/3 \\
\code{p2--CCp1--p2} & hexa & 0.1 & 2.225 & 2.760 & 2.235 & 2.230 & 2.248 & 2.187 & 2.036 & 2.369 & 2.043 & 2.151 & 2.046 & 2.094 & 201.35 & 20980.1 & 0 & 3.00/3 \\
\code{p2--CCp1--p2} & triang. no estr. & 0.0 & 2.121 & 2.535 & 2.161 & 2.102 & 2.166 & 2.040 & 1.990 & 2.162 & 2.003 & 2.058 & 2.012 & 2.018 & 1037.7 & 22673.4 & 0 & 3.00/3 \\
\code{p2--CCp1--p2} & triang. no estr. & 0.1 & 2.108 & 2.582 & 2.154 & 2.101 & 2.160 & 2.039 & 1.987 & 2.199 & 1.998 & 2.058 & 2.003 & 2.018 & 1032.1 & 24958.1 & 0 & 3.00/3 \\
\hline
\code{p2--CCp1--p3} & hexa & 0.0 & 2.213 & 2.709 & 2.224 & 2.248 & 2.235 & 2.214 & 2.031 & 2.308 & 2.039 & 2.168 & 2.041 & 2.113 & 174.28 & 19730.4 & 0 & 3.00/3 \\
\code{p2--CCp1--p3} & hexa & 0.1 & 2.225 & 2.760 & 2.235 & 2.248 & 2.248 & 2.214 & 2.036 & 2.369 & 2.043 & 2.168 & 2.046 & 2.113 & 208.11 & 21044.0 & 0 & 3.00/3 \\
\code{p2--CCp1--p3} & triang. no estr. & 0.0 & 2.122 & 2.535 & 2.161 & 2.127 & 2.166 & 2.050 & 1.990 & 2.162 & 2.003 & 2.074 & 2.012 & 2.023 & 1010.9 & 22791.0 & 0 & 3.00/3 \\
\code{p2--CCp1--p3} & triang. no estr. & 0.1 & 2.108 & 2.582 & 2.154 & 2.126 & 2.160 & 2.050 & 1.987 & 2.199 & 1.998 & 2.074 & 2.003 & 2.023 & 1065.4 & 25021.3 & 0 & 3.00/3 \\
\hline
\code{p2--CCp2--p1} & hexa & 0.0 & 2.150 & 2.609 & 2.166 & 3.039 & 2.166 & 2.253 & 2.023 & 2.214 & 2.029 & 2.986 & 2.030 & 2.044 & 213.58 & 19889.8 & 0 & 3.00/3 \\
\code{p2--CCp2--p1} & hexa & 0.1 & 2.153 & 2.653 & 2.170 & 3.040 & 2.170 & 2.269 & 2.026 & 2.251 & 2.032 & 2.989 & 2.032 & 2.050 & 217.22 & 20885.5 & 0 & 3.00/3 \\
\code{p2--CCp2--p1} & triang. no estr. & 0.0 & 2.058 & 2.376 & 2.090 & 2.090 & 2.086 & 2.086 & 1.991 & 2.077 & 1.998 & 1.998 & 2.002 & 2.002 & 1138.4 & 23567.7 & 0 & 3.00/3 \\
\code{p2--CCp2--p1} & triang. no estr. & 0.1 & 2.045 & 2.405 & 2.079 & 2.079 & 2.075 & 2.075 & 1.989 & 2.091 & 1.996 & 1.996 & 1.997 & 1.997 & 1144.6 & 24899.3 & 0 & 3.00/3 \\
\hline
\code{p2--CCp2--p2} & hexa & 0.0 & 2.214 & 2.709 & 2.236 & 3.063 & 2.235 & 2.434 & 2.032 & 2.308 & 2.041 & 2.999 & 2.041 & 2.086 & 212.93 & 20102.2 & 0 & 3.00/3 \\
\code{p2--CCp2--p2} & hexa & 0.1 & 2.226 & 2.761 & 2.250 & 3.063 & 2.249 & 2.475 & 2.037 & 2.371 & 2.046 & 3.000 & 2.047 & 2.105 & 224.72 & 21011.9 & 0 & 3.00/3 \\
\code{p2--CCp2--p2} & triang. no estr. & 0.0 & 2.122 & 2.536 & 2.172 & 3.012 & 2.166 & 2.317 & 1.990 & 2.162 & 2.003 & 2.991 & 2.012 & 2.038 & 1164.8 & 23788.5 & 0 & 3.00/3 \\
\code{p2--CCp2--p2} & triang. no estr. & 0.1 & 2.108 & 2.583 & 2.166 & 3.012 & 2.160 & 2.339 & 1.987 & 2.200 & 1.999 & 2.993 & 2.002 & 2.039 & 1153.8 & 25037.3 & 0 & 3.00/3 \\
\hline
\code{p2--CCp2--p3} & hexa & 0.0 & 2.214 & 2.709 & 2.236 & 3.058 & 2.235 & 2.434 & 2.032 & 2.308 & 2.041 & 2.998 & 2.041 & 2.088 & 212.31 & 20200.3 & 0 & 3.00/3 \\
\code{p2--CCp2--p3} & hexa & 0.1 & 2.226 & 2.761 & 2.250 & 3.058 & 2.249 & 2.474 & 2.037 & 2.371 & 2.046 & 2.999 & 2.047 & 2.106 & 218.90 & 21067.4 & 0 & 3.00/3 \\
\code{p2--CCp2--p3} & triang. no estr. & 0.0 & 2.122 & 2.536 & 2.172 & 3.009 & 2.166 & 2.314 & 1.990 & 2.162 & 2.003 & 2.991 & 2.012 & 2.039 & 1149.3 & 23888.6 & 0 & 3.00/3 \\
\code{p2--CCp2--p3} & triang. no estr. & 0.1 & 2.108 & 2.583 & 2.166 & 3.010 & 2.160 & 2.336 & 1.987 & 2.200 & 1.999 & 2.993 & 2.002 & 2.040 & 1185.5 & 25120.3 & 0 & 3.00/3 \\
\hline
\code{p2--CCp3--p1} & hexa & 0.0 & 2.150 & 2.609 & 2.166 & 3.135 & 2.166 & 2.183 & 2.023 & 2.214 & 2.029 & 1.992 & 2.030 & 2.029 & 277.42 & 21135.1 & 0 & 3.00/3 \\
\code{p2--CCp3--p1} & hexa & 0.1 & 2.153 & 2.653 & 2.170 & 3.206 & 2.170 & 2.186 & 2.026 & 2.251 & 2.032 & 2.002 & 2.032 & 2.031 & 304.67 & 21532.0 & 0 & 3.00/3 \\
\code{p2--CCp3--p1} & triang. no estr. & 0.0 & 2.058 & 2.376 & 2.090 & 2.090 & 2.086 & 2.086 & 1.991 & 2.077 & 1.998 & 1.998 & 2.002 & 2.002 & 1287.9 & 26014.3 & 0 & 3.00/3 \\
\code{p2--CCp3--p1} & triang. no estr. & 0.1 & 2.045 & 2.405 & 2.079 & 2.079 & 2.075 & 2.075 & 1.989 & 2.091 & 1.996 & 1.996 & 1.997 & 1.997 & 1342.9 & 26843.0 & 0 & 3.00/3 \\
\hline
\code{p2--CCp3--p2} & hexa & 0.0 & 2.214 & 2.709 & 2.236 & 3.796 & 2.235 & 2.299 & 2.032 & 2.308 & 2.041 & 2.178 & 2.041 & 2.039 & 293.32 & 21317.7 & 0 & 3.00/3 \\
\code{p2--CCp3--p2} & hexa & 0.1 & 2.226 & 2.761 & 2.249 & 3.882 & 2.249 & 2.314 & 2.037 & 2.371 & 2.046 & 2.315 & 2.046 & 2.044 & 304.93 & 21712.7 & 0 & 3.00/3 \\
\code{p2--CCp3--p2} & triang. no estr. & 0.0 & 2.122 & 2.536 & 2.172 & 3.303 & 2.166 & 2.231 & 1.990 & 2.162 & 2.003 & 1.988 & 2.012 & 2.010 & 1272.6 & 26231.3 & 0 & 3.00/3 \\
\code{p2--CCp3--p2} & triang. no estr. & 0.1 & 2.108 & 2.583 & 2.166 & 3.390 & 2.160 & 2.225 & 1.987 & 2.200 & 1.999 & 2.010 & 2.002 & 2.001 & 1357.1 & 27057.7 & 0 & 3.00/3 \\
\hline
\code{p2--CCp3--p3} & hexa & 0.0 & 2.214 & 2.709 & 2.236 & 3.789 & 2.235 & 2.298 & 2.032 & 2.308 & 2.041 & 2.133 & 2.041 & 2.039 & 296.50 & 21401.5 & 0 & 3.00/3 \\
\code{p2--CCp3--p3} & hexa & 0.1 & 2.226 & 2.761 & 2.249 & 3.877 & 2.249 & 2.313 & 2.037 & 2.371 & 2.046 & 2.262 & 2.046 & 2.044 & 308.41 & 21792.9 & 0 & 3.00/3 \\
\code{p2--CCp3--p3} & triang. no estr. & 0.0 & 2.122 & 2.536 & 2.172 & 3.267 & 2.166 & 2.229 & 1.990 & 2.162 & 2.003 & 1.957 & 2.012 & 2.010 & 1416.7 & 26349.1 & 0 & 3.00/3 \\
\code{p2--CCp3--p3} & triang. no estr. & 0.1 & 2.108 & 2.583 & 2.166 & 3.356 & 2.160 & 2.223 & 1.987 & 2.200 & 1.999 & 1.972 & 2.002 & 2.001 & 1354.9 & 27168.7 & 0 & 3.00/3 \\
\hline
\code{p2--GP--p1} & hexa & 0.0 & 2.150 & 2.609 & 2.003 & 2.689 & 2.102 & 2.694 & 2.023 & 2.214 & 2.001 & 2.304 & 2.011 & 2.312 & 136.92 & 18575.6 & 0 & 3.00/3 \\
\code{p2--GP--p1} & hexa & 0.1 & 2.153 & 2.653 & 2.003 & 2.728 & 2.100 & 2.732 & 2.026 & 2.251 & 2.001 & 2.350 & 2.011 & 2.358 & 160.65 & 20323.0 & 0 & 3.00/3 \\
\code{p2--GP--p1} & triang. no estr. & 0.0 & 2.058 & 2.376 & 2.090 & 2.090 & 2.086 & 2.086 & 1.991 & 2.077 & 1.998 & 1.998 & 2.002 & 2.002 & 913.79 & 21572.9 & 0 & 3.00/3 \\
\code{p2--GP--p1} & triang. no estr. & 0.1 & 2.045 & 2.405 & 2.079 & 2.079 & 2.075 & 2.075 & 1.989 & 2.091 & 1.996 & 1.996 & 1.997 & 1.997 & 933.55 & 24422.7 & 0 & 3.00/3 \\
\hline
\code{p2--GP--p2} & hexa & 0.0 & 2.214 & 2.709 & 2.003 & 2.924 & 2.114 & 2.923 & 2.032 & 2.308 & 2.000 & 2.644 & 2.007 & 2.650 & 182.72 & 18784.6 & 0 & 3.00/3 \\
\code{p2--GP--p2} & hexa & 0.1 & 2.226 & 2.761 & 2.003 & 2.952 & 2.112 & 2.951 & 2.037 & 2.371 & 2.000 & 2.706 & 2.007 & 2.711 & 232.36 & 20423.3 & 0 & 3.00/3 \\
\code{p2--GP--p2} & triang. no estr. & 0.0 & 2.122 & 2.536 & 1.992 & 2.805 & 2.037 & 2.807 & 1.990 & 2.162 & 1.998 & 2.489 & 2.003 & 2.499 & 912.84 & 21817.9 & 0 & 3.00/3 \\
\code{p2--GP--p2} & triang. no estr. & 0.1 & 2.108 & 2.583 & 1.992 & 2.840 & 2.035 & 2.842 & 1.987 & 2.200 & 1.998 & 2.553 & 2.002 & 2.561 & 1016.7 & 24550.7 & 0 & 3.00/3 \\
\hline
\code{p2--GP--p3} & hexa & 0.0 & 2.214 & 2.709 & 2.004 & 2.923 & 2.114 & 2.922 & 2.032 & 2.308 & 2.000 & 2.650 & 2.007 & 2.656 & 213.98 & 18882.5 & 0 & 3.00/3 \\
\code{p2--GP--p3} & hexa & 0.1 & 2.226 & 2.761 & 2.004 & 2.950 & 2.113 & 2.950 & 2.037 & 2.371 & 2.000 & 2.711 & 2.007 & 2.716 & 266.03 & 20500.2 & 0 & 3.00/3 \\
\code{p2--GP--p3} & triang. no estr. & 0.0 & 2.122 & 2.536 & 1.993 & 2.809 & 2.037 & 2.812 & 1.990 & 2.162 & 1.998 & 2.500 & 2.003 & 2.510 & 938.35 & 21942.5 & 0 & 3.00/3 \\
\code{p2--GP--p3} & triang. no estr. & 0.1 & 2.108 & 2.583 & 1.993 & 2.843 & 2.035 & 2.846 & 1.987 & 2.200 & 1.998 & 2.564 & 2.002 & 2.572 & 1037.7 & 24620.3 & 0 & 3.00/3 \\
\hline
\code{p3--na--NO} & hexa & 0.0 & 2.125 & 2.485 & 2.132 & 2.132 & 2.135 & 2.135 & 2.013 & 2.050 & 2.025 & 2.025 & 2.024 & 2.024 & 173.20 & 21790.9 & 0 & 3.00/3 \\
\code{p3--na--NO} & hexa & 0.1 & 2.108 & 2.481 & 2.114 & 2.114 & 2.117 & 2.117 & 2.010 & 2.048 & 2.020 & 2.020 & 2.020 & 2.020 & 164.48 & 24491.5 & 0 & 3.00/3 \\
\code{p3--na--NO} & triang. no estr. & 0.0 & 2.041 & 2.209 & 2.061 & 2.061 & 2.060 & 2.060 & 1.999 & 2.015 & 2.008 & 2.008 & 2.008 & 2.008 & 964.35 & 27515.3 & 0 & 3.00/3 \\
\code{p3--na--NO} & triang. no estr. & 0.1 & 2.029 & 2.204 & 2.046 & 2.046 & 2.046 & 2.046 & 1.999 & 2.014 & 2.008 & 2.008 & 2.007 & 2.007 & 994.75 & 31523.2 & 0 & 3.00/3 \\
\hline
\code{p3--CCp1--p1} & hexa & 0.0 & 2.185 & 2.634 & 2.188 & 2.354 & 2.201 & 2.339 & 2.020 & 2.072 & 2.038 & 2.290 & 2.040 & 2.248 & 237.07 & 22963.0 & 0 & 3.00/3 \\
\code{p3--CCp1--p1} & hexa & 0.1 & 2.160 & 2.629 & 2.164 & 2.354 & 2.176 & 2.339 & 2.016 & 2.068 & 2.031 & 2.290 & 2.032 & 2.248 & 235.69 & 25220.2 & 0 & 3.00/3 \\
\code{p3--CCp1--p1} & triang. no estr. & 0.0 & 2.041 & 2.209 & 2.061 & 2.061 & 2.060 & 2.060 & 1.999 & 2.015 & 2.008 & 2.008 & 2.008 & 2.008 & 1256.8 & 28007.2 & 0 & 3.00/3 \\
\code{p3--CCp1--p1} & triang. no estr. & 0.1 & 2.029 & 2.204 & 2.046 & 2.046 & 2.046 & 2.046 & 1.999 & 2.014 & 2.008 & 2.008 & 2.007 & 2.007 & 1171.8 & 31951.8 & 0 & 3.00/3 \\
\hline
\code{p3--CCp1--p2} & hexa & 0.0 & 2.620 & 3.309 & 2.640 & 2.229 & 2.603 & 2.184 & 2.117 & 2.267 & 2.264 & 2.150 & 2.211 & 2.094 & 255.61 & 23079.3 & 0 & 3.00/3 \\
\code{p3--CCp1--p2} & hexa & 0.1 & 2.564 & 3.307 & 2.592 & 2.229 & 2.554 & 2.184 & 2.084 & 2.237 & 2.206 & 2.150 & 2.160 & 2.094 & 240.70 & 25330.0 & 0 & 3.00/3 \\
\code{p3--CCp1--p2} & triang. no estr. & 0.0 & 2.396 & 2.791 & 2.510 & 2.101 & 2.460 & 2.038 & 2.003 & 2.026 & 2.041 & 2.058 & 2.014 & 2.017 & 1207.2 & 28162.8 & 0 & 3.00/3 \\
\code{p3--CCp1--p2} & triang. no estr. & 0.1 & 2.330 & 2.777 & 2.443 & 2.101 & 2.397 & 2.038 & 1.999 & 2.022 & 2.031 & 2.058 & 2.006 & 2.017 & 1169.7 & 32100.9 & 0 & 3.00/3 \\
\hline
\code{p3--CCp1--p3} & hexa & 0.0 & 2.620 & 3.309 & 2.640 & 2.247 & 2.603 & 2.210 & 2.117 & 2.267 & 2.264 & 2.167 & 2.211 & 2.112 & 275.00 & 23143.0 & 0 & 3.00/3 \\
\code{p3--CCp1--p3} & hexa & 0.1 & 2.564 & 3.307 & 2.592 & 2.247 & 2.554 & 2.210 & 2.084 & 2.237 & 2.206 & 2.167 & 2.160 & 2.112 & 262.22 & 25414.1 & 0 & 3.00/3 \\
\code{p3--CCp1--p3} & triang. no estr. & 0.0 & 2.396 & 2.791 & 2.509 & 2.126 & 2.460 & 2.048 & 2.003 & 2.026 & 2.041 & 2.074 & 2.014 & 2.023 & 1189.6 & 28235.2 & 0 & 3.00/3 \\
\code{p3--CCp1--p3} & triang. no estr. & 0.1 & 2.330 & 2.777 & 2.443 & 2.126 & 2.397 & 2.048 & 2.000 & 2.023 & 2.031 & 2.074 & 2.006 & 2.023 & 1252.4 & 32160.5 & 0 & 3.00/3 \\
\hline
\code{p3--CCp2--p1} & hexa & 0.0 & 2.187 & 2.637 & 2.200 & 3.042 & 2.203 & 2.643 & 2.020 & 2.073 & 2.041 & 2.999 & 2.040 & 2.257 & 293.35 & 23015.5 & 0 & 3.00/3 \\
\code{p3--CCp2--p1} & hexa & 0.1 & 2.162 & 2.631 & 2.175 & 3.042 & 2.177 & 2.638 & 2.016 & 2.069 & 2.033 & 2.999 & 2.033 & 2.252 & 285.17 & 25238.3 & 0 & 3.00/3 \\
\code{p3--CCp2--p1} & triang. no estr. & 0.0 & 2.041 & 2.209 & 2.061 & 2.061 & 2.060 & 2.060 & 1.999 & 2.015 & 2.008 & 2.008 & 2.008 & 2.008 & 1278.2 & 28367.6 & 0 & 3.00/3 \\
\code{p3--CCp2--p1} & triang. no estr. & 0.1 & 2.029 & 2.204 & 2.046 & 2.046 & 2.046 & 2.046 & 1.999 & 2.014 & 2.008 & 2.008 & 2.007 & 2.007 & 1352.5 & 31980.0 & 0 & 3.00/3 \\
\hline
\code{p3--CCp2--p2} & hexa & 0.0 & 3.539 & 3.992 & 3.284 & 3.063 & 3.303 & 3.110 & 4.323 & 4.278 & 4.078 & 3.002 & 4.085 & 3.055 & 277.42 & 23151.1 & 0 & 3.00/3 \\
\code{p3--CCp2--p2} & hexa & 0.1 & 3.570 & 4.025 & 3.312 & 3.063 & 3.332 & 3.109 & 4.298 & 4.271 & 4.078 & 3.002 & 4.086 & 3.049 & 272.86 & 25356.0 & 0 & 3.00/3 \\
\code{p3--CCp2--p2} & triang. no estr. & 0.0 & 4.057 & 4.096 & 3.842 & 3.013 & 3.861 & 3.085 & 4.476 & 4.165 & 4.249 & 2.997 & 4.256 & 3.023 & 1294.1 & 28527.1 & 0 & 3.00/3 \\
\code{p3--CCp2--p2} & triang. no estr. & 0.1 & 4.077 & 4.105 & 3.880 & 3.013 & 3.901 & 3.082 & 4.460 & 4.139 & 4.196 & 2.997 & 4.211 & 3.022 & 1435.7 & 32105.5 & 0 & 3.00/3 \\
\hline
\code{p3--CCp2--p3} & hexa & 0.0 & 3.539 & 3.992 & 3.284 & 3.058 & 3.303 & 3.102 & 4.323 & 4.278 & 4.078 & 3.001 & 4.085 & 3.050 & 262.11 & 23226.5 & 0 & 3.00/3 \\
\code{p3--CCp2--p3} & hexa & 0.1 & 3.570 & 4.025 & 3.312 & 3.058 & 3.332 & 3.101 & 4.298 & 4.271 & 4.078 & 3.001 & 4.086 & 3.044 & 289.51 & 25404.4 & 0 & 3.00/3 \\
\code{p3--CCp2--p3} & triang. no estr. & 0.0 & 4.057 & 4.096 & 3.842 & 3.011 & 3.861 & 3.072 & 4.476 & 4.165 & 4.250 & 2.997 & 4.257 & 3.019 & 1323.3 & 28624.5 & 0 & 3.00/3 \\
\code{p3--CCp2--p3} & triang. no estr. & 0.1 & 4.077 & 4.105 & 3.880 & 3.011 & 3.901 & 3.070 & 4.460 & 4.139 & 4.196 & 2.997 & 4.211 & 3.017 & 1410.4 & 32177.0 & 0 & 3.00/3 \\
\hline
\code{p3--CCp3--p1} & hexa & 0.0 & 2.186 & 2.637 & 2.200 & 3.993 & 2.202 & 2.225 & 2.020 & 2.073 & 2.041 & 2.812 & 2.040 & 2.040 & 326.31 & 23721.1 & 0 & 3.00/3 \\
\code{p3--CCp3--p1} & hexa & 0.1 & 2.161 & 2.631 & 2.174 & 3.993 & 2.177 & 2.201 & 2.016 & 2.069 & 2.033 & 2.809 & 2.033 & 2.033 & 356.50 & 25270.9 & 0 & 3.00/3 \\
\code{p3--CCp3--p1} & triang. no estr. & 0.0 & 2.041 & 2.209 & 2.061 & 2.061 & 2.060 & 2.060 & 1.999 & 2.015 & 2.008 & 2.008 & 2.008 & 2.008 & 1358.5 & 30641.3 & 0 & 3.00/3 \\
\code{p3--CCp3--p1} & triang. no estr. & 0.1 & 2.029 & 2.204 & 2.046 & 2.046 & 2.046 & 2.046 & 1.999 & 2.014 & 2.008 & 2.008 & 2.007 & 2.007 & 1550.3 & 32104.3 & 0 & 3.00/3 \\
\hline
\code{p3--CCp3--p2} & hexa & 0.0 & 3.538 & 3.993 & 3.290 & 4.231 & 3.304 & 3.441 & 4.323 & 4.279 & 4.079 & 4.284 & 4.086 & 4.100 & 333.92 & 23866.2 & 0 & 3.00/3 \\
\code{p3--CCp3--p2} & hexa & 0.1 & 3.569 & 4.026 & 3.319 & 4.231 & 3.333 & 3.485 & 4.299 & 4.271 & 4.081 & 4.283 & 4.087 & 4.101 & 389.68 & 25377.7 & 0 & 3.00/3 \\
\code{p3--CCp3--p2} & triang. no estr. & 0.0 & 4.058 & 4.096 & 3.846 & 4.109 & 3.862 & 3.934 & 4.476 & 4.165 & 4.252 & 4.049 & 4.256 & 4.239 & 1394.4 & 30860.8 & 0 & 3.00/3 \\
\code{p3--CCp3--p2} & triang. no estr. & 0.1 & 4.078 & 4.105 & 3.884 & 4.108 & 3.903 & 3.988 & 4.460 & 4.139 & 4.200 & 4.049 & 4.211 & 4.188 & 1577.0 & 32238.8 & 0 & 3.00/3 \\
\hline
\code{p3--CCp3--p3} & hexa & 0.0 & 3.538 & 3.993 & 3.290 & 4.241 & 3.304 & 3.438 & 4.323 & 4.279 & 4.079 & 4.302 & 4.086 & 4.100 & 367.11 & 23951.7 & 0 & 3.00/3 \\
\code{p3--CCp3--p3} & hexa & 0.1 & 3.569 & 4.026 & 3.319 & 4.241 & 3.333 & 3.481 & 4.299 & 4.271 & 4.081 & 4.300 & 4.087 & 4.100 & 393.42 & 25437.6 & 0 & 3.00/3 \\
\code{p3--CCp3--p3} & triang. no estr. & 0.0 & 4.057 & 4.096 & 3.846 & 4.114 & 3.862 & 3.932 & 4.477 & 4.165 & 4.252 & 4.050 & 4.257 & 4.240 & 1419.5 & 30971.1 & 0 & 3.00/3 \\
\code{p3--CCp3--p3} & triang. no estr. & 0.1 & 4.077 & 4.105 & 3.884 & 4.114 & 3.903 & 3.986 & 4.460 & 4.139 & 4.200 & 4.050 & 4.211 & 4.189 & 1585.1 & 32323.1 & 0 & 3.00/3 \\
\hline
\code{p3--GP--p1} & hexa & 0.0 & 2.187 & 2.637 & 2.001 & 2.598 & 1.989 & 2.596 & 2.020 & 2.073 & 2.000 & 2.068 & 1.999 & 2.066 & 266.14 & 21856.4 & 0 & 3.00/3 \\
\code{p3--GP--p1} & hexa & 0.1 & 2.161 & 2.631 & 2.001 & 2.592 & 1.988 & 2.589 & 2.016 & 2.069 & 2.000 & 2.060 & 1.999 & 2.059 & 237.47 & 24532.6 & 0 & 3.00/3 \\
\code{p3--GP--p1} & triang. no estr. & 0.0 & 2.041 & 2.209 & 2.061 & 2.061 & 2.060 & 2.060 & 1.999 & 2.015 & 2.008 & 2.008 & 2.008 & 2.008 & 1040.7 & 27557.2 & 0 & 3.00/3 \\
\code{p3--GP--p1} & triang. no estr. & 0.1 & 2.029 & 2.204 & 2.046 & 2.046 & 2.046 & 2.046 & 1.999 & 2.014 & 2.008 & 2.008 & 2.007 & 2.007 & 1082.2 & 31544.0 & 0 & 3.00/3 \\
\hline
\code{p3--GP--p2} & hexa & 0.0 & 3.538 & 3.992 & 1.999 & 4.003 & 1.997 & 4.010 & 4.323 & 4.279 & 2.000 & 4.194 & 2.001 & 4.195 & 256.36 & 21952.8 & 0 & 3.00/3 \\
\code{p3--GP--p2} & hexa & 0.1 & 3.569 & 4.026 & 1.999 & 4.033 & 1.997 & 4.040 & 4.299 & 4.271 & 2.000 & 4.212 & 2.000 & 4.213 & 293.86 & 24648.2 & 0 & 3.00/3 \\
\code{p3--GP--p2} & triang. no estr. & 0.0 & 4.057 & 4.096 & 1.992 & 4.077 & 2.002 & 4.083 & 4.476 & 4.165 & 1.998 & 4.204 & 2.003 & 4.205 & 1039.0 & 27685.6 & 0 & 3.00/3 \\
\code{p3--GP--p2} & triang. no estr. & 0.1 & 4.077 & 4.105 & 1.992 & 4.094 & 2.001 & 4.098 & 4.460 & 4.139 & 1.998 & 4.195 & 2.003 & 4.198 & 1148.6 & 31684.3 & 0 & 3.00/3 \\
\hline
\code{p3--GP--p3} & hexa & 0.0 & 3.538 & 3.992 & 1.999 & 4.000 & 1.998 & 4.008 & 4.323 & 4.279 & 2.000 & 4.198 & 2.001 & 4.200 & 311.89 & 22022.3 & 0 & 3.00/3 \\
\code{p3--GP--p3} & hexa & 0.1 & 3.569 & 4.026 & 1.999 & 4.031 & 1.997 & 4.039 & 4.299 & 4.271 & 2.000 & 4.217 & 2.001 & 4.219 & 334.26 & 24706.1 & 0 & 3.00/3 \\
\code{p3--GP--p3} & triang. no estr. & 0.0 & 4.057 & 4.096 & 1.992 & 4.079 & 2.002 & 4.086 & 4.476 & 4.165 & 1.998 & 4.227 & 2.003 & 4.229 & 1166.7 & 27749.1 & 0 & 3.00/3 \\
\code{p3--GP--p3} & triang. no estr. & 0.1 & 4.077 & 4.105 & 1.992 & 4.098 & 2.002 & 4.103 & 4.460 & 4.139 & 1.998 & 4.219 & 2.003 & 4.223 & 1198.0 & 31741.2 & 0 & 3.00/3 \\
\bottomrule
\end{longtable}
\normalsize


\subsubsection{\code{diagonalIion}}
\paragraph{$\Delta t=10^{-2}$}
\tiny
\setlength{\tabcolsep}{1.4pt}
\begin{longtable}{lllrrrrrrrrrrrrrrrr}
\caption{Resultados manufacturados 2D para \code{diagonalIion}, \code{crankNicolson}, masa \code{consistent}, estados \code{RKF45} y $\Delta t=10^{-2}$. El primer bloque de convergencia se ajusta sobre $N=10$--$100$; el segundo sobre $N=70$--$100$. La primera columna registra la terna \code{Vm-states-Iion} efectivamente usada en la corrida. \textbf{RB} es la cantidad de pasos de tiempo que agotaron el presupuesto de iteraciones no lineales, acumulada sobre los 91 tamaños de malla de la fila (455 pasos de tiempo en total por fila). Con \code{nonlinearAcceptUnconverged=false} un paso no convergido se revierte a los valores del inicio del paso mientras el reloj avanza igual. \textbf{$n_{\mathrm{nl}}$} reporta la mediana y el máximo de iteraciones no lineales por paso sobre esos mismos pasos; el presupuesto configurado es \code{nonlinearIterations}$=2000$.}\label{tab:es-results-2d-diagonaliion-1p0em02}\\
\toprule
& & & \multicolumn{6}{c}{Ajuste $N=10$--$100$} & \multicolumn{6}{c}{Ajuste $N=70$--$100$} & & & & \\
\cmidrule(lr){4-9}\cmidrule(lr){10-15}
\code{Vm-states-Iion} & Tipo de malla & $\alpha$ & $V_m$ & $V_m^G$ & $u_1$ & $u_1^G$ & $u_2$ & $u_2^G$ & $V_m$ & $V_m^G$ & $u_1$ & $u_1^G$ & $u_2$ & $u_2^G$ & $t_{\Sigma}$ [s] & RSS$_{\Sigma}$ [MB] & RB & $n_{\mathrm{nl}}$ \\
\midrule
\endfirsthead
\toprule
& & & \multicolumn{6}{c}{Ajuste $N=10$--$100$} & \multicolumn{6}{c}{Ajuste $N=70$--$100$} & & & & \\
\cmidrule(lr){4-9}\cmidrule(lr){10-15}
\code{Vm-states-Iion} & Tipo de malla & $\alpha$ & $V_m$ & $V_m^G$ & $u_1$ & $u_1^G$ & $u_2$ & $u_2^G$ & $V_m$ & $V_m^G$ & $u_1$ & $u_1^G$ & $u_2$ & $u_2^G$ & $t_{\Sigma}$ [s] & RSS$_{\Sigma}$ [MB] & RB & $n_{\mathrm{nl}}$ \\
\midrule
\endhead
\code{NO--na--NO} & hexa & 0.0 & 1.992 & 1.992 & 2.135 & 2.135 & 2.134 & 2.134 & 1.985 & 1.985 & 2.500 & 2.500 & 2.497 & 2.497 & 10.461 & 16447.3 & 0 & 3.00/3 \\
\code{NO--na--NO} & hexa & 0.1 & 1.849 & 1.849 & 1.612 & 1.612 & 1.618 & 1.618 & 2.278 & 2.278 & 2.159 & 2.159 & 2.163 & 2.163 & 34.701 & 18361.8 & 0 & 3.00/3 \\
\code{NO--na--NO} & triang. no estr. & 0.0 & 0.730 & 0.730 & 0.778 & 0.778 & 0.784 & 0.784 & 0.862 & 0.862 & 1.037 & 1.037 & 1.107 & 1.107 & 85.422 & 17192.2 & 0 & 3.00/3 \\
\code{NO--na--NO} & triang. no estr. & 0.1 & 0.722 & 0.722 & 0.772 & 0.772 & 0.777 & 0.777 & 0.871 & 0.871 & 1.027 & 1.027 & 1.097 & 1.097 & 158.63 & 20812.7 & 0 & 3.00/3 \\
\hline
\code{NO--CCp1--p1} & hexa & 0.0 & 1.993 & 1.993 & 2.136 & 2.354 & 2.134 & 2.333 & 1.985 & 1.985 & 2.504 & 2.289 & 2.497 & 2.242 & 26.905 & 16918.2 & 0 & 3.00/3 \\
\code{NO--CCp1--p1} & hexa & 0.1 & 1.849 & 1.849 & 1.611 & 2.354 & 1.618 & 2.310 & 2.278 & 2.278 & 2.159 & 2.289 & 2.163 & 2.225 & 53.309 & 18484.4 & 0 & 3.00/3 \\
\code{NO--CCp1--p1} & triang. no estr. & 0.0 & 0.730 & 0.730 & 0.778 & 0.778 & 0.784 & 0.784 & 0.862 & 0.862 & 1.037 & 1.037 & 1.107 & 1.107 & 152.08 & 17929.7 & 0 & 3.00/3 \\
\code{NO--CCp1--p1} & triang. no estr. & 0.1 & 0.722 & 0.722 & 0.772 & 0.772 & 0.777 & 0.777 & 0.871 & 0.871 & 1.027 & 1.027 & 1.097 & 1.097 & 202.40 & 20927.9 & 0 & 3.00/3 \\
\hline
\code{NO--CCp1--p2} & hexa & 0.0 & 1.993 & 1.993 & 2.139 & 2.229 & 2.134 & 2.184 & 1.985 & 1.985 & 2.513 & 2.150 & 2.495 & 2.094 & 36.302 & 17151.2 & 0 & 3.00/3 \\
\code{NO--CCp1--p2} & hexa & 0.1 & 1.849 & 1.849 & 1.609 & 2.229 & 1.619 & 2.181 & 2.278 & 2.278 & 2.159 & 2.149 & 2.163 & 2.091 & 56.018 & 18610.3 & 0 & 3.00/3 \\
\code{NO--CCp1--p2} & triang. no estr. & 0.0 & 0.730 & 0.730 & 0.778 & 2.111 & 0.784 & 1.164 & 0.862 & 0.862 & 1.037 & 2.101 & 1.107 & 1.150 & 128.42 & 18180.6 & 0 & 3.00/3 \\
\code{NO--CCp1--p2} & triang. no estr. & 0.1 & 0.722 & 0.722 & 0.772 & 2.112 & 0.777 & 1.172 & 0.871 & 0.871 & 1.027 & 2.105 & 1.097 & 1.143 & 193.06 & 21050.2 & 0 & 3.00/3 \\
\hline
\code{NO--CCp1--p3} & hexa & 0.0 & 1.993 & 1.993 & 2.139 & 2.247 & 2.134 & 2.211 & 1.985 & 1.985 & 2.513 & 2.167 & 2.495 & 2.112 & 39.682 & 17266.4 & 0 & 3.00/3 \\
\code{NO--CCp1--p3} & hexa & 0.1 & 1.849 & 1.849 & 1.609 & 2.247 & 1.619 & 2.206 & 2.278 & 2.278 & 2.159 & 2.166 & 2.163 & 2.109 & 53.827 & 18659.0 & 0 & 3.00/3 \\
\code{NO--CCp1--p3} & triang. no estr. & 0.0 & 0.730 & 0.730 & 0.778 & 2.139 & 0.784 & 1.132 & 0.862 & 0.862 & 1.037 & 2.131 & 1.107 & 1.144 & 158.11 & 18311.1 & 0 & 3.00/3 \\
\code{NO--CCp1--p3} & triang. no estr. & 0.1 & 0.722 & 0.722 & 0.772 & 2.140 & 0.777 & 1.139 & 0.871 & 0.871 & 1.027 & 2.135 & 1.097 & 1.136 & 210.31 & 21126.0 & 0 & 3.00/3 \\
\hline
\code{NO--CCp2--p1} & hexa & 0.0 & 1.992 & 1.992 & 2.141 & 3.042 & 2.132 & 2.624 & 1.985 & 1.985 & 2.517 & 2.999 & 2.491 & 2.603 & 40.777 & 17394.3 & 0 & 3.00/3 \\
\code{NO--CCp2--p1} & hexa & 0.1 & 1.849 & 1.849 & 1.607 & 3.038 & 1.618 & 1.948 & 2.278 & 2.278 & 2.158 & 2.987 & 2.162 & 2.187 & 57.797 & 18603.4 & 0 & 3.00/3 \\
\code{NO--CCp2--p1} & triang. no estr. & 0.0 & 0.730 & 0.730 & 0.778 & 0.778 & 0.784 & 0.784 & 0.862 & 0.862 & 1.037 & 1.037 & 1.107 & 1.107 & 173.84 & 19021.5 & 0 & 3.00/3 \\
\code{NO--CCp2--p1} & triang. no estr. & 0.1 & 0.722 & 0.722 & 0.772 & 0.772 & 0.777 & 0.777 & 0.871 & 0.871 & 1.027 & 1.027 & 1.097 & 1.097 & 225.85 & 21033.4 & 0 & 3.00/3 \\
\hline
\code{NO--CCp2--p2} & hexa & 0.0 & 1.991 & 1.991 & 2.145 & 3.063 & 2.120 & 2.764 & 1.986 & 1.986 & 2.519 & 3.002 & 2.444 & 2.629 & 42.048 & 17646.0 & 0 & 3.00/3 \\
\code{NO--CCp2--p2} & hexa & 0.1 & 1.851 & 1.851 & 1.600 & 3.061 & 1.624 & 2.068 & 2.274 & 2.274 & 2.155 & 2.996 & 2.162 & 2.210 & 59.158 & 18838.7 & 0 & 3.00/3 \\
\code{NO--CCp2--p2} & triang. no estr. & 0.0 & 0.730 & 0.730 & 0.778 & 1.941 & 0.784 & 0.776 & 0.862 & 0.862 & 1.037 & 1.095 & 1.107 & 1.097 & 162.86 & 19271.0 & 0 & 3.00/3 \\
\code{NO--CCp2--p2} & triang. no estr. & 0.1 & 0.722 & 0.722 & 0.772 & 1.957 & 0.777 & 0.773 & 0.871 & 0.871 & 1.027 & 1.090 & 1.097 & 1.087 & 250.97 & 21225.4 & 0 & 3.00/3 \\
\hline
\code{NO--CCp2--p3} & hexa & 0.0 & 1.991 & 1.991 & 2.145 & 3.058 & 2.120 & 2.763 & 1.986 & 1.986 & 2.519 & 3.001 & 2.444 & 2.632 & 40.253 & 17775.1 & 0 & 3.00/3 \\
\code{NO--CCp2--p3} & hexa & 0.1 & 1.851 & 1.851 & 1.600 & 3.056 & 1.624 & 2.069 & 2.274 & 2.274 & 2.155 & 2.995 & 2.162 & 2.211 & 57.600 & 18965.1 & 0 & 3.00/3 \\
\code{NO--CCp2--p3} & triang. no estr. & 0.0 & 0.730 & 0.730 & 0.778 & 1.952 & 0.784 & 0.776 & 0.862 & 0.862 & 1.037 & 1.098 & 1.107 & 1.096 & 171.94 & 19401.9 & 0 & 3.00/3 \\
\code{NO--CCp2--p3} & triang. no estr. & 0.1 & 0.722 & 0.722 & 0.772 & 1.968 & 0.777 & 0.773 & 0.871 & 0.871 & 1.027 & 1.093 & 1.097 & 1.087 & 226.38 & 21364.4 & 0 & 3.00/3 \\
\hline
\code{NO--CCp3--p1} & hexa & 0.0 & 1.992 & 1.992 & 2.139 & 4.063 & 2.132 & 2.165 & 1.985 & 1.985 & 2.517 & 3.183 & 2.491 & 2.491 & 55.598 & 18727.0 & 0 & 3.00/3 \\
\code{NO--CCp3--p1} & hexa & 0.1 & 1.849 & 1.849 & 1.605 & 3.118 & 1.618 & 1.661 & 2.278 & 2.278 & 2.158 & 2.340 & 2.163 & 2.189 & 77.383 & 19869.3 & 0 & 3.00/3 \\
\code{NO--CCp3--p1} & triang. no estr. & 0.0 & 0.730 & 0.730 & 0.778 & 0.778 & 0.784 & 0.784 & 0.862 & 0.862 & 1.037 & 1.037 & 1.107 & 1.107 & 206.33 & 21447.4 & 0 & 3.00/3 \\
\code{NO--CCp3--p1} & triang. no estr. & 0.1 & 0.722 & 0.722 & 0.772 & 0.772 & 0.777 & 0.777 & 0.871 & 0.871 & 1.027 & 1.027 & 1.097 & 1.097 & 261.62 & 23304.6 & 0 & 3.00/3 \\
\hline
\code{NO--CCp3--p2} & hexa & 0.0 & 1.992 & 1.992 & 2.141 & 4.277 & 2.120 & 2.216 & 1.985 & 1.985 & 2.519 & 4.233 & 2.444 & 2.443 & 71.803 & 19000.5 & 0 & 3.00/3 \\
\code{NO--CCp3--p2} & hexa & 0.1 & 1.851 & 1.851 & 1.594 & 3.421 & 1.622 & 1.722 & 2.274 & 2.274 & 2.155 & 2.589 & 2.162 & 2.221 & 75.873 & 20100.4 & 0 & 3.00/3 \\
\code{NO--CCp3--p2} & triang. no estr. & 0.0 & 0.730 & 0.730 & 0.778 & 1.071 & 0.784 & 0.796 & 0.862 & 0.862 & 1.037 & 1.037 & 1.107 & 1.109 & 214.00 & 21742.7 & 0 & 3.00/3 \\
\code{NO--CCp3--p2} & triang. no estr. & 0.1 & 0.722 & 0.722 & 0.772 & 1.081 & 0.777 & 0.798 & 0.871 & 0.871 & 1.027 & 1.028 & 1.097 & 1.100 & 282.58 & 23548.4 & 0 & 3.00/3 \\
\hline
\code{NO--CCp3--p3} & hexa & 0.0 & 1.992 & 1.992 & 2.141 & 4.293 & 2.120 & 2.214 & 1.985 & 1.985 & 2.519 & 4.254 & 2.444 & 2.443 & 60.643 & 19121.6 & 0 & 3.00/3 \\
\code{NO--CCp3--p3} & hexa & 0.1 & 1.851 & 1.851 & 1.594 & 3.409 & 1.622 & 1.721 & 2.274 & 2.274 & 2.155 & 2.578 & 2.162 & 2.221 & 67.099 & 20203.7 & 0 & 3.00/3 \\
\code{NO--CCp3--p3} & triang. no estr. & 0.0 & 0.730 & 0.730 & 0.778 & 1.054 & 0.784 & 0.796 & 0.862 & 0.862 & 1.037 & 1.037 & 1.107 & 1.109 & 243.59 & 21897.1 & 0 & 3.00/3 \\
\code{NO--CCp3--p3} & triang. no estr. & 0.1 & 0.722 & 0.722 & 0.772 & 1.064 & 0.777 & 0.797 & 0.871 & 0.871 & 1.027 & 1.028 & 1.097 & 1.100 & 198.26 & 23661.4 & 0 & 3.00/3 \\
\hline
\code{NO--GP--p1} & hexa & 0.0 & 1.992 & 1.992 & 2.000 & 2.057 & 2.005 & 2.061 & 1.985 & 1.985 & 2.002 & 2.031 & 2.023 & 2.031 & 13.699 & 16534.5 & 0 & 3.00/3 \\
\code{NO--GP--p1} & hexa & 0.1 & 1.849 & 1.849 & 2.000 & 2.049 & 1.997 & 2.053 & 2.278 & 2.278 & 2.002 & 2.039 & 2.030 & 2.039 & 45.175 & 18416.0 & 0 & 3.00/3 \\
\code{NO--GP--p1} & triang. no estr. & 0.0 & 0.730 & 0.730 & 0.778 & 0.778 & 0.784 & 0.784 & 0.862 & 0.862 & 1.037 & 1.037 & 1.107 & 1.107 & 87.018 & 17338.7 & 0 & 3.00/3 \\
\code{NO--GP--p1} & triang. no estr. & 0.1 & 0.722 & 0.722 & 0.772 & 0.772 & 0.777 & 0.777 & 0.871 & 0.871 & 1.027 & 1.027 & 1.097 & 1.097 & 165.88 & 20895.6 & 0 & 3.00/3 \\
\hline
\code{NO--GP--p2} & hexa & 0.0 & 1.992 & 1.992 & 1.999 & 2.015 & 2.002 & 2.017 & 1.986 & 1.986 & 2.001 & 2.010 & 2.011 & 2.010 & 28.704 & 16813.7 & 0 & 3.00/3 \\
\code{NO--GP--p2} & hexa & 0.1 & 1.851 & 1.851 & 1.999 & 2.015 & 2.001 & 2.017 & 2.274 & 2.274 & 2.001 & 2.011 & 2.013 & 2.011 & 52.057 & 18561.0 & 0 & 3.00/3 \\
\code{NO--GP--p2} & triang. no estr. & 0.0 & 0.730 & 0.730 & 1.878 & 1.418 & 1.102 & 1.411 & 0.862 & 0.862 & 1.678 & 1.219 & 1.112 & 1.269 & 133.70 & 17598.5 & 0 & 3.00/3 \\
\code{NO--GP--p2} & triang. no estr. & 0.1 & 0.722 & 0.722 & 1.880 & 1.427 & 1.107 & 1.420 & 0.871 & 0.871 & 1.682 & 1.220 & 1.104 & 1.270 & 133.85 & 21034.4 & 0 & 3.00/3 \\
\hline
\code{NO--GP--p3} & hexa & 0.0 & 1.992 & 1.992 & 1.999 & 2.019 & 2.002 & 2.021 & 1.986 & 1.986 & 2.001 & 2.011 & 2.011 & 2.011 & 47.744 & 16969.7 & 0 & 3.00/3 \\
\code{NO--GP--p3} & hexa & 0.1 & 1.851 & 1.851 & 1.999 & 2.019 & 2.001 & 2.021 & 2.274 & 2.274 & 2.001 & 2.013 & 2.013 & 2.013 & 50.120 & 18613.3 & 0 & 3.00/3 \\
\code{NO--GP--p3} & triang. no estr. & 0.0 & 0.730 & 0.730 & 1.878 & 1.389 & 1.102 & 1.384 & 0.862 & 0.862 & 1.678 & 1.206 & 1.112 & 1.258 & 71.347 & 17710.6 & 0 & 3.00/3 \\
\code{NO--GP--p3} & triang. no estr. & 0.1 & 0.722 & 0.722 & 1.881 & 1.398 & 1.107 & 1.392 & 0.871 & 0.871 & 1.682 & 1.207 & 1.104 & 1.259 & 143.32 & 21099.4 & 0 & 3.00/3 \\
\hline
\code{p1--na--NO} & hexa & 0.0 & 1.911 & 2.146 & 1.616 & 1.616 & 1.626 & 1.626 & 2.612 & 2.168 & 2.381 & 2.381 & 2.387 & 2.387 & 80.953 & 17858.1 & 0 & 3.00/3 \\
\code{p1--na--NO} & hexa & 0.1 & 1.971 & 2.163 & 1.677 & 1.677 & 1.687 & 1.687 & 2.614 & 2.141 & 2.416 & 2.416 & 2.422 & 2.422 & 94.616 & 19059.8 & 0 & 3.00/3 \\
\code{p1--na--NO} & triang. no estr. & 0.0 & 2.271 & 2.126 & 2.004 & 2.004 & 2.002 & 2.002 & 2.384 & 2.115 & 2.552 & 2.552 & 2.549 & 2.549 & 590.16 & 19872.9 & 0 & 3.00/3 \\
\code{p1--na--NO} & triang. no estr. & 0.1 & 2.301 & 2.106 & 2.092 & 2.092 & 2.090 & 2.090 & 2.222 & 2.054 & 2.524 & 2.524 & 2.524 & 2.524 & 559.07 & 21538.7 & 0 & 3.00/3 \\
\hline
\code{p1--CCp1--p1} & hexa & 0.0 & 1.911 & 2.146 & 1.616 & 2.349 & 1.626 & 1.955 & 2.612 & 2.168 & 2.381 & 2.294 & 2.387 & 2.369 & 91.258 & 18293.1 & 0 & 3.00/3 \\
\code{p1--CCp1--p1} & hexa & 0.1 & 1.971 & 2.163 & 1.677 & 2.351 & 1.687 & 2.045 & 2.614 & 2.141 & 2.416 & 2.293 & 2.422 & 2.383 & 104.24 & 19267.6 & 0 & 3.00/3 \\
\code{p1--CCp1--p1} & triang. no estr. & 0.0 & 2.271 & 2.126 & 2.004 & 2.004 & 2.002 & 2.002 & 2.384 & 2.115 & 2.552 & 2.552 & 2.549 & 2.549 & 623.89 & 20378.7 & 0 & 3.00/3 \\
\code{p1--CCp1--p1} & triang. no estr. & 0.1 & 2.301 & 2.106 & 2.092 & 2.092 & 2.090 & 2.090 & 2.222 & 2.054 & 2.524 & 2.524 & 2.524 & 2.524 & 591.30 & 21752.6 & 0 & 3.00/3 \\
\hline
\code{p1--CCp1--p2} & hexa & 0.0 & 1.911 & 2.146 & 1.616 & 2.228 & 1.626 & 2.057 & 2.612 & 2.168 & 2.381 & 2.153 & 2.387 & 2.232 & 91.419 & 18505.1 & 0 & 3.00/3 \\
\code{p1--CCp1--p2} & hexa & 0.1 & 1.971 & 2.163 & 1.676 & 2.228 & 1.687 & 2.104 & 2.614 & 2.141 & 2.416 & 2.152 & 2.422 & 2.208 & 107.48 & 19395.6 & 0 & 3.00/3 \\
\code{p1--CCp1--p2} & triang. no estr. & 0.0 & 2.271 & 2.126 & 2.004 & 2.102 & 2.002 & 2.036 & 2.384 & 2.115 & 2.552 & 2.063 & 2.549 & 2.285 & 611.35 & 20633.4 & 0 & 3.00/3 \\
\code{p1--CCp1--p2} & triang. no estr. & 0.1 & 2.302 & 2.106 & 2.092 & 2.102 & 2.090 & 2.066 & 2.222 & 2.054 & 2.524 & 2.061 & 2.524 & 2.190 & 595.05 & 21917.4 & 0 & 3.00/3 \\
\hline
\code{p1--CCp1--p3} & hexa & 0.0 & 1.911 & 2.146 & 1.616 & 2.246 & 1.626 & 2.063 & 2.612 & 2.168 & 2.381 & 2.170 & 2.387 & 2.255 & 97.575 & 18602.7 & 0 & 3.00/3 \\
\code{p1--CCp1--p3} & hexa & 0.1 & 1.971 & 2.163 & 1.676 & 2.246 & 1.687 & 2.115 & 2.614 & 2.141 & 2.416 & 2.169 & 2.422 & 2.234 & 111.94 & 19454.3 & 0 & 3.00/3 \\
\code{p1--CCp1--p3} & triang. no estr. & 0.0 & 2.271 & 2.126 & 2.004 & 2.127 & 2.002 & 2.040 & 2.384 & 2.115 & 2.552 & 2.080 & 2.549 & 2.315 & 608.61 & 20768.4 & 0 & 3.00/3 \\
\code{p1--CCp1--p3} & triang. no estr. & 0.1 & 2.302 & 2.106 & 2.092 & 2.127 & 2.090 & 2.075 & 2.222 & 2.054 & 2.524 & 2.078 & 2.524 & 2.217 & 593.84 & 21991.6 & 0 & 3.00/3 \\
\hline
\code{p1--CCp2--p1} & hexa & 0.0 & 1.911 & 2.146 & 1.616 & 2.904 & 1.626 & 1.665 & 2.612 & 2.168 & 2.381 & 2.783 & 2.387 & 2.395 & 109.67 & 18769.5 & 0 & 3.00/3 \\
\code{p1--CCp2--p1} & hexa & 0.1 & 1.971 & 2.163 & 1.676 & 2.954 & 1.687 & 1.736 & 2.614 & 2.141 & 2.416 & 2.859 & 2.422 & 2.431 & 122.16 & 19370.8 & 0 & 3.00/3 \\
\code{p1--CCp2--p1} & triang. no estr. & 0.0 & 2.271 & 2.126 & 2.004 & 2.004 & 2.002 & 2.002 & 2.384 & 2.115 & 2.552 & 2.552 & 2.549 & 2.549 & 661.75 & 21487.3 & 0 & 3.00/3 \\
\code{p1--CCp2--p1} & triang. no estr. & 0.1 & 2.301 & 2.106 & 2.092 & 2.092 & 2.090 & 2.090 & 2.222 & 2.054 & 2.524 & 2.524 & 2.524 & 2.524 & 625.71 & 22098.7 & 0 & 3.00/3 \\
\hline
\code{p1--CCp2--p2} & hexa & 0.0 & 1.912 & 2.146 & 1.616 & 2.990 & 1.627 & 1.688 & 2.612 & 2.168 & 2.381 & 2.882 & 2.387 & 2.400 & 120.42 & 18967.0 & 0 & 3.00/3 \\
\code{p1--CCp2--p2} & hexa & 0.1 & 1.971 & 2.163 & 1.676 & 3.018 & 1.687 & 1.765 & 2.613 & 2.141 & 2.416 & 2.931 & 2.422 & 2.437 & 122.15 & 19561.4 & 0 & 3.00/3 \\
\code{p1--CCp2--p2} & triang. no estr. & 0.0 & 2.271 & 2.126 & 2.004 & 2.677 & 2.002 & 1.977 & 2.384 & 2.115 & 2.552 & 2.669 & 2.549 & 2.528 & 649.54 & 21723.7 & 0 & 3.00/3 \\
\code{p1--CCp2--p2} & triang. no estr. & 0.1 & 2.301 & 2.106 & 2.092 & 2.799 & 2.090 & 2.062 & 2.222 & 2.054 & 2.524 & 2.735 & 2.523 & 2.511 & 631.25 & 22285.7 & 0 & 3.00/3 \\
\hline
\code{p1--CCp2--p3} & hexa & 0.0 & 1.912 & 2.146 & 1.616 & 2.987 & 1.627 & 1.688 & 2.612 & 2.168 & 2.381 & 2.884 & 2.387 & 2.400 & 121.12 & 19064.7 & 0 & 3.00/3 \\
\code{p1--CCp2--p3} & hexa & 0.1 & 1.971 & 2.163 & 1.676 & 3.015 & 1.687 & 1.765 & 2.613 & 2.141 & 2.416 & 2.932 & 2.422 & 2.437 & 117.46 & 19617.3 & 0 & 3.00/3 \\
\code{p1--CCp2--p3} & triang. no estr. & 0.0 & 2.271 & 2.126 & 2.004 & 2.681 & 2.002 & 1.977 & 2.384 & 2.115 & 2.552 & 2.672 & 2.549 & 2.527 & 650.20 & 21848.2 & 0 & 3.00/3 \\
\code{p1--CCp2--p3} & triang. no estr. & 0.1 & 2.301 & 2.106 & 2.092 & 2.802 & 2.090 & 2.062 & 2.222 & 2.054 & 2.524 & 2.740 & 2.523 & 2.510 & 629.57 & 22409.5 & 0 & 3.00/3 \\
\hline
\code{p1--CCp3--p1} & hexa & 0.0 & 1.911 & 2.146 & 1.616 & 2.034 & 1.626 & 1.630 & 2.612 & 2.168 & 2.381 & 2.378 & 2.387 & 2.407 & 148.41 & 20023.7 & 0 & 3.00/3 \\
\code{p1--CCp3--p1} & hexa & 0.1 & 1.971 & 2.163 & 1.677 & 2.165 & 1.687 & 1.692 & 2.614 & 2.141 & 2.416 & 2.409 & 2.422 & 2.443 & 137.46 & 20539.4 & 0 & 3.00/3 \\
\code{p1--CCp3--p1} & triang. no estr. & 0.0 & 2.271 & 2.126 & 2.004 & 2.004 & 2.002 & 2.002 & 2.384 & 2.115 & 2.552 & 2.552 & 2.549 & 2.549 & 694.57 & 23794.5 & 0 & 3.00/3 \\
\code{p1--CCp3--p1} & triang. no estr. & 0.1 & 2.301 & 2.106 & 2.092 & 2.092 & 2.090 & 2.090 & 2.222 & 2.054 & 2.524 & 2.524 & 2.524 & 2.524 & 677.84 & 24272.0 & 0 & 3.00/3 \\
\hline
\code{p1--CCp3--p2} & hexa & 0.0 & 1.912 & 2.146 & 1.617 & 2.239 & 1.627 & 1.634 & 2.612 & 2.168 & 2.381 & 2.377 & 2.387 & 2.426 & 135.39 & 20235.5 & 0 & 3.00/3 \\
\code{p1--CCp3--p2} & hexa & 0.1 & 1.971 & 2.163 & 1.677 & 2.394 & 1.688 & 1.698 & 2.613 & 2.141 & 2.416 & 2.402 & 2.422 & 2.464 & 139.88 & 20752.4 & 0 & 3.00/3 \\
\code{p1--CCp3--p2} & triang. no estr. & 0.0 & 2.271 & 2.126 & 2.004 & 2.031 & 2.002 & 1.964 & 2.384 & 2.115 & 2.552 & 2.527 & 2.549 & 2.528 & 738.52 & 24032.2 & 0 & 3.00/3 \\
\code{p1--CCp3--p2} & triang. no estr. & 0.1 & 2.301 & 2.106 & 2.092 & 2.141 & 2.090 & 2.052 & 2.222 & 2.054 & 2.524 & 2.511 & 2.523 & 2.516 & 676.05 & 24507.1 & 0 & 3.00/3 \\
\hline
\code{p1--CCp3--p3} & hexa & 0.0 & 1.912 & 2.146 & 1.617 & 2.225 & 1.627 & 1.634 & 2.612 & 2.168 & 2.381 & 2.376 & 2.387 & 2.426 & 146.52 & 20334.8 & 0 & 3.00/3 \\
\code{p1--CCp3--p3} & hexa & 0.1 & 1.971 & 2.163 & 1.677 & 2.379 & 1.688 & 1.697 & 2.613 & 2.141 & 2.416 & 2.402 & 2.422 & 2.463 & 136.92 & 20848.2 & 0 & 3.00/3 \\
\code{p1--CCp3--p3} & triang. no estr. & 0.0 & 2.271 & 2.126 & 2.004 & 2.027 & 2.002 & 1.964 & 2.384 & 2.115 & 2.552 & 2.527 & 2.549 & 2.528 & 771.47 & 24158.5 & 0 & 3.00/3 \\
\code{p1--CCp3--p3} & triang. no estr. & 0.1 & 2.301 & 2.106 & 2.092 & 2.136 & 2.090 & 2.052 & 2.222 & 2.054 & 2.524 & 2.512 & 2.523 & 2.516 & 663.06 & 24640.8 & 0 & 3.00/3 \\
\hline
\code{p1--GP--p1} & hexa & 0.0 & 1.911 & 2.146 & 1.998 & 2.154 & 1.838 & 2.160 & 2.612 & 2.168 & 2.007 & 2.316 & 2.258 & 2.316 & 85.742 & 17904.5 & 0 & 3.00/3 \\
\code{p1--GP--p1} & hexa & 0.1 & 1.971 & 2.163 & 1.999 & 2.221 & 1.898 & 2.227 & 2.614 & 2.141 & 2.005 & 2.319 & 2.234 & 2.318 & 106.80 & 19111.6 & 0 & 3.00/3 \\
\code{p1--GP--p1} & triang. no estr. & 0.0 & 2.271 & 2.126 & 2.004 & 2.004 & 2.002 & 2.002 & 2.384 & 2.115 & 2.552 & 2.552 & 2.549 & 2.549 & 598.75 & 19898.8 & 0 & 3.00/3 \\
\code{p1--GP--p1} & triang. no estr. & 0.1 & 2.301 & 2.106 & 2.092 & 2.092 & 2.090 & 2.090 & 2.222 & 2.054 & 2.524 & 2.524 & 2.524 & 2.524 & 560.43 & 21570.5 & 0 & 3.00/3 \\
\hline
\code{p1--GP--p2} & hexa & 0.0 & 1.912 & 2.146 & 1.999 & 2.182 & 1.951 & 2.181 & 2.612 & 2.168 & 2.002 & 2.161 & 2.105 & 2.157 & 101.45 & 18110.1 & 0 & 3.00/3 \\
\code{p1--GP--p2} & hexa & 0.1 & 1.971 & 2.163 & 1.999 & 2.201 & 1.975 & 2.199 & 2.613 & 2.141 & 2.001 & 2.156 & 2.080 & 2.151 & 122.27 & 19231.4 & 0 & 3.00/3 \\
\code{p1--GP--p2} & triang. no estr. & 0.0 & 2.271 & 2.126 & 1.992 & 2.069 & 2.012 & 2.068 & 2.384 & 2.115 & 2.004 & 2.184 & 2.284 & 2.176 & 625.13 & 20140.2 & 0 & 3.00/3 \\
\code{p1--GP--p2} & triang. no estr. & 0.1 & 2.301 & 2.106 & 1.992 & 2.089 & 2.038 & 2.087 & 2.222 & 2.054 & 2.002 & 2.117 & 2.186 & 2.113 & 570.64 & 21716.5 & 0 & 3.00/3 \\
\hline
\code{p1--GP--p3} & hexa & 0.0 & 1.912 & 2.146 & 1.999 & 2.196 & 1.951 & 2.195 & 2.612 & 2.168 & 2.002 & 2.177 & 2.105 & 2.173 & 108.74 & 18227.1 & 0 & 3.00/3 \\
\code{p1--GP--p3} & hexa & 0.1 & 1.971 & 2.163 & 1.999 & 2.217 & 1.975 & 2.215 & 2.613 & 2.141 & 2.001 & 2.172 & 2.080 & 2.167 & 117.87 & 19282.9 & 0 & 3.00/3 \\
\code{p1--GP--p3} & triang. no estr. & 0.0 & 2.271 & 2.126 & 1.992 & 2.082 & 2.012 & 2.081 & 2.384 & 2.115 & 2.004 & 2.219 & 2.284 & 2.210 & 663.03 & 20261.8 & 0 & 3.00/3 \\
\code{p1--GP--p3} & triang. no estr. & 0.1 & 2.301 & 2.106 & 1.992 & 2.109 & 2.038 & 2.106 & 2.222 & 2.054 & 2.002 & 2.144 & 2.186 & 2.139 & 563.90 & 21789.6 & 0 & 3.00/3 \\
\hline
\code{p2--na--NO} & hexa & 0.0 & 2.127 & 2.567 & 2.176 & 2.176 & 2.176 & 2.176 & 2.016 & 2.179 & 2.139 & 2.139 & 2.139 & 2.139 & 98.709 & 18691.7 & 0 & 3.00/3 \\
\code{p2--na--NO} & hexa & 0.1 & 2.128 & 2.607 & 2.180 & 2.180 & 2.181 & 2.181 & 2.018 & 2.209 & 2.153 & 2.153 & 2.153 & 2.153 & 115.37 & 20296.1 & 0 & 3.00/3 \\
\code{p2--na--NO} & triang. no estr. & 0.0 & 2.056 & 2.374 & 2.166 & 2.166 & 2.161 & 2.161 & 1.982 & 2.068 & 2.260 & 2.260 & 2.263 & 2.263 & 620.05 & 21912.9 & 0 & 3.00/3 \\
\code{p2--na--NO} & triang. no estr. & 0.1 & 2.043 & 2.402 & 2.162 & 2.162 & 2.158 & 2.158 & 1.980 & 2.081 & 2.286 & 2.286 & 2.285 & 2.285 & 587.89 & 24376.3 & 0 & 3.00/3 \\
\hline
\code{p2--CCp1--p1} & hexa & 0.0 & 2.149 & 2.608 & 2.201 & 2.356 & 2.203 & 2.337 & 2.019 & 2.209 & 2.154 & 2.292 & 2.154 & 2.244 & 117.40 & 19455.1 & 0 & 3.00/3 \\
\code{p2--CCp1--p1} & hexa & 0.1 & 2.151 & 2.651 & 2.208 & 2.356 & 2.211 & 2.338 & 2.021 & 2.247 & 2.171 & 2.291 & 2.172 & 2.246 & 119.24 & 20864.2 & 0 & 3.00/3 \\
\code{p2--CCp1--p1} & triang. no estr. & 0.0 & 2.056 & 2.374 & 2.166 & 2.166 & 2.161 & 2.161 & 1.982 & 2.068 & 2.260 & 2.260 & 2.263 & 2.263 & 727.82 & 22654.0 & 0 & 3.00/3 \\
\code{p2--CCp1--p1} & triang. no estr. & 0.1 & 2.043 & 2.402 & 2.162 & 2.162 & 2.158 & 2.158 & 1.980 & 2.081 & 2.286 & 2.286 & 2.285 & 2.285 & 615.73 & 24950.3 & 0 & 3.00/3 \\
\hline
\code{p2--CCp1--p2} & hexa & 0.0 & 2.212 & 2.708 & 2.272 & 2.230 & 2.283 & 2.188 & 2.026 & 2.303 & 2.199 & 2.150 & 2.206 & 2.095 & 123.45 & 19688.0 & 0 & 3.00/3 \\
\code{p2--CCp1--p2} & hexa & 0.1 & 2.223 & 2.759 & 2.291 & 2.229 & 2.304 & 2.188 & 2.030 & 2.364 & 2.229 & 2.150 & 2.238 & 2.095 & 121.64 & 20990.5 & 0 & 3.00/3 \\
\code{p2--CCp1--p2} & triang. no estr. & 0.0 & 2.118 & 2.532 & 2.273 & 2.101 & 2.279 & 2.040 & 1.977 & 2.149 & 2.394 & 2.057 & 2.419 & 2.018 & 726.78 & 22876.9 & 0 & 3.00/3 \\
\code{p2--CCp1--p2} & triang. no estr. & 0.1 & 2.104 & 2.578 & 2.283 & 2.101 & 2.292 & 2.039 & 1.972 & 2.185 & 2.455 & 2.057 & 2.479 & 2.018 & 596.61 & 25104.2 & 0 & 3.00/3 \\
\hline
\code{p2--CCp1--p3} & hexa & 0.0 & 2.212 & 2.708 & 2.272 & 2.248 & 2.283 & 2.215 & 2.026 & 2.303 & 2.199 & 2.167 & 2.206 & 2.114 & 132.77 & 19802.8 & 0 & 3.00/3 \\
\code{p2--CCp1--p3} & hexa & 0.1 & 2.223 & 2.759 & 2.291 & 2.248 & 2.304 & 2.214 & 2.030 & 2.364 & 2.229 & 2.167 & 2.238 & 2.114 & 120.06 & 21058.9 & 0 & 3.00/3 \\
\code{p2--CCp1--p3} & triang. no estr. & 0.0 & 2.118 & 2.532 & 2.273 & 2.126 & 2.279 & 2.050 & 1.977 & 2.149 & 2.394 & 2.074 & 2.419 & 2.023 & 777.40 & 22967.1 & 0 & 3.00/3 \\
\code{p2--CCp1--p3} & triang. no estr. & 0.1 & 2.104 & 2.578 & 2.283 & 2.126 & 2.292 & 2.050 & 1.972 & 2.185 & 2.455 & 2.074 & 2.479 & 2.023 & 627.52 & 25163.3 & 0 & 3.00/3 \\
\hline
\code{p2--CCp2--p1} & hexa & 0.0 & 2.149 & 2.608 & 2.203 & 3.039 & 2.203 & 2.289 & 2.019 & 2.210 & 2.155 & 2.989 & 2.155 & 2.169 & 140.24 & 19917.8 & 0 & 3.00/3 \\
\code{p2--CCp2--p1} & hexa & 0.1 & 2.152 & 2.651 & 2.211 & 3.040 & 2.211 & 2.310 & 2.021 & 2.247 & 2.173 & 2.991 & 2.172 & 2.189 & 125.89 & 20909.1 & 0 & 3.00/3 \\
\code{p2--CCp2--p1} & triang. no estr. & 0.0 & 2.056 & 2.374 & 2.166 & 2.166 & 2.161 & 2.161 & 1.982 & 2.068 & 2.260 & 2.260 & 2.263 & 2.263 & 754.03 & 23653.5 & 0 & 3.00/3 \\
\code{p2--CCp2--p1} & triang. no estr. & 0.1 & 2.043 & 2.402 & 2.162 & 2.162 & 2.158 & 2.158 & 1.980 & 2.081 & 2.286 & 2.286 & 2.285 & 2.285 & 617.64 & 25105.4 & 0 & 3.00/3 \\
\hline
\code{p2--CCp2--p2} & hexa & 0.0 & 2.212 & 2.708 & 2.284 & 3.063 & 2.283 & 2.480 & 2.026 & 2.303 & 2.206 & 3.000 & 2.206 & 2.250 & 141.14 & 20114.9 & 0 & 3.00/3 \\
\code{p2--CCp2--p2} & hexa & 0.1 & 2.224 & 2.760 & 2.305 & 3.063 & 2.305 & 2.527 & 2.030 & 2.365 & 2.239 & 3.001 & 2.239 & 2.294 & 134.56 & 21038.7 & 0 & 3.00/3 \\
\code{p2--CCp2--p2} & triang. no estr. & 0.0 & 2.118 & 2.532 & 2.286 & 3.012 & 2.279 & 2.428 & 1.977 & 2.149 & 2.413 & 2.995 & 2.420 & 2.443 & 727.80 & 23875.1 & 0 & 3.00/3 \\
\code{p2--CCp2--p2} & triang. no estr. & 0.1 & 2.104 & 2.579 & 2.300 & 3.012 & 2.293 & 2.467 & 1.972 & 2.186 & 2.481 & 2.996 & 2.482 & 2.511 & 613.39 & 25257.0 & 0 & 3.00/3 \\
\hline
\code{p2--CCp2--p3} & hexa & 0.0 & 2.212 & 2.708 & 2.284 & 3.058 & 2.283 & 2.480 & 2.026 & 2.303 & 2.206 & 2.999 & 2.206 & 2.251 & 140.66 & 20243.0 & 0 & 3.00/3 \\
\code{p2--CCp2--p3} & hexa & 0.1 & 2.224 & 2.760 & 2.305 & 3.058 & 2.305 & 2.527 & 2.030 & 2.365 & 2.239 & 3.000 & 2.239 & 2.296 & 143.18 & 21098.7 & 0 & 3.00/3 \\
\code{p2--CCp2--p3} & triang. no estr. & 0.0 & 2.118 & 2.532 & 2.286 & 3.010 & 2.279 & 2.425 & 1.977 & 2.149 & 2.413 & 2.995 & 2.420 & 2.443 & 725.90 & 23969.0 & 0 & 3.00/3 \\
\code{p2--CCp2--p3} & triang. no estr. & 0.1 & 2.104 & 2.579 & 2.300 & 3.010 & 2.293 & 2.465 & 1.972 & 2.186 & 2.481 & 2.996 & 2.482 & 2.512 & 622.84 & 25335.9 & 0 & 3.00/3 \\
\hline
\code{p2--CCp3--p1} & hexa & 0.0 & 2.149 & 2.608 & 2.203 & 3.175 & 2.203 & 2.220 & 2.019 & 2.210 & 2.155 & 2.123 & 2.155 & 2.154 & 167.65 & 21118.3 & 0 & 3.00/3 \\
\code{p2--CCp3--p1} & hexa & 0.1 & 2.152 & 2.651 & 2.211 & 3.250 & 2.211 & 2.227 & 2.021 & 2.247 & 2.173 & 2.151 & 2.172 & 2.171 & 198.67 & 21530.3 & 0 & 3.00/3 \\
\code{p2--CCp3--p1} & triang. no estr. & 0.0 & 2.056 & 2.374 & 2.166 & 2.166 & 2.161 & 2.161 & 1.982 & 2.068 & 2.260 & 2.260 & 2.263 & 2.263 & 773.41 & 25977.5 & 0 & 3.00/3 \\
\code{p2--CCp3--p1} & triang. no estr. & 0.1 & 2.043 & 2.402 & 2.162 & 2.162 & 2.158 & 2.158 & 1.980 & 2.081 & 2.286 & 2.286 & 2.285 & 2.285 & 668.09 & 26830.7 & 0 & 3.00/3 \\
\hline
\code{p2--CCp3--p2} & hexa & 0.0 & 2.212 & 2.708 & 2.284 & 3.846 & 2.283 & 2.347 & 2.026 & 2.303 & 2.206 & 2.397 & 2.206 & 2.204 & 162.87 & 21297.1 & 0 & 3.00/3 \\
\code{p2--CCp3--p2} & hexa & 0.1 & 2.224 & 2.760 & 2.305 & 3.936 & 2.304 & 2.369 & 2.030 & 2.365 & 2.239 & 2.584 & 2.238 & 2.237 & 184.34 & 21724.7 & 0 & 3.00/3 \\
\code{p2--CCp3--p2} & triang. no estr. & 0.0 & 2.118 & 2.532 & 2.286 & 3.422 & 2.279 & 2.345 & 1.977 & 2.149 & 2.413 & 2.437 & 2.420 & 2.419 & 793.59 & 26150.4 & 0 & 3.00/3 \\
\code{p2--CCp3--p2} & triang. no estr. & 0.1 & 2.104 & 2.579 & 2.300 & 3.528 & 2.293 & 2.358 & 1.972 & 2.186 & 2.481 & 2.550 & 2.482 & 2.482 & 653.69 & 27025.5 & 0 & 3.00/3 \\
\hline
\code{p2--CCp3--p3} & hexa & 0.0 & 2.212 & 2.708 & 2.284 & 3.839 & 2.283 & 2.346 & 2.026 & 2.303 & 2.206 & 2.351 & 2.206 & 2.204 & 159.34 & 21389.3 & 0 & 3.00/3 \\
\code{p2--CCp3--p3} & hexa & 0.1 & 2.224 & 2.760 & 2.305 & 3.933 & 2.304 & 2.368 & 2.030 & 2.365 & 2.239 & 2.531 & 2.238 & 2.237 & 169.65 & 21784.7 & 0 & 3.00/3 \\
\code{p2--CCp3--p3} & triang. no estr. & 0.0 & 2.118 & 2.532 & 2.286 & 3.389 & 2.279 & 2.343 & 1.977 & 2.149 & 2.413 & 2.404 & 2.420 & 2.419 & 803.40 & 26267.3 & 0 & 3.00/3 \\
\code{p2--CCp3--p3} & triang. no estr. & 0.1 & 2.104 & 2.579 & 2.300 & 3.497 & 2.293 & 2.356 & 1.972 & 2.186 & 2.481 & 2.509 & 2.482 & 2.482 & 596.37 & 27166.6 & 0 & 3.00/3 \\
\hline
\code{p2--GP--p1} & hexa & 0.0 & 2.149 & 2.608 & 2.004 & 2.716 & 2.109 & 2.720 & 2.019 & 2.210 & 2.002 & 2.416 & 2.034 & 2.422 & 114.99 & 18745.8 & 0 & 3.00/3 \\
\code{p2--GP--p1} & hexa & 0.1 & 2.152 & 2.651 & 2.004 & 2.757 & 2.107 & 2.760 & 2.021 & 2.247 & 2.002 & 2.470 & 2.033 & 2.476 & 131.63 & 20347.1 & 0 & 3.00/3 \\
\code{p2--GP--p1} & triang. no estr. & 0.0 & 2.056 & 2.374 & 2.166 & 2.166 & 2.161 & 2.161 & 1.982 & 2.068 & 2.260 & 2.260 & 2.263 & 2.263 & 643.71 & 21946.4 & 0 & 3.00/3 \\
\code{p2--GP--p1} & triang. no estr. & 0.1 & 2.043 & 2.402 & 2.162 & 2.162 & 2.158 & 2.158 & 1.980 & 2.081 & 2.286 & 2.286 & 2.285 & 2.285 & 626.12 & 24414.7 & 0 & 3.00/3 \\
\hline
\code{p2--GP--p2} & hexa & 0.0 & 2.212 & 2.708 & 2.004 & 2.942 & 2.118 & 2.941 & 2.026 & 2.303 & 2.001 & 2.734 & 2.020 & 2.738 & 133.29 & 18953.3 & 0 & 3.00/3 \\
\code{p2--GP--p2} & hexa & 0.1 & 2.224 & 2.760 & 2.004 & 2.969 & 2.116 & 2.968 & 2.030 & 2.365 & 2.001 & 2.794 & 2.019 & 2.797 & 154.46 & 20458.3 & 0 & 3.00/3 \\
\code{p2--GP--p2} & triang. no estr. & 0.0 & 2.118 & 2.532 & 1.993 & 2.860 & 2.044 & 2.862 & 1.977 & 2.149 & 2.000 & 2.751 & 2.026 & 2.757 & 655.40 & 22216.4 & 0 & 3.00/3 \\
\code{p2--GP--p2} & triang. no estr. & 0.1 & 2.104 & 2.579 & 1.993 & 2.895 & 2.042 & 2.896 & 1.972 & 2.186 & 2.000 & 2.820 & 2.025 & 2.823 & 649.16 & 24557.3 & 0 & 3.00/3 \\
\hline
\code{p2--GP--p3} & hexa & 0.0 & 2.212 & 2.708 & 2.004 & 2.940 & 2.118 & 2.940 & 2.026 & 2.303 & 2.001 & 2.739 & 2.020 & 2.743 & 129.78 & 19057.3 & 0 & 3.00/3 \\
\code{p2--GP--p3} & hexa & 0.1 & 2.224 & 2.760 & 2.004 & 2.967 & 2.117 & 2.966 & 2.030 & 2.365 & 2.001 & 2.798 & 2.019 & 2.801 & 196.31 & 20537.1 & 0 & 3.00/3 \\
\code{p2--GP--p3} & triang. no estr. & 0.0 & 2.118 & 2.532 & 1.993 & 2.863 & 2.044 & 2.865 & 1.977 & 2.149 & 2.000 & 2.758 & 2.026 & 2.764 & 683.61 & 22289.2 & 0 & 3.00/3 \\
\code{p2--GP--p3} & triang. no estr. & 0.1 & 2.104 & 2.579 & 1.993 & 2.897 & 2.041 & 2.898 & 1.972 & 2.186 & 2.000 & 2.825 & 2.025 & 2.829 & 626.99 & 24625.8 & 0 & 3.00/3 \\
\hline
\code{p3--na--NO} & hexa & 0.0 & 2.121 & 2.481 & 2.256 & 2.256 & 2.258 & 2.258 & 1.999 & 2.037 & 2.472 & 2.472 & 2.469 & 2.469 & 127.96 & 21803.9 & 0 & 3.00/3 \\
\code{p3--na--NO} & hexa & 0.1 & 2.104 & 2.476 & 2.238 & 2.238 & 2.240 & 2.240 & 1.996 & 2.034 & 2.467 & 2.467 & 2.464 & 2.464 & 175.14 & 24520.5 & 0 & 3.00/3 \\
\code{p3--na--NO} & triang. no estr. & 0.0 & 2.033 & 2.201 & 2.348 & 2.348 & 2.345 & 2.345 & 1.972 & 1.988 & 3.183 & 3.183 & 3.172 & 3.172 & 725.87 & 27537.3 & 0 & 3.00/3 \\
\code{p3--na--NO} & triang. no estr. & 0.1 & 2.021 & 2.196 & 2.333 & 2.333 & 2.331 & 2.331 & 1.972 & 1.987 & 3.182 & 3.182 & 3.171 & 3.171 & 700.89 & 31526.0 & 0 & 3.00/3 \\
\hline
\code{p3--CCp1--p1} & hexa & 0.0 & 2.180 & 2.629 & 2.347 & 2.353 & 2.360 & 2.339 & 2.003 & 2.055 & 2.608 & 2.289 & 2.617 & 2.247 & 157.29 & 22965.0 & 0 & 3.00/3 \\
\code{p3--CCp1--p1} & hexa & 0.1 & 2.155 & 2.624 & 2.323 & 2.353 & 2.335 & 2.339 & 1.998 & 2.051 & 2.600 & 2.289 & 2.608 & 2.247 & 167.80 & 25249.5 & 0 & 3.00/3 \\
\code{p3--CCp1--p1} & triang. no estr. & 0.0 & 2.033 & 2.201 & 2.348 & 2.348 & 2.345 & 2.345 & 1.972 & 1.988 & 3.183 & 3.183 & 3.172 & 3.172 & 853.69 & 28339.2 & 0 & 3.00/3 \\
\code{p3--CCp1--p1} & triang. no estr. & 0.1 & 2.021 & 2.196 & 2.333 & 2.333 & 2.331 & 2.331 & 1.972 & 1.987 & 3.182 & 3.182 & 3.171 & 3.171 & 638.38 & 31988.0 & 0 & 3.00/3 \\
\hline
\code{p3--CCp1--p2} & hexa & 0.0 & 2.672 & 3.361 & 2.672 & 2.229 & 2.497 & 2.184 & 2.142 & 2.316 & 1.405 & 2.150 & 1.266 & 2.093 & 166.44 & 23100.5 & 0 & 3.00/3 \\
\code{p3--CCp1--p2} & hexa & 0.1 & 2.615 & 3.360 & 2.623 & 2.229 & 2.442 & 2.184 & 2.103 & 2.281 & 1.306 & 2.150 & 1.209 & 2.093 & 170.38 & 25367.0 & 0 & 3.00/3 \\
\code{p3--CCp1--p2} & triang. no estr. & 0.0 & 2.439 & 2.843 & 2.197 & 2.101 & 2.046 & 2.038 & 1.999 & 2.026 & 0.132 & 2.057 & 0.463 & 2.016 & 850.23 & 28492.7 & 0 & 3.00/3 \\
\code{p3--CCp1--p2} & triang. no estr. & 0.1 & 2.369 & 2.829 & 2.121 & 2.101 & 1.976 & 2.038 & 1.995 & 2.023 & 0.125 & 2.057 & 0.459 & 2.016 & 645.23 & 32105.2 & 0 & 3.00/3 \\
\hline
\code{p3--CCp1--p3} & hexa & 0.0 & 2.672 & 3.361 & 2.672 & 2.247 & 2.497 & 2.210 & 2.142 & 2.316 & 1.405 & 2.166 & 1.266 & 2.111 & 161.31 & 23166.0 & 0 & 3.00/3 \\
\code{p3--CCp1--p3} & hexa & 0.1 & 2.615 & 3.360 & 2.623 & 2.247 & 2.442 & 2.210 & 2.103 & 2.281 & 1.306 & 2.166 & 1.209 & 2.111 & 173.48 & 25411.9 & 0 & 3.00/3 \\
\code{p3--CCp1--p3} & triang. no estr. & 0.0 & 2.439 & 2.843 & 2.196 & 2.126 & 2.045 & 2.048 & 1.999 & 2.026 & 0.132 & 2.073 & 0.463 & 2.021 & 861.85 & 28573.6 & 0 & 3.00/3 \\
\code{p3--CCp1--p3} & triang. no estr. & 0.1 & 2.369 & 2.829 & 2.121 & 2.126 & 1.976 & 2.048 & 1.995 & 2.023 & 0.125 & 2.073 & 0.459 & 2.021 & 920.66 & 32166.9 & 0 & 3.00/3 \\
\hline
\code{p3--CCp2--p1} & hexa & 0.0 & 2.181 & 2.631 & 2.362 & 3.042 & 2.362 & 2.762 & 2.003 & 2.055 & 2.629 & 2.999 & 2.619 & 2.745 & 171.36 & 23027.8 & 0 & 3.00/3 \\
\code{p3--CCp2--p1} & hexa & 0.1 & 2.156 & 2.626 & 2.337 & 3.042 & 2.337 & 2.758 & 1.998 & 2.051 & 2.620 & 2.999 & 2.611 & 2.740 & 201.57 & 25251.4 & 0 & 3.00/3 \\
\code{p3--CCp2--p1} & triang. no estr. & 0.0 & 2.033 & 2.201 & 2.348 & 2.348 & 2.345 & 2.345 & 1.972 & 1.988 & 3.183 & 3.183 & 3.172 & 3.172 & 902.08 & 28536.7 & 0 & 3.00/3 \\
\code{p3--CCp2--p1} & triang. no estr. & 0.1 & 2.021 & 2.196 & 2.333 & 2.333 & 2.331 & 2.331 & 1.972 & 1.987 & 3.182 & 3.182 & 3.171 & 3.171 & 846.35 & 31989.9 & 0 & 3.00/3 \\
\hline
\code{p3--CCp2--p2} & hexa & 0.0 & 3.520 & 3.971 & 2.698 & 3.063 & 2.707 & 3.071 & 4.171 & 4.123 & 0.669 & 3.002 & 0.641 & 2.724 & 156.32 & 23188.1 & 0 & 3.00/3 \\
\code{p3--CCp2--p2} & hexa & 0.1 & 3.547 & 4.002 & 2.648 & 3.063 & 2.657 & 3.069 & 4.104 & 4.100 & 0.514 & 3.002 & 0.488 & 2.717 & 238.50 & 25382.0 & 0 & 3.00/3 \\
\code{p3--CCp2--p2} & triang. no estr. & 0.0 & 3.748 & 3.924 & 2.081 & 3.013 & 2.060 & 2.855 & 2.165 & 2.958 & 0.011 & 2.997 & 0.007 & 1.544 & 936.56 & 28757.3 & 0 & 3.00/3 \\
\code{p3--CCp2--p2} & triang. no estr. & 0.1 & 3.692 & 3.923 & 1.983 & 3.013 & 1.961 & 2.852 & 1.802 & 2.887 & 0.004 & 2.997 & 0.001 & 1.543 & 822.29 & 32135.0 & 0 & 3.00/3 \\
\hline
\code{p3--CCp2--p3} & hexa & 0.0 & 3.520 & 3.971 & 2.698 & 3.058 & 2.707 & 3.064 & 4.171 & 4.123 & 0.669 & 3.001 & 0.641 & 2.729 & 152.33 & 23274.7 & 0 & 3.00/3 \\
\code{p3--CCp2--p3} & hexa & 0.1 & 3.547 & 4.002 & 2.648 & 3.058 & 2.657 & 3.062 & 4.104 & 4.100 & 0.514 & 3.001 & 0.488 & 2.721 & 241.40 & 25437.1 & 0 & 3.00/3 \\
\code{p3--CCp2--p3} & triang. no estr. & 0.0 & 3.748 & 3.923 & 2.081 & 3.011 & 2.060 & 2.849 & 2.165 & 2.958 & 0.011 & 2.996 & 0.007 & 1.567 & 944.25 & 28825.7 & 0 & 3.00/3 \\
\code{p3--CCp2--p3} & triang. no estr. & 0.1 & 3.692 & 3.923 & 1.982 & 3.011 & 1.961 & 2.846 & 1.802 & 2.887 & 0.004 & 2.996 & 0.001 & 1.565 & 806.35 & 32197.5 & 0 & 3.00/3 \\
\hline
\code{p3--CCp3--p1} & hexa & 0.0 & 2.181 & 2.631 & 2.362 & 4.082 & 2.362 & 2.385 & 2.003 & 2.056 & 2.629 & 3.391 & 2.619 & 2.619 & 193.31 & 23618.3 & 0 & 3.00/3 \\
\code{p3--CCp3--p1} & hexa & 0.1 & 2.156 & 2.626 & 2.337 & 4.082 & 2.337 & 2.361 & 1.998 & 2.051 & 2.620 & 3.388 & 2.611 & 2.611 & 189.17 & 25269.6 & 0 & 3.00/3 \\
\code{p3--CCp3--p1} & triang. no estr. & 0.0 & 2.033 & 2.201 & 2.348 & 2.348 & 2.345 & 2.345 & 1.972 & 1.988 & 3.183 & 3.183 & 3.172 & 3.172 & 873.64 & 30657.9 & 0 & 3.00/3 \\
\code{p3--CCp3--p1} & triang. no estr. & 0.1 & 2.021 & 2.196 & 2.333 & 2.333 & 2.331 & 2.331 & 1.972 & 1.987 & 3.182 & 3.182 & 3.171 & 3.171 & 823.45 & 32119.3 & 0 & 3.00/3 \\
\hline
\code{p3--CCp3--p2} & hexa & 0.0 & 3.520 & 3.971 & 2.705 & 4.177 & 2.709 & 2.869 & 4.171 & 4.124 & 0.671 & 3.889 & 0.642 & 0.730 & 196.25 & 23779.8 & 0 & 3.00/3 \\
\code{p3--CCp3--p2} & hexa & 0.1 & 3.547 & 4.002 & 2.656 & 4.177 & 2.658 & 2.840 & 4.105 & 4.100 & 0.516 & 3.888 & 0.490 & 0.578 & 239.00 & 25414.8 & 0 & 3.00/3 \\
\code{p3--CCp3--p2} & triang. no estr. & 0.0 & 3.749 & 3.924 & 2.085 & 3.907 & 2.062 & 2.112 & 2.165 & 2.958 & 0.011 & 2.695 & 0.007 & 0.004 & 844.77 & 30890.6 & 0 & 3.00/3 \\
\code{p3--CCp3--p2} & triang. no estr. & 0.1 & 3.693 & 3.924 & 1.987 & 3.906 & 1.963 & 2.027 & 1.802 & 2.887 & 0.004 & 2.694 & 0.001 & -0.001 & 758.28 & 32277.2 & 0 & 3.00/3 \\
\hline
\code{p3--CCp3--p3} & hexa & 0.0 & 3.520 & 3.971 & 2.705 & 4.182 & 2.709 & 2.866 & 4.171 & 4.124 & 0.671 & 3.875 & 0.642 & 0.729 & 178.20 & 23856.2 & 0 & 3.00/3 \\
\code{p3--CCp3--p3} & hexa & 0.1 & 3.547 & 4.002 & 2.656 & 4.182 & 2.658 & 2.835 & 4.105 & 4.100 & 0.516 & 3.873 & 0.490 & 0.577 & 215.29 & 25464.6 & 0 & 3.00/3 \\
\code{p3--CCp3--p3} & triang. no estr. & 0.0 & 3.749 & 3.924 & 2.085 & 3.889 & 2.061 & 2.108 & 2.165 & 2.958 & 0.011 & 2.592 & 0.007 & 0.004 & 905.41 & 30993.9 & 0 & 3.00/3 \\
\code{p3--CCp3--p3} & triang. no estr. & 0.1 & 3.693 & 3.924 & 1.986 & 3.889 & 1.963 & 2.022 & 1.802 & 2.887 & 0.004 & 2.591 & 0.001 & -0.001 & 718.57 & 32342.4 & 0 & 3.00/3 \\
\hline
\code{p3--GP--p1} & hexa & 0.0 & 2.181 & 2.631 & 2.002 & 2.762 & 1.992 & 2.757 & 2.003 & 2.056 & 2.002 & 2.661 & 2.008 & 2.650 & 151.23 & 21847.0 & 0 & 3.00/3 \\
\code{p3--GP--p1} & hexa & 0.1 & 2.156 & 2.626 & 2.002 & 2.756 & 1.991 & 2.752 & 1.998 & 2.051 & 2.002 & 2.653 & 2.008 & 2.642 & 196.90 & 24560.6 & 0 & 3.00/3 \\
\code{p3--GP--p1} & triang. no estr. & 0.0 & 2.033 & 2.201 & 2.348 & 2.348 & 2.345 & 2.345 & 1.972 & 1.988 & 3.183 & 3.183 & 3.172 & 3.172 & 792.64 & 27566.0 & 0 & 3.00/3 \\
\code{p3--GP--p1} & triang. no estr. & 0.1 & 2.021 & 2.196 & 2.333 & 2.333 & 2.331 & 2.331 & 1.972 & 1.987 & 3.182 & 3.182 & 3.171 & 3.171 & 711.74 & 31533.8 & 0 & 3.00/3 \\
\hline
\code{p3--GP--p2} & hexa & 0.0 & 3.520 & 3.971 & 2.000 & 3.744 & 1.999 & 3.748 & 4.171 & 4.124 & 2.001 & 1.458 & 2.006 & 1.396 & 178.12 & 21950.5 & 0 & 3.00/3 \\
\code{p3--GP--p2} & hexa & 0.1 & 3.547 & 4.002 & 1.999 & 3.757 & 1.998 & 3.761 & 4.105 & 4.100 & 2.001 & 1.348 & 2.006 & 1.284 & 181.72 & 24663.2 & 0 & 3.00/3 \\
\code{p3--GP--p2} & triang. no estr. & 0.0 & 3.749 & 3.924 & 1.992 & 2.844 & 2.005 & 2.843 & 2.165 & 2.958 & 2.000 & -0.092 & 2.012 & -0.097 & 974.21 & 27701.8 & 0 & 3.00/3 \\
\code{p3--GP--p2} & triang. no estr. & 0.1 & 3.693 & 3.924 & 1.992 & 2.836 & 2.005 & 2.836 & 1.802 & 2.887 & 2.000 & -0.098 & 2.012 & -0.103 & 907.50 & 31675.3 & 0 & 3.00/3 \\
\hline
\code{p3--GP--p3} & hexa & 0.0 & 3.520 & 3.971 & 2.000 & 3.738 & 1.999 & 3.742 & 4.171 & 4.124 & 2.001 & 1.420 & 2.006 & 1.355 & 206.65 & 22042.1 & 0 & 3.00/3 \\
\code{p3--GP--p3} & hexa & 0.1 & 3.547 & 4.002 & 2.000 & 3.751 & 1.999 & 3.755 & 4.105 & 4.100 & 2.001 & 1.305 & 2.006 & 1.239 & 227.33 & 24736.5 & 0 & 3.00/3 \\
\code{p3--GP--p3} & triang. no estr. & 0.0 & 3.748 & 3.924 & 1.992 & 2.812 & 2.005 & 2.811 & 2.165 & 2.958 & 2.000 & -0.108 & 2.012 & -0.114 & 802.37 & 27772.6 & 0 & 3.00/3 \\
\code{p3--GP--p3} & triang. no estr. & 0.1 & 3.693 & 3.923 & 1.992 & 2.803 & 2.005 & 2.803 & 1.802 & 2.887 & 2.000 & -0.115 & 2.012 & -0.120 & 733.20 & 31735.5 & 0 & 3.00/3 \\
\bottomrule
\end{longtable}
\normalsize

\paragraph{$\Delta t=10^{-3}$}
\tiny
\setlength{\tabcolsep}{1.4pt}
\begin{longtable}{lllrrrrrrrrrrrrrrrr}
\caption{Resultados manufacturados 2D para \code{diagonalIion}, \code{crankNicolson}, masa \code{consistent}, estados \code{RKF45} y $\Delta t=10^{-3}$. El primer bloque de convergencia se ajusta sobre $N=10$--$100$; el segundo sobre $N=70$--$100$. La primera columna registra la terna \code{Vm-states-Iion} efectivamente usada en la corrida. \textbf{RB} es la cantidad de pasos de tiempo que agotaron el presupuesto de iteraciones no lineales, acumulada sobre los 91 tamaños de malla de la fila (4550 pasos de tiempo en total por fila). Con \code{nonlinearAcceptUnconverged=false} un paso no convergido se revierte a los valores del inicio del paso mientras el reloj avanza igual. \textbf{$n_{\mathrm{nl}}$} reporta la mediana y el máximo de iteraciones no lineales por paso sobre esos mismos pasos; el presupuesto configurado es \code{nonlinearIterations}$=2000$.}\label{tab:es-results-2d-diagonaliion-1p0em03}\\
\toprule
& & & \multicolumn{6}{c}{Ajuste $N=10$--$100$} & \multicolumn{6}{c}{Ajuste $N=70$--$100$} & & & & \\
\cmidrule(lr){4-9}\cmidrule(lr){10-15}
\code{Vm-states-Iion} & Tipo de malla & $\alpha$ & $V_m$ & $V_m^G$ & $u_1$ & $u_1^G$ & $u_2$ & $u_2^G$ & $V_m$ & $V_m^G$ & $u_1$ & $u_1^G$ & $u_2$ & $u_2^G$ & $t_{\Sigma}$ [s] & RSS$_{\Sigma}$ [MB] & RB & $n_{\mathrm{nl}}$ \\
\midrule
\endfirsthead
\toprule
& & & \multicolumn{6}{c}{Ajuste $N=10$--$100$} & \multicolumn{6}{c}{Ajuste $N=70$--$100$} & & & & \\
\cmidrule(lr){4-9}\cmidrule(lr){10-15}
\code{Vm-states-Iion} & Tipo de malla & $\alpha$ & $V_m$ & $V_m^G$ & $u_1$ & $u_1^G$ & $u_2$ & $u_2^G$ & $V_m$ & $V_m^G$ & $u_1$ & $u_1^G$ & $u_2$ & $u_2^G$ & $t_{\Sigma}$ [s] & RSS$_{\Sigma}$ [MB] & RB & $n_{\mathrm{nl}}$ \\
\midrule
\endhead
\code{NO--na--NO} & hexa & 0.0 & 1.997 & 1.997 & 1.998 & 1.998 & 1.998 & 1.998 & 2.000 & 2.000 & 2.004 & 2.004 & 2.004 & 2.004 & 66.387 & 16461.7 & 0 & 2.00/2 \\
\code{NO--na--NO} & hexa & 0.1 & 1.851 & 1.851 & 1.607 & 1.607 & 1.613 & 1.613 & 2.294 & 2.294 & 2.140 & 2.140 & 2.142 & 2.142 & 121.55 & 18342.1 & 0 & 2.00/2 \\
\code{NO--na--NO} & triang. no estr. & 0.0 & 0.798 & 0.798 & 0.777 & 0.777 & 0.783 & 0.783 & 0.963 & 0.963 & 1.038 & 1.038 & 1.107 & 1.107 & 375.50 & 17145.7 & 0 & 2.00/2 \\
\code{NO--na--NO} & triang. no estr. & 0.1 & 0.785 & 0.785 & 0.771 & 0.771 & 0.776 & 0.776 & 0.953 & 0.953 & 1.027 & 1.027 & 1.097 & 1.097 & 537.35 & 20800.8 & 0 & 2.00/2 \\
\hline
\code{NO--CCp1--p1} & hexa & 0.0 & 1.997 & 1.997 & 1.999 & 2.354 & 1.998 & 2.333 & 2.000 & 2.000 & 2.004 & 2.290 & 2.004 & 2.242 & 104.98 & 16905.3 & 0 & 2.00/2 \\
\code{NO--CCp1--p1} & hexa & 0.1 & 1.851 & 1.851 & 1.608 & 2.354 & 1.613 & 2.310 & 2.294 & 2.294 & 2.140 & 2.290 & 2.142 & 2.226 & 158.70 & 18479.0 & 0 & 2.00/2 \\
\code{NO--CCp1--p1} & triang. no estr. & 0.0 & 0.798 & 0.798 & 0.777 & 0.777 & 0.783 & 0.783 & 0.963 & 0.963 & 1.038 & 1.038 & 1.107 & 1.107 & 486.71 & 17867.2 & 0 & 2.00/2 \\
\code{NO--CCp1--p1} & triang. no estr. & 0.1 & 0.785 & 0.785 & 0.771 & 0.771 & 0.776 & 0.776 & 0.953 & 0.953 & 1.027 & 1.027 & 1.097 & 1.097 & 595.77 & 20904.5 & 0 & 2.00/2 \\
\hline
\code{NO--CCp1--p2} & hexa & 0.0 & 1.997 & 1.997 & 2.000 & 2.229 & 1.999 & 2.184 & 2.000 & 2.000 & 2.004 & 2.150 & 2.004 & 2.094 & 115.96 & 17132.7 & 0 & 2.00/2 \\
\code{NO--CCp1--p2} & hexa & 0.1 & 1.851 & 1.851 & 1.608 & 2.229 & 1.614 & 2.181 & 2.295 & 2.295 & 2.140 & 2.150 & 2.142 & 2.092 & 177.55 & 18601.7 & 0 & 2.00/2 \\
\code{NO--CCp1--p2} & triang. no estr. & 0.0 & 0.798 & 0.798 & 0.777 & 2.111 & 0.783 & 1.163 & 0.963 & 0.963 & 1.038 & 2.102 & 1.107 & 1.150 & 554.31 & 18116.1 & 0 & 2.00/2 \\
\code{NO--CCp1--p2} & triang. no estr. & 0.1 & 0.785 & 0.785 & 0.771 & 2.112 & 0.776 & 1.171 & 0.953 & 0.953 & 1.027 & 2.105 & 1.097 & 1.143 & 726.29 & 21044.1 & 0 & 2.00/2 \\
\hline
\code{NO--CCp1--p3} & hexa & 0.0 & 1.997 & 1.997 & 2.000 & 2.247 & 1.999 & 2.211 & 2.000 & 2.000 & 2.004 & 2.167 & 2.004 & 2.113 & 115.51 & 17272.0 & 0 & 2.00/2 \\
\code{NO--CCp1--p3} & hexa & 0.1 & 1.851 & 1.851 & 1.608 & 2.247 & 1.614 & 2.206 & 2.295 & 2.295 & 2.140 & 2.167 & 2.142 & 2.110 & 210.02 & 18665.7 & 0 & 2.00/2 \\
\code{NO--CCp1--p3} & triang. no estr. & 0.0 & 0.798 & 0.798 & 0.777 & 2.139 & 0.783 & 1.131 & 0.963 & 0.963 & 1.038 & 2.131 & 1.107 & 1.144 & 522.54 & 18249.6 & 0 & 2.00/2 \\
\code{NO--CCp1--p3} & triang. no estr. & 0.1 & 0.785 & 0.785 & 0.771 & 2.140 & 0.776 & 1.138 & 0.953 & 0.953 & 1.027 & 2.135 & 1.097 & 1.137 & 605.11 & 21107.3 & 0 & 2.00/2 \\
\hline
\code{NO--CCp2--p1} & hexa & 0.0 & 1.997 & 1.997 & 2.000 & 3.042 & 1.997 & 2.510 & 1.999 & 1.999 & 2.004 & 2.998 & 2.003 & 2.157 & 163.49 & 17393.5 & 0 & 2.00/2 \\
\code{NO--CCp2--p1} & hexa & 0.1 & 1.851 & 1.851 & 1.603 & 3.038 & 1.613 & 1.942 & 2.294 & 2.294 & 2.140 & 2.987 & 2.142 & 2.167 & 258.60 & 18604.2 & 0 & 2.00/2 \\
\code{NO--CCp2--p1} & triang. no estr. & 0.0 & 0.798 & 0.798 & 0.777 & 0.777 & 0.783 & 0.783 & 0.963 & 0.963 & 1.038 & 1.038 & 1.107 & 1.107 & 633.81 & 18972.1 & 0 & 2.00/2 \\
\code{NO--CCp2--p1} & triang. no estr. & 0.1 & 0.785 & 0.785 & 0.771 & 0.771 & 0.776 & 0.776 & 0.953 & 0.953 & 1.027 & 1.027 & 1.097 & 1.097 & 848.91 & 21040.1 & 0 & 2.00/2 \\
\hline
\code{NO--CCp2--p2} & hexa & 0.0 & 1.995 & 1.995 & 2.003 & 3.063 & 1.997 & 2.671 & 1.999 & 1.999 & 2.004 & 3.002 & 2.003 & 2.246 & 198.35 & 17648.5 & 0 & 2.00/2 \\
\code{NO--CCp2--p2} & hexa & 0.1 & 1.853 & 1.853 & 1.597 & 3.061 & 1.619 & 2.063 & 2.289 & 2.289 & 2.140 & 2.996 & 2.141 & 2.192 & 269.27 & 18821.5 & 0 & 2.00/2 \\
\code{NO--CCp2--p2} & triang. no estr. & 0.0 & 0.799 & 0.799 & 0.777 & 1.940 & 0.783 & 0.775 & 0.963 & 0.963 & 1.038 & 1.095 & 1.107 & 1.097 & 779.96 & 19260.5 & 0 & 2.00/2 \\
\code{NO--CCp2--p2} & triang. no estr. & 0.1 & 0.786 & 0.786 & 0.771 & 1.955 & 0.776 & 0.772 & 0.953 & 0.953 & 1.027 & 1.090 & 1.097 & 1.087 & 754.35 & 21228.6 & 0 & 2.00/2 \\
\hline
\code{NO--CCp2--p3} & hexa & 0.0 & 1.995 & 1.995 & 2.003 & 3.058 & 1.997 & 2.670 & 1.999 & 1.999 & 2.004 & 3.001 & 2.003 & 2.251 & 219.51 & 17774.7 & 0 & 2.00/2 \\
\code{NO--CCp2--p3} & hexa & 0.1 & 1.853 & 1.853 & 1.597 & 3.056 & 1.619 & 2.064 & 2.289 & 2.289 & 2.140 & 2.995 & 2.141 & 2.193 & 258.90 & 18968.1 & 0 & 2.00/2 \\
\code{NO--CCp2--p3} & triang. no estr. & 0.0 & 0.799 & 0.799 & 0.777 & 1.950 & 0.783 & 0.775 & 0.963 & 0.963 & 1.038 & 1.098 & 1.107 & 1.097 & 634.20 & 19406.0 & 0 & 2.00/2 \\
\code{NO--CCp2--p3} & triang. no estr. & 0.1 & 0.786 & 0.786 & 0.771 & 1.966 & 0.776 & 0.772 & 0.953 & 0.953 & 1.027 & 1.092 & 1.097 & 1.087 & 913.75 & 21372.5 & 0 & 2.00/2 \\
\hline
\code{NO--CCp3--p1} & hexa & 0.0 & 1.997 & 1.997 & 1.998 & 3.944 & 1.997 & 2.030 & 1.999 & 1.999 & 2.004 & 2.505 & 2.004 & 2.004 & 334.52 & 18754.3 & 0 & 2.00/2 \\
\code{NO--CCp3--p1} & hexa & 0.1 & 1.851 & 1.851 & 1.600 & 3.114 & 1.613 & 1.657 & 2.294 & 2.294 & 2.140 & 2.324 & 2.142 & 2.170 & 378.04 & 19859.4 & 0 & 2.00/2 \\
\code{NO--CCp3--p1} & triang. no estr. & 0.0 & 0.798 & 0.798 & 0.777 & 0.777 & 0.783 & 0.783 & 0.963 & 0.963 & 1.038 & 1.038 & 1.107 & 1.107 & 1067.1 & 21482.2 & 0 & 2.00/2 \\
\code{NO--CCp3--p1} & triang. no estr. & 0.1 & 0.785 & 0.785 & 0.771 & 0.771 & 0.776 & 0.776 & 0.953 & 0.953 & 1.027 & 1.027 & 1.097 & 1.097 & 1049.8 & 23311.5 & 0 & 2.00/2 \\
\hline
\code{NO--CCp3--p2} & hexa & 0.0 & 1.996 & 1.996 & 1.998 & 4.244 & 1.996 & 2.092 & 1.999 & 1.999 & 2.004 & 3.862 & 2.003 & 2.002 & 342.67 & 19008.4 & 0 & 2.00/2 \\
\code{NO--CCp3--p2} & hexa & 0.1 & 1.853 & 1.853 & 1.590 & 3.420 & 1.617 & 1.718 & 2.289 & 2.289 & 2.140 & 2.581 & 2.141 & 2.203 & 351.87 & 20077.6 & 0 & 2.00/2 \\
\code{NO--CCp3--p2} & triang. no estr. & 0.0 & 0.799 & 0.799 & 0.777 & 1.070 & 0.783 & 0.795 & 0.963 & 0.963 & 1.038 & 1.038 & 1.107 & 1.110 & 1034.3 & 21759.2 & 0 & 2.00/2 \\
\code{NO--CCp3--p2} & triang. no estr. & 0.1 & 0.786 & 0.786 & 0.771 & 1.080 & 0.776 & 0.797 & 0.953 & 0.953 & 1.027 & 1.028 & 1.097 & 1.100 & 1164.8 & 23565.8 & 0 & 2.00/2 \\
\hline
\code{NO--CCp3--p3} & hexa & 0.0 & 1.996 & 1.996 & 1.998 & 4.257 & 1.996 & 2.090 & 1.999 & 1.999 & 2.004 & 3.849 & 2.003 & 2.002 & 383.46 & 19125.2 & 0 & 2.00/2 \\
\code{NO--CCp3--p3} & hexa & 0.1 & 1.853 & 1.853 & 1.590 & 3.408 & 1.617 & 1.717 & 2.289 & 2.289 & 2.140 & 2.570 & 2.141 & 2.203 & 369.57 & 20194.4 & 0 & 2.00/2 \\
\code{NO--CCp3--p3} & triang. no estr. & 0.0 & 0.799 & 0.799 & 0.777 & 1.053 & 0.783 & 0.795 & 0.963 & 0.963 & 1.038 & 1.038 & 1.107 & 1.110 & 824.24 & 21903.0 & 0 & 2.00/2 \\
\code{NO--CCp3--p3} & triang. no estr. & 0.1 & 0.786 & 0.786 & 0.771 & 1.062 & 0.776 & 0.796 & 0.953 & 0.953 & 1.027 & 1.028 & 1.097 & 1.100 & 1012.6 & 23686.4 & 0 & 2.00/2 \\
\hline
\code{NO--GP--p1} & hexa & 0.0 & 1.997 & 1.997 & 1.999 & 2.050 & 1.998 & 2.053 & 2.000 & 2.000 & 2.000 & 2.007 & 2.000 & 2.007 & 83.484 & 16618.9 & 0 & 2.00/2 \\
\code{NO--GP--p1} & hexa & 0.1 & 1.851 & 1.851 & 1.999 & 2.042 & 1.990 & 2.046 & 2.294 & 2.294 & 2.000 & 2.015 & 2.008 & 2.016 & 169.69 & 18446.5 & 0 & 2.00/2 \\
\code{NO--GP--p1} & triang. no estr. & 0.0 & 0.798 & 0.798 & 0.777 & 0.777 & 0.783 & 0.783 & 0.963 & 0.963 & 1.038 & 1.038 & 1.107 & 1.107 & 499.64 & 17283.5 & 0 & 2.00/2 \\
\code{NO--GP--p1} & triang. no estr. & 0.1 & 0.785 & 0.785 & 0.771 & 0.771 & 0.776 & 0.776 & 0.953 & 0.953 & 1.027 & 1.027 & 1.097 & 1.097 & 657.51 & 20901.9 & 0 & 2.00/2 \\
\hline
\code{NO--GP--p2} & hexa & 0.0 & 1.997 & 1.997 & 1.999 & 2.013 & 1.998 & 2.015 & 1.999 & 1.999 & 2.000 & 2.002 & 2.000 & 2.002 & 202.57 & 16847.9 & 0 & 2.00/2 \\
\code{NO--GP--p2} & hexa & 0.1 & 1.853 & 1.853 & 1.999 & 2.013 & 1.998 & 2.014 & 2.290 & 2.290 & 2.000 & 2.003 & 2.001 & 2.003 & 267.50 & 18578.8 & 0 & 2.00/2 \\
\code{NO--GP--p2} & triang. no estr. & 0.0 & 0.799 & 0.799 & 1.878 & 1.417 & 1.101 & 1.410 & 0.963 & 0.963 & 1.677 & 1.218 & 1.112 & 1.269 & 347.95 & 17526.6 & 0 & 2.00/2 \\
\code{NO--GP--p2} & triang. no estr. & 0.1 & 0.786 & 0.786 & 1.880 & 1.426 & 1.105 & 1.419 & 0.953 & 0.953 & 1.680 & 1.219 & 1.104 & 1.270 & 548.75 & 21024.5 & 0 & 2.00/2 \\
\hline
\code{NO--GP--p3} & hexa & 0.0 & 1.997 & 1.997 & 1.999 & 2.016 & 1.999 & 2.018 & 1.999 & 1.999 & 2.000 & 2.002 & 2.000 & 2.002 & 233.61 & 16957.7 & 0 & 2.00/2 \\
\code{NO--GP--p3} & hexa & 0.1 & 1.853 & 1.853 & 1.999 & 2.016 & 1.998 & 2.018 & 2.290 & 2.290 & 2.000 & 2.004 & 2.002 & 2.004 & 260.87 & 18624.0 & 0 & 2.00/2 \\
\code{NO--GP--p3} & triang. no estr. & 0.0 & 0.799 & 0.799 & 1.878 & 1.388 & 1.101 & 1.383 & 0.963 & 0.963 & 1.677 & 1.205 & 1.112 & 1.258 & 405.61 & 17643.7 & 0 & 2.00/2 \\
\code{NO--GP--p3} & triang. no estr. & 0.1 & 0.786 & 0.786 & 1.880 & 1.397 & 1.106 & 1.391 & 0.953 & 0.953 & 1.680 & 1.206 & 1.104 & 1.258 & 700.56 & 21097.0 & 0 & 2.00/2 \\
\hline
\code{p1--na--NO} & hexa & 0.0 & 1.912 & 2.147 & 1.616 & 1.616 & 1.626 & 1.626 & 2.611 & 2.168 & 2.380 & 2.380 & 2.385 & 2.385 & 161.28 & 17839.1 & 0 & 2.00/2 \\
\code{p1--na--NO} & hexa & 0.1 & 1.971 & 2.163 & 1.677 & 1.677 & 1.687 & 1.687 & 2.615 & 2.141 & 2.414 & 2.414 & 2.419 & 2.419 & 189.26 & 19067.4 & 0 & 2.00/2 \\
\code{p1--na--NO} & triang. no estr. & 0.0 & 2.292 & 2.131 & 2.004 & 2.004 & 2.003 & 2.003 & 2.600 & 2.164 & 2.545 & 2.545 & 2.542 & 2.542 & 773.26 & 19883.4 & 0 & 2.00/2 \\
\code{p1--na--NO} & triang. no estr. & 0.1 & 2.345 & 2.112 & 2.092 & 2.092 & 2.089 & 2.089 & 2.552 & 2.095 & 2.516 & 2.516 & 2.515 & 2.515 & 976.30 & 21554.6 & 0 & 2.00/2 \\
\hline
\code{p1--CCp1--p1} & hexa & 0.0 & 1.912 & 2.147 & 1.616 & 2.350 & 1.626 & 1.955 & 2.611 & 2.168 & 2.380 & 2.295 & 2.385 & 2.369 & 231.15 & 18280.8 & 0 & 2.00/2 \\
\code{p1--CCp1--p1} & hexa & 0.1 & 1.971 & 2.163 & 1.677 & 2.351 & 1.687 & 2.045 & 2.615 & 2.141 & 2.414 & 2.294 & 2.419 & 2.382 & 248.72 & 19264.1 & 0 & 2.00/2 \\
\code{p1--CCp1--p1} & triang. no estr. & 0.0 & 2.292 & 2.131 & 2.004 & 2.004 & 2.003 & 2.003 & 2.600 & 2.164 & 2.545 & 2.545 & 2.542 & 2.542 & 973.13 & 20403.7 & 0 & 2.00/2 \\
\code{p1--CCp1--p1} & triang. no estr. & 0.1 & 2.345 & 2.112 & 2.092 & 2.092 & 2.089 & 2.089 & 2.552 & 2.095 & 2.516 & 2.516 & 2.515 & 2.515 & 1046.9 & 21754.1 & 0 & 2.00/2 \\
\hline
\code{p1--CCp1--p2} & hexa & 0.0 & 1.912 & 2.147 & 1.616 & 2.228 & 1.626 & 2.057 & 2.611 & 2.168 & 2.380 & 2.154 & 2.385 & 2.233 & 236.12 & 18508.2 & 0 & 2.00/2 \\
\code{p1--CCp1--p2} & hexa & 0.1 & 1.971 & 2.163 & 1.676 & 2.228 & 1.686 & 2.104 & 2.615 & 2.141 & 2.414 & 2.153 & 2.419 & 2.209 & 259.07 & 19395.4 & 0 & 2.00/2 \\
\code{p1--CCp1--p2} & triang. no estr. & 0.0 & 2.292 & 2.131 & 2.004 & 2.103 & 2.003 & 2.036 & 2.600 & 2.164 & 2.545 & 2.063 & 2.542 & 2.283 & 1044.4 & 20638.8 & 0 & 2.00/2 \\
\code{p1--CCp1--p2} & triang. no estr. & 0.1 & 2.345 & 2.112 & 2.092 & 2.103 & 2.089 & 2.066 & 2.552 & 2.095 & 2.516 & 2.061 & 2.515 & 2.188 & 1052.1 & 21916.4 & 0 & 2.00/2 \\
\hline
\code{p1--CCp1--p3} & hexa & 0.0 & 1.912 & 2.147 & 1.616 & 2.246 & 1.626 & 2.063 & 2.611 & 2.168 & 2.380 & 2.171 & 2.385 & 2.256 & 236.80 & 18612.1 & 0 & 2.00/2 \\
\code{p1--CCp1--p3} & hexa & 0.1 & 1.971 & 2.163 & 1.676 & 2.246 & 1.686 & 2.115 & 2.615 & 2.141 & 2.414 & 2.170 & 2.419 & 2.234 & 258.99 & 19447.7 & 0 & 2.00/2 \\
\code{p1--CCp1--p3} & triang. no estr. & 0.0 & 2.292 & 2.131 & 2.004 & 2.127 & 2.003 & 2.040 & 2.600 & 2.164 & 2.545 & 2.081 & 2.542 & 2.312 & 1049.2 & 20766.1 & 0 & 2.00/2 \\
\code{p1--CCp1--p3} & triang. no estr. & 0.1 & 2.345 & 2.112 & 2.092 & 2.128 & 2.089 & 2.076 & 2.552 & 2.095 & 2.516 & 2.078 & 2.515 & 2.216 & 1107.6 & 21983.2 & 0 & 2.00/2 \\
\hline
\code{p1--CCp2--p1} & hexa & 0.0 & 1.912 & 2.147 & 1.616 & 2.903 & 1.626 & 1.665 & 2.611 & 2.168 & 2.380 & 2.782 & 2.385 & 2.394 & 263.19 & 18761.8 & 0 & 2.00/2 \\
\code{p1--CCp2--p1} & hexa & 0.1 & 1.971 & 2.163 & 1.676 & 2.953 & 1.687 & 1.736 & 2.615 & 2.141 & 2.414 & 2.858 & 2.419 & 2.428 & 293.62 & 19371.8 & 0 & 2.00/2 \\
\code{p1--CCp2--p1} & triang. no estr. & 0.0 & 2.292 & 2.131 & 2.004 & 2.004 & 2.003 & 2.003 & 2.600 & 2.164 & 2.545 & 2.545 & 2.542 & 2.542 & 1109.4 & 21489.3 & 0 & 2.00/2 \\
\code{p1--CCp2--p1} & triang. no estr. & 0.1 & 2.345 & 2.112 & 2.092 & 2.092 & 2.089 & 2.089 & 2.552 & 2.095 & 2.516 & 2.516 & 2.515 & 2.515 & 1214.9 & 22111.9 & 0 & 2.00/2 \\
\hline
\code{p1--CCp2--p2} & hexa & 0.0 & 1.912 & 2.147 & 1.616 & 2.990 & 1.627 & 1.689 & 2.611 & 2.168 & 2.380 & 2.882 & 2.385 & 2.400 & 271.94 & 18974.1 & 0 & 2.00/2 \\
\code{p1--CCp2--p2} & hexa & 0.1 & 1.972 & 2.163 & 1.676 & 3.018 & 1.687 & 1.765 & 2.615 & 2.141 & 2.414 & 2.931 & 2.419 & 2.435 & 314.85 & 19557.6 & 0 & 2.00/2 \\
\code{p1--CCp2--p2} & triang. no estr. & 0.0 & 2.292 & 2.131 & 2.004 & 2.676 & 2.003 & 1.978 & 2.600 & 2.164 & 2.546 & 2.665 & 2.542 & 2.522 & 1119.2 & 21722.9 & 0 & 2.00/2 \\
\code{p1--CCp2--p2} & triang. no estr. & 0.1 & 2.345 & 2.112 & 2.092 & 2.799 & 2.089 & 2.063 & 2.552 & 2.095 & 2.516 & 2.731 & 2.515 & 2.504 & 1251.5 & 22294.3 & 0 & 2.00/2 \\
\hline
\code{p1--CCp2--p3} & hexa & 0.0 & 1.912 & 2.147 & 1.616 & 2.987 & 1.627 & 1.688 & 2.611 & 2.168 & 2.380 & 2.884 & 2.385 & 2.400 & 289.83 & 19073.9 & 0 & 2.00/2 \\
\code{p1--CCp2--p3} & hexa & 0.1 & 1.972 & 2.163 & 1.676 & 3.015 & 1.687 & 1.765 & 2.615 & 2.141 & 2.414 & 2.931 & 2.419 & 2.435 & 328.75 & 19618.5 & 0 & 2.00/2 \\
\code{p1--CCp2--p3} & triang. no estr. & 0.0 & 2.292 & 2.131 & 2.004 & 2.681 & 2.003 & 1.978 & 2.600 & 2.164 & 2.546 & 2.668 & 2.542 & 2.522 & 1128.9 & 21852.8 & 0 & 2.00/2 \\
\code{p1--CCp2--p3} & triang. no estr. & 0.1 & 2.345 & 2.112 & 2.092 & 2.802 & 2.089 & 2.062 & 2.552 & 2.095 & 2.516 & 2.736 & 2.515 & 2.504 & 1283.7 & 22413.7 & 0 & 2.00/2 \\
\hline
\code{p1--CCp3--p1} & hexa & 0.0 & 1.912 & 2.147 & 1.617 & 2.034 & 1.626 & 1.630 & 2.611 & 2.168 & 2.380 & 2.378 & 2.385 & 2.406 & 395.23 & 20040.1 & 0 & 2.00/2 \\
\code{p1--CCp3--p1} & hexa & 0.1 & 1.971 & 2.163 & 1.677 & 2.164 & 1.687 & 1.692 & 2.615 & 2.141 & 2.414 & 2.408 & 2.419 & 2.441 & 436.47 & 20506.7 & 0 & 2.00/2 \\
\code{p1--CCp3--p1} & triang. no estr. & 0.0 & 2.292 & 2.131 & 2.004 & 2.004 & 2.003 & 2.003 & 2.600 & 2.164 & 2.545 & 2.545 & 2.542 & 2.542 & 1416.4 & 23811.3 & 0 & 2.00/2 \\
\code{p1--CCp3--p1} & triang. no estr. & 0.1 & 2.345 & 2.112 & 2.092 & 2.092 & 2.089 & 2.089 & 2.552 & 2.095 & 2.516 & 2.516 & 2.515 & 2.515 & 1506.6 & 24280.8 & 0 & 2.00/2 \\
\hline
\code{p1--CCp3--p2} & hexa & 0.0 & 1.913 & 2.147 & 1.617 & 2.239 & 1.627 & 1.635 & 2.611 & 2.168 & 2.380 & 2.377 & 2.385 & 2.426 & 409.23 & 20223.1 & 0 & 2.00/2 \\
\code{p1--CCp3--p2} & hexa & 0.1 & 1.972 & 2.163 & 1.677 & 2.394 & 1.688 & 1.698 & 2.615 & 2.141 & 2.414 & 2.402 & 2.419 & 2.463 & 432.50 & 20725.2 & 0 & 2.00/2 \\
\code{p1--CCp3--p2} & triang. no estr. & 0.0 & 2.292 & 2.131 & 2.004 & 2.031 & 2.003 & 1.965 & 2.600 & 2.164 & 2.546 & 2.521 & 2.542 & 2.523 & 1390.3 & 24059.1 & 0 & 2.00/2 \\
\code{p1--CCp3--p2} & triang. no estr. & 0.1 & 2.345 & 2.112 & 2.092 & 2.141 & 2.089 & 2.052 & 2.552 & 2.095 & 2.516 & 2.506 & 2.515 & 2.510 & 1558.2 & 24526.5 & 0 & 2.00/2 \\
\hline
\code{p1--CCp3--p3} & hexa & 0.0 & 1.913 & 2.147 & 1.617 & 2.224 & 1.627 & 1.634 & 2.611 & 2.168 & 2.380 & 2.376 & 2.385 & 2.426 & 418.15 & 20356.2 & 0 & 2.00/2 \\
\code{p1--CCp3--p3} & hexa & 0.1 & 1.972 & 2.163 & 1.677 & 2.379 & 1.688 & 1.698 & 2.615 & 2.141 & 2.414 & 2.401 & 2.419 & 2.462 & 447.84 & 20842.8 & 0 & 2.00/2 \\
\code{p1--CCp3--p3} & triang. no estr. & 0.0 & 2.292 & 2.131 & 2.004 & 2.028 & 2.003 & 1.965 & 2.600 & 2.164 & 2.546 & 2.522 & 2.542 & 2.523 & 1368.8 & 24169.1 & 0 & 2.00/2 \\
\code{p1--CCp3--p3} & triang. no estr. & 0.1 & 2.345 & 2.112 & 2.092 & 2.136 & 2.089 & 2.052 & 2.552 & 2.095 & 2.516 & 2.506 & 2.515 & 2.510 & 1450.6 & 24662.7 & 0 & 2.00/2 \\
\hline
\code{p1--GP--p1} & hexa & 0.0 & 1.912 & 2.147 & 1.997 & 2.152 & 1.836 & 2.158 & 2.611 & 2.168 & 2.005 & 2.309 & 2.249 & 2.308 & 204.40 & 17929.4 & 0 & 2.00/2 \\
\code{p1--GP--p1} & hexa & 0.1 & 1.971 & 2.163 & 1.999 & 2.219 & 1.895 & 2.225 & 2.615 & 2.141 & 2.003 & 2.309 & 2.222 & 2.308 & 241.70 & 19121.1 & 0 & 2.00/2 \\
\code{p1--GP--p1} & triang. no estr. & 0.0 & 2.292 & 2.131 & 2.004 & 2.004 & 2.003 & 2.003 & 2.600 & 2.164 & 2.545 & 2.545 & 2.542 & 2.542 & 842.03 & 19912.6 & 0 & 2.00/2 \\
\code{p1--GP--p1} & triang. no estr. & 0.1 & 2.345 & 2.112 & 2.092 & 2.092 & 2.089 & 2.089 & 2.552 & 2.095 & 2.516 & 2.516 & 2.515 & 2.515 & 1027.9 & 21576.0 & 0 & 2.00/2 \\
\hline
\code{p1--GP--p2} & hexa & 0.0 & 1.913 & 2.147 & 1.998 & 2.181 & 1.949 & 2.179 & 2.611 & 2.168 & 2.001 & 2.155 & 2.096 & 2.151 & 320.49 & 18117.7 & 0 & 2.00/2 \\
\code{p1--GP--p2} & hexa & 0.1 & 1.972 & 2.163 & 1.999 & 2.199 & 1.972 & 2.197 & 2.615 & 2.141 & 2.001 & 2.150 & 2.071 & 2.145 & 327.66 & 19246.6 & 0 & 2.00/2 \\
\code{p1--GP--p2} & triang. no estr. & 0.0 & 2.292 & 2.131 & 1.991 & 2.067 & 2.009 & 2.066 & 2.600 & 2.164 & 2.002 & 2.176 & 2.268 & 2.168 & 931.04 & 20140.0 & 0 & 2.00/2 \\
\code{p1--GP--p2} & triang. no estr. & 0.1 & 2.345 & 2.112 & 1.992 & 2.087 & 2.034 & 2.085 & 2.552 & 2.095 & 2.000 & 2.109 & 2.167 & 2.105 & 1145.7 & 21718.7 & 0 & 2.00/2 \\
\hline
\code{p1--GP--p3} & hexa & 0.0 & 1.913 & 2.147 & 1.999 & 2.194 & 1.949 & 2.193 & 2.611 & 2.168 & 2.001 & 2.170 & 2.096 & 2.166 & 334.71 & 18237.7 & 0 & 2.00/2 \\
\code{p1--GP--p3} & hexa & 0.1 & 1.972 & 2.163 & 1.999 & 2.215 & 1.973 & 2.213 & 2.615 & 2.141 & 2.001 & 2.165 & 2.071 & 2.160 & 331.70 & 19282.7 & 0 & 2.00/2 \\
\code{p1--GP--p3} & triang. no estr. & 0.0 & 2.292 & 2.131 & 1.992 & 2.080 & 2.009 & 2.079 & 2.600 & 2.164 & 2.002 & 2.209 & 2.268 & 2.200 & 969.28 & 20260.7 & 0 & 2.00/2 \\
\code{p1--GP--p3} & triang. no estr. & 0.1 & 2.345 & 2.112 & 1.992 & 2.106 & 2.034 & 2.104 & 2.552 & 2.095 & 2.000 & 2.134 & 2.167 & 2.129 & 1166.4 & 21785.8 & 0 & 2.00/2 \\
\hline
\code{p2--na--NO} & hexa & 0.0 & 2.128 & 2.568 & 2.142 & 2.142 & 2.143 & 2.143 & 2.020 & 2.183 & 2.025 & 2.025 & 2.026 & 2.026 & 177.84 & 18691.0 & 0 & 2.00/2 \\
\code{p2--na--NO} & hexa & 0.1 & 2.129 & 2.608 & 2.143 & 2.143 & 2.144 & 2.144 & 2.022 & 2.213 & 2.027 & 2.027 & 2.028 & 2.028 & 199.53 & 20243.9 & 0 & 2.00/2 \\
\code{p2--na--NO} & triang. no estr. & 0.0 & 2.058 & 2.376 & 2.090 & 2.090 & 2.086 & 2.086 & 1.991 & 2.077 & 1.998 & 1.998 & 2.002 & 2.002 & 880.50 & 21921.3 & 0 & 2.00/2 \\
\code{p2--na--NO} & triang. no estr. & 0.1 & 2.045 & 2.405 & 2.079 & 2.079 & 2.075 & 2.075 & 1.989 & 2.091 & 1.996 & 1.996 & 1.997 & 1.997 & 1011.0 & 24393.0 & 0 & 2.00/2 \\
\hline
\code{p2--CCp1--p1} & hexa & 0.0 & 2.150 & 2.609 & 2.164 & 2.356 & 2.166 & 2.334 & 2.023 & 2.214 & 2.029 & 2.293 & 2.030 & 2.232 & 253.21 & 19462.0 & 0 & 2.00/2 \\
\code{p2--CCp1--p1} & hexa & 0.1 & 2.153 & 2.652 & 2.167 & 2.356 & 2.170 & 2.335 & 2.026 & 2.251 & 2.031 & 2.293 & 2.032 & 2.235 & 289.62 & 20876.6 & 0 & 2.00/2 \\
\code{p2--CCp1--p1} & triang. no estr. & 0.0 & 2.058 & 2.376 & 2.090 & 2.090 & 2.086 & 2.086 & 1.991 & 2.077 & 1.998 & 1.998 & 2.002 & 2.002 & 1108.5 & 22641.8 & 0 & 2.00/2 \\
\code{p2--CCp1--p1} & triang. no estr. & 0.1 & 2.045 & 2.405 & 2.079 & 2.079 & 2.075 & 2.075 & 1.989 & 2.091 & 1.996 & 1.996 & 1.997 & 1.997 & 1210.2 & 24958.0 & 0 & 2.00/2 \\
\hline
\code{p2--CCp1--p2} & hexa & 0.0 & 2.213 & 2.709 & 2.224 & 2.230 & 2.235 & 2.188 & 2.031 & 2.308 & 2.039 & 2.151 & 2.041 & 2.094 & 256.27 & 19682.3 & 0 & 2.00/3 \\
\code{p2--CCp1--p2} & hexa & 0.1 & 2.225 & 2.760 & 2.235 & 2.230 & 2.248 & 2.187 & 2.036 & 2.369 & 2.043 & 2.151 & 2.046 & 2.094 & 306.21 & 20992.6 & 0 & 2.00/3 \\
\code{p2--CCp1--p2} & triang. no estr. & 0.0 & 2.121 & 2.535 & 2.161 & 2.102 & 2.166 & 2.040 & 1.990 & 2.162 & 2.003 & 2.058 & 2.012 & 2.018 & 1086.9 & 22876.5 & 0 & 2.00/2 \\
\code{p2--CCp1--p2} & triang. no estr. & 0.1 & 2.108 & 2.582 & 2.154 & 2.101 & 2.160 & 2.039 & 1.987 & 2.199 & 1.998 & 2.058 & 2.003 & 2.018 & 1327.0 & 25104.0 & 0 & 2.00/2 \\
\hline
\code{p2--CCp1--p3} & hexa & 0.0 & 2.213 & 2.709 & 2.224 & 2.248 & 2.235 & 2.214 & 2.031 & 2.308 & 2.039 & 2.168 & 2.041 & 2.113 & 260.37 & 19802.0 & 0 & 2.00/3 \\
\code{p2--CCp1--p3} & hexa & 0.1 & 2.225 & 2.760 & 2.235 & 2.248 & 2.248 & 2.214 & 2.036 & 2.369 & 2.043 & 2.168 & 2.046 & 2.113 & 317.77 & 21046.5 & 0 & 2.00/3 \\
\code{p2--CCp1--p3} & triang. no estr. & 0.0 & 2.121 & 2.535 & 2.161 & 2.127 & 2.166 & 2.050 & 1.990 & 2.162 & 2.003 & 2.074 & 2.012 & 2.023 & 1172.1 & 22980.5 & 0 & 2.00/2 \\
\code{p2--CCp1--p3} & triang. no estr. & 0.1 & 2.108 & 2.582 & 2.154 & 2.126 & 2.160 & 2.050 & 1.987 & 2.199 & 1.998 & 2.074 & 2.003 & 2.023 & 1310.8 & 25177.5 & 0 & 2.00/2 \\
\hline
\code{p2--CCp2--p1} & hexa & 0.0 & 2.150 & 2.609 & 2.166 & 3.039 & 2.166 & 2.253 & 2.023 & 2.214 & 2.029 & 2.986 & 2.030 & 2.044 & 293.23 & 19904.3 & 0 & 2.00/2 \\
\code{p2--CCp2--p1} & hexa & 0.1 & 2.153 & 2.653 & 2.170 & 3.040 & 2.170 & 2.269 & 2.026 & 2.251 & 2.032 & 2.989 & 2.032 & 2.050 & 347.11 & 20900.2 & 0 & 2.00/2 \\
\code{p2--CCp2--p1} & triang. no estr. & 0.0 & 2.058 & 2.376 & 2.090 & 2.090 & 2.086 & 2.086 & 1.991 & 2.077 & 1.998 & 1.998 & 2.002 & 2.002 & 1269.2 & 23668.3 & 0 & 2.00/2 \\
\code{p2--CCp2--p1} & triang. no estr. & 0.1 & 2.045 & 2.405 & 2.079 & 2.079 & 2.075 & 2.075 & 1.989 & 2.091 & 1.996 & 1.996 & 1.997 & 1.997 & 1460.7 & 25105.2 & 0 & 2.00/2 \\
\hline
\code{p2--CCp2--p2} & hexa & 0.0 & 2.214 & 2.709 & 2.236 & 3.063 & 2.235 & 2.434 & 2.032 & 2.308 & 2.041 & 2.999 & 2.041 & 2.086 & 302.60 & 20118.8 & 0 & 2.00/3 \\
\code{p2--CCp2--p2} & hexa & 0.1 & 2.226 & 2.761 & 2.250 & 3.063 & 2.249 & 2.475 & 2.037 & 2.371 & 2.046 & 3.000 & 2.047 & 2.105 & 348.76 & 21015.7 & 0 & 2.00/3 \\
\code{p2--CCp2--p2} & triang. no estr. & 0.0 & 2.122 & 2.536 & 2.172 & 3.012 & 2.166 & 2.317 & 1.990 & 2.162 & 2.003 & 2.991 & 2.012 & 2.038 & 1209.5 & 23871.8 & 0 & 2.00/2 \\
\code{p2--CCp2--p2} & triang. no estr. & 0.1 & 2.108 & 2.583 & 2.166 & 3.012 & 2.160 & 2.339 & 1.987 & 2.200 & 1.999 & 2.993 & 2.002 & 2.039 & 1466.0 & 25271.6 & 0 & 2.00/2 \\
\hline
\code{p2--CCp2--p3} & hexa & 0.0 & 2.214 & 2.709 & 2.236 & 3.058 & 2.235 & 2.434 & 2.032 & 2.308 & 2.041 & 2.998 & 2.041 & 2.087 & 297.91 & 20232.9 & 0 & 2.00/3 \\
\code{p2--CCp2--p3} & hexa & 0.1 & 2.226 & 2.761 & 2.250 & 3.058 & 2.249 & 2.474 & 2.037 & 2.371 & 2.046 & 2.999 & 2.047 & 2.106 & 414.45 & 21086.6 & 0 & 2.00/3 \\
\code{p2--CCp2--p3} & triang. no estr. & 0.0 & 2.122 & 2.536 & 2.172 & 3.009 & 2.166 & 2.314 & 1.990 & 2.162 & 2.003 & 2.991 & 2.012 & 2.039 & 1212.5 & 23983.9 & 0 & 2.00/2 \\
\code{p2--CCp2--p3} & triang. no estr. & 0.1 & 2.108 & 2.583 & 2.166 & 3.010 & 2.160 & 2.336 & 1.987 & 2.200 & 1.999 & 2.993 & 2.002 & 2.040 & 1487.1 & 25347.9 & 0 & 2.00/2 \\
\hline
\code{p2--CCp3--p1} & hexa & 0.0 & 2.150 & 2.609 & 2.166 & 3.135 & 2.166 & 2.183 & 2.023 & 2.214 & 2.029 & 1.992 & 2.030 & 2.029 & 420.59 & 21104.2 & 0 & 2.00/2 \\
\code{p2--CCp3--p1} & hexa & 0.1 & 2.153 & 2.653 & 2.170 & 3.206 & 2.170 & 2.186 & 2.026 & 2.251 & 2.032 & 2.002 & 2.032 & 2.031 & 577.79 & 21515.3 & 0 & 2.00/2 \\
\code{p2--CCp3--p1} & triang. no estr. & 0.0 & 2.058 & 2.376 & 2.090 & 2.090 & 2.086 & 2.086 & 1.991 & 2.077 & 1.998 & 1.998 & 2.002 & 2.002 & 1433.7 & 25987.5 & 0 & 2.00/2 \\
\code{p2--CCp3--p1} & triang. no estr. & 0.1 & 2.045 & 2.405 & 2.079 & 2.079 & 2.075 & 2.075 & 1.989 & 2.091 & 1.996 & 1.996 & 1.997 & 1.997 & 1728.1 & 26852.1 & 0 & 2.00/2 \\
\hline
\code{p2--CCp3--p2} & hexa & 0.0 & 2.214 & 2.709 & 2.236 & 3.796 & 2.235 & 2.299 & 2.032 & 2.308 & 2.041 & 2.178 & 2.041 & 2.039 & 452.14 & 21274.2 & 0 & 2.00/3 \\
\code{p2--CCp3--p2} & hexa & 0.1 & 2.226 & 2.761 & 2.249 & 3.882 & 2.249 & 2.314 & 2.037 & 2.371 & 2.046 & 2.315 & 2.046 & 2.044 & 616.38 & 21717.9 & 0 & 2.00/3 \\
\code{p2--CCp3--p2} & triang. no estr. & 0.0 & 2.122 & 2.536 & 2.172 & 3.302 & 2.166 & 2.231 & 1.990 & 2.162 & 2.003 & 1.988 & 2.012 & 2.010 & 1455.9 & 26171.3 & 0 & 2.00/2 \\
\code{p2--CCp3--p2} & triang. no estr. & 0.1 & 2.108 & 2.583 & 2.166 & 3.390 & 2.160 & 2.225 & 1.987 & 2.200 & 1.999 & 2.010 & 2.002 & 2.001 & 1726.2 & 27073.5 & 0 & 2.00/2 \\
\hline
\code{p2--CCp3--p3} & hexa & 0.0 & 2.214 & 2.709 & 2.236 & 3.789 & 2.235 & 2.298 & 2.032 & 2.308 & 2.041 & 2.133 & 2.041 & 2.039 & 449.72 & 21383.6 & 0 & 2.00/3 \\
\code{p2--CCp3--p3} & hexa & 0.1 & 2.226 & 2.761 & 2.249 & 3.877 & 2.249 & 2.313 & 2.037 & 2.371 & 2.046 & 2.262 & 2.046 & 2.044 & 645.79 & 21799.0 & 0 & 2.00/3 \\
\code{p2--CCp3--p3} & triang. no estr. & 0.0 & 2.122 & 2.536 & 2.172 & 3.267 & 2.166 & 2.229 & 1.990 & 2.162 & 2.003 & 1.957 & 2.012 & 2.010 & 1428.8 & 26265.6 & 0 & 2.00/2 \\
\code{p2--CCp3--p3} & triang. no estr. & 0.1 & 2.108 & 2.583 & 2.166 & 3.356 & 2.160 & 2.223 & 1.987 & 2.200 & 1.999 & 1.972 & 2.002 & 2.001 & 1337.8 & 27180.5 & 0 & 2.00/2 \\
\hline
\code{p2--GP--p1} & hexa & 0.0 & 2.150 & 2.609 & 2.003 & 2.689 & 2.102 & 2.694 & 2.023 & 2.214 & 2.001 & 2.304 & 2.011 & 2.312 & 229.19 & 18746.2 & 0 & 2.00/2 \\
\code{p2--GP--p1} & hexa & 0.1 & 2.153 & 2.653 & 2.003 & 2.728 & 2.100 & 2.732 & 2.026 & 2.251 & 2.001 & 2.350 & 2.011 & 2.358 & 250.07 & 20328.7 & 0 & 2.00/2 \\
\code{p2--GP--p1} & triang. no estr. & 0.0 & 2.058 & 2.376 & 2.090 & 2.090 & 2.086 & 2.086 & 1.991 & 2.077 & 1.998 & 1.998 & 2.002 & 2.002 & 944.25 & 21969.6 & 0 & 2.00/2 \\
\code{p2--GP--p1} & triang. no estr. & 0.1 & 2.045 & 2.405 & 2.079 & 2.079 & 2.075 & 2.075 & 1.989 & 2.091 & 1.996 & 1.996 & 1.997 & 1.997 & 1097.3 & 24415.8 & 0 & 2.00/2 \\
\hline
\code{p2--GP--p2} & hexa & 0.0 & 2.214 & 2.709 & 2.003 & 2.924 & 2.114 & 2.923 & 2.032 & 2.308 & 2.000 & 2.644 & 2.007 & 2.650 & 315.69 & 18954.4 & 0 & 2.00/3 \\
\code{p2--GP--p2} & hexa & 0.1 & 2.226 & 2.761 & 2.003 & 2.952 & 2.112 & 2.951 & 2.037 & 2.371 & 2.000 & 2.706 & 2.007 & 2.711 & 388.09 & 20460.4 & 0 & 2.00/3 \\
\code{p2--GP--p2} & triang. no estr. & 0.0 & 2.122 & 2.536 & 1.992 & 2.805 & 2.037 & 2.807 & 1.990 & 2.162 & 1.998 & 2.489 & 2.003 & 2.499 & 1056.9 & 22218.4 & 0 & 2.00/2 \\
\code{p2--GP--p2} & triang. no estr. & 0.1 & 2.108 & 2.583 & 1.992 & 2.840 & 2.035 & 2.842 & 1.987 & 2.200 & 1.998 & 2.553 & 2.002 & 2.561 & 1200.0 & 24559.1 & 0 & 2.00/2 \\
\hline
\code{p2--GP--p3} & hexa & 0.0 & 2.214 & 2.709 & 2.004 & 2.923 & 2.114 & 2.922 & 2.032 & 2.308 & 2.000 & 2.650 & 2.007 & 2.656 & 377.79 & 19059.3 & 0 & 2.00/3 \\
\code{p2--GP--p3} & hexa & 0.1 & 2.226 & 2.761 & 2.004 & 2.950 & 2.113 & 2.950 & 2.037 & 2.371 & 2.000 & 2.711 & 2.007 & 2.716 & 499.25 & 20511.8 & 0 & 2.00/3 \\
\code{p2--GP--p3} & triang. no estr. & 0.0 & 2.122 & 2.536 & 1.993 & 2.809 & 2.037 & 2.812 & 1.990 & 2.162 & 1.998 & 2.500 & 2.003 & 2.510 & 1067.6 & 22304.9 & 0 & 2.00/2 \\
\code{p2--GP--p3} & triang. no estr. & 0.1 & 2.108 & 2.583 & 1.993 & 2.843 & 2.035 & 2.846 & 1.987 & 2.200 & 1.998 & 2.564 & 2.002 & 2.572 & 1257.2 & 24623.2 & 0 & 2.00/2 \\
\hline
\code{p3--na--NO} & hexa & 0.0 & 2.125 & 2.485 & 2.132 & 2.132 & 2.135 & 2.135 & 2.013 & 2.050 & 2.025 & 2.025 & 2.024 & 2.024 & 245.34 & 21804.9 & 0 & 2.00/2 \\
\code{p3--na--NO} & hexa & 0.1 & 2.108 & 2.481 & 2.114 & 2.114 & 2.117 & 2.117 & 2.010 & 2.048 & 2.020 & 2.020 & 2.020 & 2.020 & 377.83 & 24497.7 & 0 & 2.00/2 \\
\code{p3--na--NO} & triang. no estr. & 0.0 & 2.041 & 2.209 & 2.061 & 2.061 & 2.060 & 2.060 & 1.999 & 2.015 & 2.008 & 2.008 & 2.008 & 2.008 & 1086.6 & 27536.0 & 0 & 2.00/2 \\
\code{p3--na--NO} & triang. no estr. & 0.1 & 2.029 & 2.204 & 2.046 & 2.046 & 2.046 & 2.046 & 1.999 & 2.014 & 2.008 & 2.008 & 2.007 & 2.007 & 999.45 & 31524.1 & 0 & 2.00/2 \\
\hline
\code{p3--CCp1--p1} & hexa & 0.0 & 2.185 & 2.634 & 2.188 & 2.354 & 2.201 & 2.339 & 2.020 & 2.072 & 2.038 & 2.290 & 2.040 & 2.248 & 337.38 & 22971.3 & 0 & 2.00/2 \\
\code{p3--CCp1--p1} & hexa & 0.1 & 2.160 & 2.629 & 2.164 & 2.354 & 2.176 & 2.339 & 2.016 & 2.068 & 2.031 & 2.290 & 2.032 & 2.248 & 538.47 & 25276.2 & 0 & 2.00/2 \\
\code{p3--CCp1--p1} & triang. no estr. & 0.0 & 2.041 & 2.209 & 2.061 & 2.061 & 2.060 & 2.060 & 1.999 & 2.015 & 2.008 & 2.008 & 2.008 & 2.008 & 1249.8 & 28336.6 & 0 & 2.00/2 \\
\code{p3--CCp1--p1} & triang. no estr. & 0.1 & 2.029 & 2.204 & 2.046 & 2.046 & 2.046 & 2.046 & 1.999 & 2.014 & 2.008 & 2.008 & 2.007 & 2.007 & 1225.3 & 31983.9 & 0 & 2.00/2 \\
\hline
\code{p3--CCp1--p2} & hexa & 0.0 & 2.620 & 3.309 & 2.640 & 2.229 & 2.603 & 2.184 & 2.117 & 2.267 & 2.264 & 2.150 & 2.211 & 2.094 & 357.27 & 23092.7 & 0 & 2.00/3 \\
\code{p3--CCp1--p2} & hexa & 0.1 & 2.564 & 3.307 & 2.592 & 2.229 & 2.554 & 2.184 & 2.084 & 2.237 & 2.206 & 2.150 & 2.160 & 2.094 & 569.95 & 25382.7 & 0 & 2.00/3 \\
\code{p3--CCp1--p2} & triang. no estr. & 0.0 & 2.396 & 2.791 & 2.510 & 2.101 & 2.460 & 2.038 & 2.003 & 2.026 & 2.041 & 2.058 & 2.014 & 2.017 & 1285.8 & 28499.6 & 0 & 2.00/2 \\
\code{p3--CCp1--p2} & triang. no estr. & 0.1 & 2.330 & 2.777 & 2.443 & 2.101 & 2.397 & 2.038 & 1.999 & 2.022 & 2.031 & 2.058 & 2.006 & 2.017 & 1349.2 & 32108.1 & 0 & 2.00/2 \\
\hline
\code{p3--CCp1--p3} & hexa & 0.0 & 2.620 & 3.309 & 2.640 & 2.247 & 2.603 & 2.210 & 2.117 & 2.267 & 2.264 & 2.167 & 2.211 & 2.112 & 395.31 & 23152.0 & 0 & 2.00/3 \\
\code{p3--CCp1--p3} & hexa & 0.1 & 2.564 & 3.307 & 2.592 & 2.247 & 2.554 & 2.210 & 2.084 & 2.237 & 2.206 & 2.167 & 2.160 & 2.112 & 560.65 & 25436.5 & 0 & 2.00/3 \\
\code{p3--CCp1--p3} & triang. no estr. & 0.0 & 2.396 & 2.791 & 2.509 & 2.126 & 2.460 & 2.048 & 2.003 & 2.026 & 2.041 & 2.074 & 2.014 & 2.023 & 1443.9 & 28577.5 & 0 & 2.00/2 \\
\code{p3--CCp1--p3} & triang. no estr. & 0.1 & 2.330 & 2.777 & 2.443 & 2.126 & 2.397 & 2.048 & 2.000 & 2.023 & 2.031 & 2.074 & 2.006 & 2.023 & 1751.9 & 32158.6 & 0 & 2.00/2 \\
\hline
\code{p3--CCp2--p1} & hexa & 0.0 & 2.187 & 2.637 & 2.200 & 3.042 & 2.203 & 2.642 & 2.020 & 2.073 & 2.041 & 2.999 & 2.040 & 2.257 & 414.93 & 23046.5 & 0 & 2.00/2 \\
\code{p3--CCp2--p1} & hexa & 0.1 & 2.162 & 2.631 & 2.175 & 3.042 & 2.177 & 2.638 & 2.016 & 2.069 & 2.033 & 2.999 & 2.033 & 2.251 & 538.68 & 25254.3 & 0 & 2.00/2 \\
\code{p3--CCp2--p1} & triang. no estr. & 0.0 & 2.041 & 2.209 & 2.061 & 2.061 & 2.060 & 2.060 & 1.999 & 2.015 & 2.008 & 2.008 & 2.008 & 2.008 & 1635.1 & 28534.6 & 0 & 2.00/2 \\
\code{p3--CCp2--p1} & triang. no estr. & 0.1 & 2.029 & 2.204 & 2.046 & 2.046 & 2.046 & 2.046 & 1.999 & 2.014 & 2.008 & 2.008 & 2.007 & 2.007 & 1687.7 & 31978.8 & 0 & 2.00/2 \\
\hline
\code{p3--CCp2--p2} & hexa & 0.0 & 3.539 & 3.992 & 3.284 & 3.063 & 3.303 & 3.110 & 4.323 & 4.278 & 4.078 & 3.002 & 4.085 & 3.055 & 444.41 & 23186.7 & 0 & 2.00/3 \\
\code{p3--CCp2--p2} & hexa & 0.1 & 3.570 & 4.025 & 3.312 & 3.063 & 3.332 & 3.109 & 4.298 & 4.271 & 4.078 & 3.002 & 4.086 & 3.049 & 516.67 & 25373.1 & 0 & 2.00/3 \\
\code{p3--CCp2--p2} & triang. no estr. & 0.0 & 4.057 & 4.096 & 3.842 & 3.013 & 3.861 & 3.085 & 4.475 & 4.165 & 4.251 & 2.997 & 4.258 & 3.023 & 1695.1 & 28749.0 & 0 & 2.00/2 \\
\code{p3--CCp2--p2} & triang. no estr. & 0.1 & 4.077 & 4.105 & 3.880 & 3.013 & 3.901 & 3.082 & 4.459 & 4.138 & 4.198 & 2.997 & 4.214 & 3.022 & 1613.4 & 32141.5 & 0 & 2.00/2 \\
\hline
\code{p3--CCp2--p3} & hexa & 0.0 & 3.539 & 3.992 & 3.284 & 3.058 & 3.303 & 3.102 & 4.323 & 4.278 & 4.078 & 3.001 & 4.085 & 3.050 & 450.49 & 23294.3 & 0 & 2.00/3 \\
\code{p3--CCp2--p3} & hexa & 0.1 & 3.570 & 4.025 & 3.312 & 3.058 & 3.332 & 3.101 & 4.298 & 4.271 & 4.078 & 3.001 & 4.086 & 3.044 & 521.37 & 25432.6 & 0 & 2.00/3 \\
\code{p3--CCp2--p3} & triang. no estr. & 0.0 & 4.057 & 4.096 & 3.842 & 3.011 & 3.861 & 3.072 & 4.476 & 4.165 & 4.251 & 2.997 & 4.259 & 3.019 & 1574.2 & 28818.3 & 0 & 2.00/2 \\
\code{p3--CCp2--p3} & triang. no estr. & 0.1 & 4.077 & 4.105 & 3.880 & 3.011 & 3.901 & 3.070 & 4.459 & 4.138 & 4.198 & 2.997 & 4.214 & 3.017 & 1663.2 & 32205.2 & 0 & 2.00/2 \\
\hline
\code{p3--CCp3--p1} & hexa & 0.0 & 2.186 & 2.637 & 2.200 & 3.993 & 2.202 & 2.225 & 2.020 & 2.073 & 2.041 & 2.812 & 2.040 & 2.040 & 549.23 & 23615.9 & 0 & 2.00/2 \\
\code{p3--CCp3--p1} & hexa & 0.1 & 2.161 & 2.631 & 2.174 & 3.993 & 2.177 & 2.201 & 2.016 & 2.069 & 2.033 & 2.809 & 2.033 & 2.032 & 659.38 & 25283.6 & 0 & 2.00/2 \\
\code{p3--CCp3--p1} & triang. no estr. & 0.0 & 2.041 & 2.209 & 2.061 & 2.061 & 2.060 & 2.060 & 1.999 & 2.015 & 2.008 & 2.008 & 2.008 & 2.008 & 1648.9 & 30662.0 & 0 & 2.00/2 \\
\code{p3--CCp3--p1} & triang. no estr. & 0.1 & 2.029 & 2.204 & 2.046 & 2.046 & 2.046 & 2.046 & 1.999 & 2.014 & 2.008 & 2.008 & 2.007 & 2.007 & 1792.5 & 32119.3 & 0 & 2.00/2 \\
\hline
\code{p3--CCp3--p2} & hexa & 0.0 & 3.538 & 3.993 & 3.290 & 4.231 & 3.304 & 3.441 & 4.323 & 4.279 & 4.079 & 4.284 & 4.086 & 4.100 & 543.80 & 23800.5 & 0 & 2.00/3 \\
\code{p3--CCp3--p2} & hexa & 0.1 & 3.569 & 4.026 & 3.319 & 4.231 & 3.333 & 3.485 & 4.299 & 4.271 & 4.081 & 4.283 & 4.087 & 4.101 & 701.99 & 25398.9 & 0 & 2.00/3 \\
\code{p3--CCp3--p2} & triang. no estr. & 0.0 & 4.058 & 4.096 & 3.846 & 4.109 & 3.863 & 3.934 & 4.476 & 4.165 & 4.253 & 4.049 & 4.258 & 4.241 & 1698.7 & 30891.7 & 0 & 2.00/2 \\
\code{p3--CCp3--p2} & triang. no estr. & 0.1 & 4.078 & 4.105 & 3.885 & 4.109 & 3.903 & 3.989 & 4.460 & 4.139 & 4.202 & 4.049 & 4.214 & 4.191 & 1679.5 & 32279.4 & 0 & 2.00/2 \\
\hline
\code{p3--CCp3--p3} & hexa & 0.0 & 3.538 & 3.993 & 3.290 & 4.241 & 3.304 & 3.438 & 4.323 & 4.279 & 4.079 & 4.302 & 4.086 & 4.100 & 549.65 & 23864.2 & 0 & 2.00/3 \\
\code{p3--CCp3--p3} & hexa & 0.1 & 3.569 & 4.026 & 3.319 & 4.241 & 3.333 & 3.481 & 4.299 & 4.271 & 4.081 & 4.300 & 4.087 & 4.100 & 753.98 & 25461.3 & 0 & 2.00/3 \\
\code{p3--CCp3--p3} & triang. no estr. & 0.0 & 4.057 & 4.096 & 3.846 & 4.114 & 3.862 & 3.932 & 4.476 & 4.165 & 4.254 & 4.050 & 4.259 & 4.242 & 1879.3 & 31010.9 & 0 & 2.00/2 \\
\code{p3--CCp3--p3} & triang. no estr. & 0.1 & 4.077 & 4.105 & 3.884 & 4.114 & 3.903 & 3.986 & 4.460 & 4.139 & 4.202 & 4.050 & 4.214 & 4.192 & 1453.5 & 32339.0 & 0 & 2.00/2 \\
\hline
\code{p3--GP--p1} & hexa & 0.0 & 2.187 & 2.637 & 2.001 & 2.598 & 1.989 & 2.596 & 2.020 & 2.073 & 2.000 & 2.068 & 1.999 & 2.066 & 358.46 & 21855.3 & 0 & 2.00/2 \\
\code{p3--GP--p1} & hexa & 0.1 & 2.161 & 2.631 & 2.001 & 2.592 & 1.988 & 2.589 & 2.016 & 2.069 & 2.000 & 2.060 & 1.999 & 2.059 & 534.78 & 24560.3 & 0 & 2.00/2 \\
\code{p3--GP--p1} & triang. no estr. & 0.0 & 2.041 & 2.209 & 2.061 & 2.061 & 2.060 & 2.060 & 1.999 & 2.015 & 2.008 & 2.008 & 2.008 & 2.008 & 1511.9 & 27558.5 & 0 & 2.00/2 \\
\code{p3--GP--p1} & triang. no estr. & 0.1 & 2.029 & 2.204 & 2.046 & 2.046 & 2.046 & 2.046 & 1.999 & 2.014 & 2.008 & 2.008 & 2.007 & 2.007 & 1438.6 & 31546.0 & 0 & 2.00/2 \\
\hline
\code{p3--GP--p2} & hexa & 0.0 & 3.538 & 3.992 & 1.999 & 4.003 & 1.997 & 4.010 & 4.323 & 4.279 & 2.000 & 4.194 & 2.001 & 4.195 & 430.39 & 21968.9 & 0 & 2.00/3 \\
\code{p3--GP--p2} & hexa & 0.1 & 3.569 & 4.026 & 1.999 & 4.033 & 1.997 & 4.040 & 4.299 & 4.271 & 2.000 & 4.212 & 2.000 & 4.213 & 512.98 & 24660.5 & 0 & 2.00/3 \\
\code{p3--GP--p2} & triang. no estr. & 0.0 & 4.057 & 4.096 & 1.992 & 4.077 & 2.002 & 4.083 & 4.476 & 4.165 & 1.998 & 4.204 & 2.003 & 4.204 & 1565.6 & 27688.4 & 0 & 2.00/2 \\
\code{p3--GP--p2} & triang. no estr. & 0.1 & 4.077 & 4.105 & 1.992 & 4.094 & 2.001 & 4.098 & 4.460 & 4.138 & 1.998 & 4.195 & 2.003 & 4.198 & 1452.0 & 31683.3 & 0 & 2.00/2 \\
\hline
\code{p3--GP--p3} & hexa & 0.0 & 3.538 & 3.992 & 1.999 & 4.000 & 1.998 & 4.008 & 4.323 & 4.279 & 2.000 & 4.198 & 2.001 & 4.200 & 451.44 & 22029.4 & 0 & 2.00/3 \\
\code{p3--GP--p3} & hexa & 0.1 & 3.569 & 4.026 & 1.999 & 4.031 & 1.997 & 4.039 & 4.299 & 4.271 & 2.000 & 4.217 & 2.001 & 4.219 & 669.66 & 24733.5 & 0 & 2.00/3 \\
\code{p3--GP--p3} & triang. no estr. & 0.0 & 4.057 & 4.096 & 1.992 & 4.079 & 2.002 & 4.086 & 4.476 & 4.165 & 1.998 & 4.227 & 2.003 & 4.229 & 1653.7 & 27773.2 & 0 & 2.00/2 \\
\code{p3--GP--p3} & triang. no estr. & 0.1 & 4.077 & 4.105 & 1.992 & 4.098 & 2.002 & 4.103 & 4.460 & 4.138 & 1.998 & 4.218 & 2.003 & 4.222 & 1605.1 & 31748.6 & 0 & 2.00/2 \\
\bottomrule
\end{longtable}
\normalsize


\subsubsection{\code{JFNK}}
\paragraph{$\Delta t=10^{-2}$}
\tiny
\setlength{\tabcolsep}{1.4pt}
\begin{longtable}{lllrrrrrrrrrrrrrrrr}
\caption{Resultados manufacturados 2D para \code{JFNK}, \code{crankNicolson}, masa \code{consistent}, estados \code{RKF45} y $\Delta t=10^{-2}$. El primer bloque de convergencia se ajusta sobre $N=10$--$100$; el segundo sobre $N=70$--$100$. La primera columna registra la terna \code{Vm-states-Iion} efectivamente usada en la corrida. \textbf{RB} es la cantidad de pasos de tiempo que agotaron el presupuesto de iteraciones no lineales, acumulada sobre los 91 tamaños de malla de la fila (455 pasos de tiempo en total por fila). Con \code{nonlinearAcceptUnconverged=false} un paso no convergido se revierte a los valores del inicio del paso mientras el reloj avanza igual. \textbf{$n_{\mathrm{nl}}$} reporta la mediana y el máximo de iteraciones no lineales por paso sobre esos mismos pasos; el presupuesto configurado es \code{nonlinearIterations}$=2000$.}\label{tab:es-results-2d-jfnk-1p0em02}\\
\toprule
& & & \multicolumn{6}{c}{Ajuste $N=10$--$100$} & \multicolumn{6}{c}{Ajuste $N=70$--$100$} & & & & \\
\cmidrule(lr){4-9}\cmidrule(lr){10-15}
\code{Vm-states-Iion} & Tipo de malla & $\alpha$ & $V_m$ & $V_m^G$ & $u_1$ & $u_1^G$ & $u_2$ & $u_2^G$ & $V_m$ & $V_m^G$ & $u_1$ & $u_1^G$ & $u_2$ & $u_2^G$ & $t_{\Sigma}$ [s] & RSS$_{\Sigma}$ [MB] & RB & $n_{\mathrm{nl}}$ \\
\midrule
\endfirsthead
\toprule
& & & \multicolumn{6}{c}{Ajuste $N=10$--$100$} & \multicolumn{6}{c}{Ajuste $N=70$--$100$} & & & & \\
\cmidrule(lr){4-9}\cmidrule(lr){10-15}
\code{Vm-states-Iion} & Tipo de malla & $\alpha$ & $V_m$ & $V_m^G$ & $u_1$ & $u_1^G$ & $u_2$ & $u_2^G$ & $V_m$ & $V_m^G$ & $u_1$ & $u_1^G$ & $u_2$ & $u_2^G$ & $t_{\Sigma}$ [s] & RSS$_{\Sigma}$ [MB] & RB & $n_{\mathrm{nl}}$ \\
\midrule
\endhead
\code{NO--na--NO} & hexa & 0.0 & 1.992 & 1.992 & 2.135 & 2.135 & 2.134 & 2.134 & 1.985 & 1.985 & 2.500 & 2.500 & 2.497 & 2.497 & 9.266 & 16550.6 & 0 & 2.00/3 \\
\code{NO--na--NO} & hexa & 0.1 & 1.849 & 1.849 & 1.612 & 1.612 & 1.618 & 1.618 & 2.278 & 2.278 & 2.159 & 2.159 & 2.163 & 2.163 & 30.513 & 18377.9 & 0 & 2.00/3 \\
\code{NO--na--NO} & triang. no estr. & 0.0 & 0.730 & 0.730 & 0.778 & 0.778 & 0.784 & 0.784 & 0.862 & 0.862 & 1.037 & 1.037 & 1.107 & 1.107 & 78.833 & 17294.1 & 0 & 3.00/3 \\
\code{NO--na--NO} & triang. no estr. & 0.1 & 0.722 & 0.722 & 0.772 & 0.772 & 0.777 & 0.777 & 0.871 & 0.871 & 1.027 & 1.027 & 1.097 & 1.097 & 105.57 & 20811.1 & 0 & 3.00/3 \\
\hline
\code{NO--CCp1--p1} & hexa & 0.0 & 1.993 & 1.993 & 2.136 & 2.354 & 2.134 & 2.333 & 1.985 & 1.985 & 2.504 & 2.289 & 2.497 & 2.242 & 19.227 & 17065.9 & 0 & 2.00/3 \\
\code{NO--CCp1--p1} & hexa & 0.1 & 1.849 & 1.849 & 1.611 & 2.354 & 1.618 & 2.310 & 2.278 & 2.278 & 2.159 & 2.289 & 2.163 & 2.225 & 34.874 & 18494.9 & 0 & 2.00/3 \\
\code{NO--CCp1--p1} & triang. no estr. & 0.0 & 0.730 & 0.730 & 0.778 & 0.778 & 0.784 & 0.784 & 0.862 & 0.862 & 1.037 & 1.037 & 1.107 & 1.107 & 170.83 & 18064.6 & 0 & 3.00/3 \\
\code{NO--CCp1--p1} & triang. no estr. & 0.1 & 0.722 & 0.722 & 0.772 & 0.772 & 0.777 & 0.777 & 0.871 & 0.871 & 1.027 & 1.027 & 1.097 & 1.097 & 182.33 & 20922.1 & 0 & 3.00/3 \\
\hline
\code{NO--CCp1--p2} & hexa & 0.0 & 1.993 & 1.993 & 2.139 & 2.229 & 2.134 & 2.184 & 1.985 & 1.985 & 2.513 & 2.150 & 2.495 & 2.094 & 31.171 & 17365.6 & 0 & 2.00/3 \\
\code{NO--CCp1--p2} & hexa & 0.1 & 1.849 & 1.849 & 1.609 & 2.229 & 1.619 & 2.181 & 2.278 & 2.278 & 2.159 & 2.149 & 2.163 & 2.091 & 39.662 & 18652.2 & 0 & 2.00/3 \\
\code{NO--CCp1--p2} & triang. no estr. & 0.0 & 0.730 & 0.730 & 0.778 & 2.111 & 0.784 & 1.164 & 0.862 & 0.862 & 1.037 & 2.101 & 1.107 & 1.150 & 175.21 & 18431.4 & 0 & 3.00/3 \\
\code{NO--CCp1--p2} & triang. no estr. & 0.1 & 0.722 & 0.722 & 0.772 & 2.112 & 0.777 & 1.172 & 0.871 & 0.871 & 1.027 & 2.105 & 1.097 & 1.143 & 210.74 & 21075.5 & 0 & 3.00/3 \\
\hline
\code{NO--CCp1--p3} & hexa & 0.0 & 1.993 & 1.993 & 2.139 & 2.247 & 2.134 & 2.211 & 1.985 & 1.985 & 2.513 & 2.167 & 2.495 & 2.112 & 30.669 & 17519.9 & 0 & 2.00/3 \\
\code{NO--CCp1--p3} & hexa & 0.1 & 1.849 & 1.849 & 1.609 & 2.247 & 1.619 & 2.206 & 2.278 & 2.278 & 2.159 & 2.166 & 2.163 & 2.109 & 39.208 & 18733.1 & 0 & 2.00/3 \\
\code{NO--CCp1--p3} & triang. no estr. & 0.0 & 0.730 & 0.730 & 0.778 & 2.139 & 0.784 & 1.132 & 0.862 & 0.862 & 1.037 & 2.131 & 1.107 & 1.144 & 153.75 & 18596.0 & 0 & 3.00/3 \\
\code{NO--CCp1--p3} & triang. no estr. & 0.1 & 0.722 & 0.722 & 0.772 & 2.140 & 0.777 & 1.139 & 0.871 & 0.871 & 1.027 & 2.135 & 1.097 & 1.136 & 206.41 & 21132.4 & 0 & 3.00/3 \\
\hline
\code{NO--CCp2--p1} & hexa & 0.0 & 1.992 & 1.992 & 2.141 & 3.042 & 2.132 & 2.624 & 1.985 & 1.985 & 2.517 & 2.999 & 2.491 & 2.603 & 37.647 & 17579.3 & 0 & 2.00/3 \\
\code{NO--CCp2--p1} & hexa & 0.1 & 1.849 & 1.849 & 1.607 & 3.038 & 1.618 & 1.948 & 2.278 & 2.278 & 2.158 & 2.987 & 2.162 & 2.187 & 43.859 & 18732.4 & 0 & 2.00/3 \\
\code{NO--CCp2--p1} & triang. no estr. & 0.0 & 0.730 & 0.730 & 0.778 & 0.778 & 0.784 & 0.784 & 0.862 & 0.862 & 1.037 & 1.037 & 1.107 & 1.107 & 293.34 & 19273.3 & 0 & 3.00/3 \\
\code{NO--CCp2--p1} & triang. no estr. & 0.1 & 0.722 & 0.722 & 0.772 & 0.772 & 0.777 & 0.777 & 0.871 & 0.871 & 1.027 & 1.027 & 1.097 & 1.097 & 279.89 & 21288.3 & 0 & 3.00/3 \\
\hline
\code{NO--CCp2--p2} & hexa & 0.0 & 1.991 & 1.991 & 2.145 & 3.063 & 2.120 & 2.764 & 1.986 & 1.986 & 2.519 & 3.002 & 2.444 & 2.629 & 49.227 & 17885.7 & 0 & 2.00/3 \\
\code{NO--CCp2--p2} & hexa & 0.1 & 1.851 & 1.851 & 1.600 & 3.061 & 1.624 & 2.068 & 2.274 & 2.274 & 2.155 & 2.996 & 2.162 & 2.210 & 46.755 & 19025.5 & 0 & 2.00/3 \\
\code{NO--CCp2--p2} & triang. no estr. & 0.0 & 0.730 & 0.730 & 0.778 & 1.941 & 0.784 & 0.776 & 0.862 & 0.862 & 1.037 & 1.095 & 1.107 & 1.097 & 252.35 & 19638.5 & 0 & 3.00/3 \\
\code{NO--CCp2--p2} & triang. no estr. & 0.1 & 0.722 & 0.722 & 0.772 & 1.957 & 0.777 & 0.773 & 0.871 & 0.871 & 1.027 & 1.090 & 1.097 & 1.087 & 265.88 & 21625.1 & 0 & 3.00/3 \\
\hline
\code{NO--CCp2--p3} & hexa & 0.0 & 1.991 & 1.991 & 2.145 & 3.058 & 2.120 & 2.763 & 1.986 & 1.986 & 2.519 & 3.001 & 2.444 & 2.632 & 48.850 & 18015.0 & 0 & 2.00/3 \\
\code{NO--CCp2--p3} & hexa & 0.1 & 1.851 & 1.851 & 1.600 & 3.056 & 1.624 & 2.069 & 2.274 & 2.274 & 2.155 & 2.995 & 2.162 & 2.211 & 51.083 & 19169.8 & 0 & 2.00/3 \\
\code{NO--CCp2--p3} & triang. no estr. & 0.0 & 0.730 & 0.730 & 0.778 & 1.952 & 0.784 & 0.776 & 0.862 & 0.862 & 1.037 & 1.098 & 1.107 & 1.096 & 337.96 & 19809.5 & 0 & 3.00/3 \\
\code{NO--CCp2--p3} & triang. no estr. & 0.1 & 0.722 & 0.722 & 0.772 & 1.968 & 0.777 & 0.773 & 0.871 & 0.871 & 1.027 & 1.093 & 1.097 & 1.087 & 316.88 & 21810.7 & 0 & 3.00/3 \\
\hline
\code{NO--CCp3--p1} & hexa & 0.0 & 1.992 & 1.992 & 2.139 & 4.063 & 2.132 & 2.165 & 1.985 & 1.985 & 2.517 & 3.183 & 2.491 & 2.491 & 67.019 & 18928.3 & 0 & 2.00/3 \\
\code{NO--CCp3--p1} & hexa & 0.1 & 1.849 & 1.849 & 1.605 & 3.118 & 1.618 & 1.661 & 2.278 & 2.278 & 2.158 & 2.340 & 2.163 & 2.189 & 85.420 & 20094.7 & 0 & 2.00/3 \\
\code{NO--CCp3--p1} & triang. no estr. & 0.0 & 0.730 & 0.730 & 0.778 & 0.778 & 0.784 & 0.784 & 0.862 & 0.862 & 1.037 & 1.037 & 1.107 & 1.107 & 435.49 & 21779.7 & 0 & 3.00/3 \\
\code{NO--CCp3--p1} & triang. no estr. & 0.1 & 0.722 & 0.722 & 0.772 & 0.772 & 0.777 & 0.777 & 0.871 & 0.871 & 1.027 & 1.027 & 1.097 & 1.097 & 353.60 & 23790.8 & 0 & 3.00/3 \\
\hline
\code{NO--CCp3--p2} & hexa & 0.0 & 1.992 & 1.992 & 2.141 & 4.277 & 2.120 & 2.216 & 1.985 & 1.985 & 2.519 & 4.233 & 2.444 & 2.443 & 71.573 & 19210.9 & 0 & 2.00/3 \\
\code{NO--CCp3--p2} & hexa & 0.1 & 1.851 & 1.851 & 1.594 & 3.421 & 1.622 & 1.722 & 2.274 & 2.274 & 2.155 & 2.589 & 2.162 & 2.221 & 74.356 & 20366.7 & 0 & 2.00/3 \\
\code{NO--CCp3--p2} & triang. no estr. & 0.0 & 0.730 & 0.730 & 0.778 & 1.071 & 0.784 & 0.796 & 0.862 & 0.862 & 1.037 & 1.037 & 1.107 & 1.109 & 463.11 & 22135.9 & 0 & 3.00/3 \\
\code{NO--CCp3--p2} & triang. no estr. & 0.1 & 0.722 & 0.722 & 0.772 & 1.081 & 0.777 & 0.798 & 0.871 & 0.871 & 1.027 & 1.028 & 1.097 & 1.100 & 362.07 & 24132.6 & 0 & 3.00/3 \\
\hline
\code{NO--CCp3--p3} & hexa & 0.0 & 1.992 & 1.992 & 2.141 & 4.293 & 2.120 & 2.214 & 1.985 & 1.985 & 2.519 & 4.254 & 2.444 & 2.443 & 61.808 & 19363.5 & 0 & 2.00/3 \\
\code{NO--CCp3--p3} & hexa & 0.1 & 1.851 & 1.851 & 1.594 & 3.409 & 1.622 & 1.721 & 2.274 & 2.274 & 2.155 & 2.578 & 2.162 & 2.221 & 74.433 & 20513.0 & 0 & 2.00/3 \\
\code{NO--CCp3--p3} & triang. no estr. & 0.0 & 0.730 & 0.730 & 0.778 & 1.054 & 0.784 & 0.796 & 0.862 & 0.862 & 1.037 & 1.037 & 1.107 & 1.109 & 399.21 & 22286.3 & 0 & 3.00/3 \\
\code{NO--CCp3--p3} & triang. no estr. & 0.1 & 0.722 & 0.722 & 0.772 & 1.064 & 0.777 & 0.797 & 0.871 & 0.871 & 1.027 & 1.028 & 1.097 & 1.100 & 347.65 & 24284.2 & 0 & 3.00/3 \\
\hline
\code{NO--GP--p1} & hexa & 0.0 & 1.992 & 1.992 & 2.000 & 2.057 & 2.005 & 2.061 & 1.985 & 1.985 & 2.002 & 2.031 & 2.023 & 2.031 & 13.090 & 16727.8 & 0 & 2.00/3 \\
\code{NO--GP--p1} & hexa & 0.1 & 1.849 & 1.849 & 2.000 & 2.049 & 1.997 & 2.053 & 2.278 & 2.278 & 2.002 & 2.039 & 2.030 & 2.039 & 24.074 & 18448.5 & 0 & 2.00/3 \\
\code{NO--GP--p1} & triang. no estr. & 0.0 & 0.730 & 0.730 & 0.778 & 0.778 & 0.784 & 0.784 & 0.862 & 0.862 & 1.037 & 1.037 & 1.107 & 1.107 & 95.302 & 17470.9 & 0 & 3.00/3 \\
\code{NO--GP--p1} & triang. no estr. & 0.1 & 0.722 & 0.722 & 0.772 & 0.772 & 0.777 & 0.777 & 0.871 & 0.871 & 1.027 & 1.027 & 1.097 & 1.097 & 132.10 & 20899.6 & 0 & 3.00/3 \\
\hline
\code{NO--GP--p2} & hexa & 0.0 & 1.992 & 1.992 & 1.999 & 2.015 & 2.002 & 2.017 & 1.986 & 1.986 & 2.001 & 2.010 & 2.011 & 2.010 & 29.398 & 17030.0 & 0 & 2.00/3 \\
\code{NO--GP--p2} & hexa & 0.1 & 1.851 & 1.851 & 1.999 & 2.015 & 2.001 & 2.017 & 2.274 & 2.274 & 2.001 & 2.011 & 2.013 & 2.011 & 38.154 & 18577.5 & 0 & 2.00/3 \\
\code{NO--GP--p2} & triang. no estr. & 0.0 & 0.730 & 0.730 & 1.878 & 1.418 & 1.102 & 1.411 & 0.862 & 0.862 & 1.678 & 1.219 & 1.112 & 1.269 & 195.38 & 17820.7 & 0 & 3.00/3 \\
\code{NO--GP--p2} & triang. no estr. & 0.1 & 0.722 & 0.722 & 1.880 & 1.427 & 1.107 & 1.420 & 0.871 & 0.871 & 1.682 & 1.220 & 1.104 & 1.270 & 175.66 & 21048.0 & 0 & 3.00/3 \\
\hline
\code{NO--GP--p3} & hexa & 0.0 & 1.992 & 1.992 & 1.999 & 2.019 & 2.002 & 2.021 & 1.986 & 1.986 & 2.001 & 2.011 & 2.011 & 2.011 & 48.439 & 17196.3 & 0 & 2.00/3 \\
\code{NO--GP--p3} & hexa & 0.1 & 1.851 & 1.851 & 1.999 & 2.019 & 2.001 & 2.021 & 2.274 & 2.274 & 2.001 & 2.013 & 2.013 & 2.013 & 49.173 & 18656.2 & 0 & 2.00/3 \\
\code{NO--GP--p3} & triang. no estr. & 0.0 & 0.730 & 0.730 & 1.878 & 1.389 & 1.102 & 1.384 & 0.862 & 0.862 & 1.678 & 1.206 & 1.112 & 1.258 & 201.79 & 17995.4 & 0 & 3.00/3 \\
\code{NO--GP--p3} & triang. no estr. & 0.1 & 0.722 & 0.722 & 1.881 & 1.398 & 1.107 & 1.392 & 0.871 & 0.871 & 1.682 & 1.207 & 1.104 & 1.259 & 233.70 & 21114.5 & 0 & 3.00/3 \\
\hline
\code{p1--na--NO} & hexa & 0.0 & 1.911 & 2.146 & 1.616 & 1.616 & 1.626 & 1.626 & 2.612 & 2.168 & 2.381 & 2.381 & 2.387 & 2.387 & 77.310 & 17878.7 & 0 & 2.00/3 \\
\code{p1--na--NO} & hexa & 0.1 & 1.971 & 2.163 & 1.677 & 1.677 & 1.687 & 1.687 & 2.614 & 2.141 & 2.416 & 2.416 & 2.422 & 2.422 & 99.193 & 19092.2 & 0 & 2.00/3 \\
\code{p1--na--NO} & triang. no estr. & 0.0 & 2.271 & 2.126 & 2.004 & 2.004 & 2.002 & 2.002 & 2.384 & 2.115 & 2.552 & 2.552 & 2.549 & 2.549 & 627.40 & 19950.2 & 0 & 2.00/3 \\
\code{p1--na--NO} & triang. no estr. & 0.1 & 2.301 & 2.106 & 2.092 & 2.092 & 2.090 & 2.090 & 2.222 & 2.054 & 2.524 & 2.524 & 2.524 & 2.524 & 630.44 & 21567.4 & 0 & 2.00/3 \\
\hline
\code{p1--CCp1--p1} & hexa & 0.0 & 1.911 & 2.146 & 1.616 & 2.349 & 1.626 & 1.955 & 2.612 & 2.168 & 2.381 & 2.294 & 2.387 & 2.369 & 87.842 & 18383.4 & 0 & 2.00/3 \\
\code{p1--CCp1--p1} & hexa & 0.1 & 1.971 & 2.163 & 1.677 & 2.351 & 1.687 & 2.045 & 2.614 & 2.141 & 2.416 & 2.293 & 2.422 & 2.383 & 100.72 & 19298.5 & 0 & 2.00/3 \\
\code{p1--CCp1--p1} & triang. no estr. & 0.0 & 2.271 & 2.126 & 2.004 & 2.004 & 2.002 & 2.002 & 2.384 & 2.115 & 2.552 & 2.552 & 2.549 & 2.549 & 726.75 & 20592.0 & 0 & 2.00/3 \\
\code{p1--CCp1--p1} & triang. no estr. & 0.1 & 2.301 & 2.106 & 2.092 & 2.092 & 2.090 & 2.090 & 2.222 & 2.054 & 2.524 & 2.524 & 2.524 & 2.524 & 626.27 & 21799.9 & 0 & 2.00/3 \\
\hline
\code{p1--CCp1--p2} & hexa & 0.0 & 1.911 & 2.146 & 1.616 & 2.228 & 1.626 & 2.057 & 2.612 & 2.168 & 2.381 & 2.153 & 2.387 & 2.232 & 87.717 & 18666.6 & 0 & 2.00/3 \\
\code{p1--CCp1--p2} & hexa & 0.1 & 1.971 & 2.163 & 1.676 & 2.228 & 1.687 & 2.104 & 2.614 & 2.141 & 2.416 & 2.152 & 2.422 & 2.208 & 93.410 & 19426.6 & 0 & 2.00/3 \\
\code{p1--CCp1--p2} & triang. no estr. & 0.0 & 2.271 & 2.126 & 2.004 & 2.102 & 2.002 & 2.036 & 2.384 & 2.115 & 2.552 & 2.063 & 2.549 & 2.285 & 669.65 & 20937.5 & 0 & 2.00/3 \\
\code{p1--CCp1--p2} & triang. no estr. & 0.1 & 2.302 & 2.106 & 2.092 & 2.102 & 2.090 & 2.066 & 2.222 & 2.054 & 2.524 & 2.061 & 2.524 & 2.190 & 586.52 & 22037.9 & 0 & 2.00/3 \\
\hline
\code{p1--CCp1--p3} & hexa & 0.0 & 1.911 & 2.146 & 1.616 & 2.246 & 1.626 & 2.063 & 2.612 & 2.168 & 2.381 & 2.170 & 2.387 & 2.255 & 91.858 & 18812.1 & 0 & 2.00/3 \\
\code{p1--CCp1--p3} & hexa & 0.1 & 1.971 & 2.163 & 1.676 & 2.246 & 1.687 & 2.115 & 2.614 & 2.141 & 2.416 & 2.169 & 2.422 & 2.234 & 97.556 & 19515.0 & 0 & 2.00/3 \\
\code{p1--CCp1--p3} & triang. no estr. & 0.0 & 2.271 & 2.126 & 2.004 & 2.127 & 2.002 & 2.040 & 2.384 & 2.115 & 2.552 & 2.080 & 2.549 & 2.315 & 696.72 & 21106.2 & 0 & 2.00/3 \\
\code{p1--CCp1--p3} & triang. no estr. & 0.1 & 2.302 & 2.106 & 2.092 & 2.127 & 2.090 & 2.075 & 2.222 & 2.054 & 2.524 & 2.078 & 2.524 & 2.217 & 580.06 & 22150.4 & 0 & 2.00/3 \\
\hline
\code{p1--CCp2--p1} & hexa & 0.0 & 1.911 & 2.146 & 1.616 & 2.904 & 1.626 & 1.665 & 2.612 & 2.168 & 2.381 & 2.783 & 2.387 & 2.395 & 101.76 & 18900.9 & 0 & 2.00/3 \\
\code{p1--CCp2--p1} & hexa & 0.1 & 1.971 & 2.163 & 1.676 & 2.954 & 1.687 & 1.736 & 2.614 & 2.141 & 2.416 & 2.859 & 2.422 & 2.431 & 100.97 & 19452.8 & 0 & 2.00/3 \\
\code{p1--CCp2--p1} & triang. no estr. & 0.0 & 2.271 & 2.126 & 2.004 & 2.004 & 2.002 & 2.002 & 2.384 & 2.115 & 2.552 & 2.552 & 2.549 & 2.549 & 734.25 & 21779.0 & 0 & 2.00/3 \\
\code{p1--CCp2--p1} & triang. no estr. & 0.1 & 2.301 & 2.106 & 2.092 & 2.092 & 2.090 & 2.090 & 2.222 & 2.054 & 2.524 & 2.524 & 2.524 & 2.524 & 603.52 & 22409.0 & 0 & 2.00/3 \\
\hline
\code{p1--CCp2--p2} & hexa & 0.0 & 1.912 & 2.146 & 1.616 & 2.990 & 1.627 & 1.688 & 2.612 & 2.168 & 2.381 & 2.882 & 2.387 & 2.400 & 116.50 & 19176.4 & 0 & 2.00/3 \\
\code{p1--CCp2--p2} & hexa & 0.1 & 1.971 & 2.163 & 1.676 & 3.018 & 1.687 & 1.765 & 2.613 & 2.141 & 2.416 & 2.931 & 2.422 & 2.437 & 100.84 & 19694.5 & 0 & 2.00/3 \\
\code{p1--CCp2--p2} & triang. no estr. & 0.0 & 2.271 & 2.126 & 2.004 & 2.677 & 2.002 & 1.977 & 2.384 & 2.115 & 2.552 & 2.669 & 2.549 & 2.528 & 729.70 & 22122.6 & 0 & 2.00/3 \\
\code{p1--CCp2--p2} & triang. no estr. & 0.1 & 2.301 & 2.106 & 2.092 & 2.799 & 2.090 & 2.062 & 2.222 & 2.054 & 2.524 & 2.735 & 2.523 & 2.511 & 628.08 & 22707.1 & 0 & 2.00/3 \\
\hline
\code{p1--CCp2--p3} & hexa & 0.0 & 1.912 & 2.146 & 1.616 & 2.987 & 1.627 & 1.688 & 2.612 & 2.168 & 2.381 & 2.884 & 2.387 & 2.400 & 116.33 & 19324.3 & 0 & 2.00/3 \\
\code{p1--CCp2--p3} & hexa & 0.1 & 1.971 & 2.163 & 1.676 & 3.015 & 1.687 & 1.765 & 2.613 & 2.141 & 2.416 & 2.932 & 2.422 & 2.437 & 104.74 & 19816.2 & 0 & 2.00/3 \\
\code{p1--CCp2--p3} & triang. no estr. & 0.0 & 2.271 & 2.126 & 2.004 & 2.681 & 2.002 & 1.977 & 2.384 & 2.115 & 2.552 & 2.672 & 2.549 & 2.527 & 750.46 & 22285.7 & 0 & 2.00/3 \\
\code{p1--CCp2--p3} & triang. no estr. & 0.1 & 2.301 & 2.106 & 2.092 & 2.802 & 2.090 & 2.062 & 2.222 & 2.054 & 2.524 & 2.740 & 2.523 & 2.510 & 622.87 & 22859.8 & 0 & 2.00/3 \\
\hline
\code{p1--CCp3--p1} & hexa & 0.0 & 1.911 & 2.146 & 1.616 & 2.034 & 1.626 & 1.630 & 2.612 & 2.168 & 2.381 & 2.378 & 2.387 & 2.407 & 147.06 & 20233.9 & 0 & 2.00/3 \\
\code{p1--CCp3--p1} & hexa & 0.1 & 1.971 & 2.163 & 1.677 & 2.165 & 1.687 & 1.692 & 2.614 & 2.141 & 2.416 & 2.409 & 2.422 & 2.443 & 125.22 & 20735.8 & 0 & 2.00/3 \\
\code{p1--CCp3--p1} & triang. no estr. & 0.0 & 2.271 & 2.126 & 2.004 & 2.004 & 2.002 & 2.002 & 2.384 & 2.115 & 2.552 & 2.552 & 2.549 & 2.549 & 800.43 & 24287.5 & 0 & 2.00/3 \\
\code{p1--CCp3--p1} & triang. no estr. & 0.1 & 2.301 & 2.106 & 2.092 & 2.092 & 2.090 & 2.090 & 2.222 & 2.054 & 2.524 & 2.524 & 2.524 & 2.524 & 683.45 & 24737.2 & 0 & 2.00/3 \\
\hline
\code{p1--CCp3--p2} & hexa & 0.0 & 1.912 & 2.146 & 1.617 & 2.239 & 1.627 & 1.634 & 2.612 & 2.168 & 2.381 & 2.377 & 2.387 & 2.426 & 138.21 & 20509.6 & 0 & 2.00/3 \\
\code{p1--CCp3--p2} & hexa & 0.1 & 1.971 & 2.163 & 1.677 & 2.394 & 1.688 & 1.698 & 2.613 & 2.141 & 2.416 & 2.402 & 2.422 & 2.464 & 125.55 & 21018.0 & 0 & 2.00/3 \\
\code{p1--CCp3--p2} & triang. no estr. & 0.0 & 2.271 & 2.126 & 2.004 & 2.031 & 2.002 & 1.964 & 2.384 & 2.115 & 2.552 & 2.527 & 2.549 & 2.528 & 782.76 & 24626.8 & 0 & 2.00/3 \\
\code{p1--CCp3--p2} & triang. no estr. & 0.1 & 2.301 & 2.106 & 2.092 & 2.141 & 2.090 & 2.052 & 2.222 & 2.054 & 2.524 & 2.511 & 2.523 & 2.516 & 735.35 & 25077.5 & 0 & 2.00/3 \\
\hline
\code{p1--CCp3--p3} & hexa & 0.0 & 1.912 & 2.146 & 1.617 & 2.225 & 1.627 & 1.634 & 2.612 & 2.168 & 2.381 & 2.376 & 2.387 & 2.426 & 133.40 & 20650.3 & 0 & 2.00/3 \\
\code{p1--CCp3--p3} & hexa & 0.1 & 1.971 & 2.163 & 1.677 & 2.379 & 1.688 & 1.697 & 2.613 & 2.141 & 2.416 & 2.402 & 2.422 & 2.463 & 133.26 & 21151.4 & 0 & 2.00/3 \\
\code{p1--CCp3--p3} & triang. no estr. & 0.0 & 2.271 & 2.126 & 2.004 & 2.027 & 2.002 & 1.964 & 2.384 & 2.115 & 2.552 & 2.527 & 2.549 & 2.528 & 745.42 & 24787.5 & 0 & 2.00/3 \\
\code{p1--CCp3--p3} & triang. no estr. & 0.1 & 2.301 & 2.106 & 2.092 & 2.136 & 2.090 & 2.052 & 2.222 & 2.054 & 2.524 & 2.512 & 2.523 & 2.516 & 764.59 & 25234.9 & 0 & 2.00/3 \\
\hline
\code{p1--GP--p1} & hexa & 0.0 & 1.911 & 2.146 & 1.998 & 2.154 & 1.838 & 2.160 & 2.612 & 2.168 & 2.007 & 2.316 & 2.258 & 2.316 & 81.233 & 17974.1 & 0 & 2.00/3 \\
\code{p1--GP--p1} & hexa & 0.1 & 1.971 & 2.163 & 1.999 & 2.221 & 1.898 & 2.227 & 2.614 & 2.141 & 2.005 & 2.319 & 2.234 & 2.318 & 102.12 & 19138.4 & 0 & 2.00/3 \\
\code{p1--GP--p1} & triang. no estr. & 0.0 & 2.271 & 2.126 & 2.004 & 2.004 & 2.002 & 2.002 & 2.384 & 2.115 & 2.552 & 2.552 & 2.549 & 2.549 & 564.30 & 20001.6 & 0 & 2.00/3 \\
\code{p1--GP--p1} & triang. no estr. & 0.1 & 2.301 & 2.106 & 2.092 & 2.092 & 2.090 & 2.090 & 2.222 & 2.054 & 2.524 & 2.524 & 2.524 & 2.524 & 637.92 & 21601.4 & 0 & 2.00/3 \\
\hline
\code{p1--GP--p2} & hexa & 0.0 & 1.912 & 2.146 & 1.999 & 2.182 & 1.951 & 2.181 & 2.612 & 2.168 & 2.002 & 2.161 & 2.105 & 2.157 & 98.467 & 18270.7 & 0 & 2.00/3 \\
\code{p1--GP--p2} & hexa & 0.1 & 1.971 & 2.163 & 1.999 & 2.201 & 1.975 & 2.199 & 2.613 & 2.141 & 2.001 & 2.156 & 2.080 & 2.151 & 101.53 & 19255.2 & 0 & 2.00/3 \\
\code{p1--GP--p2} & triang. no estr. & 0.0 & 2.271 & 2.126 & 1.992 & 2.069 & 2.012 & 2.068 & 2.384 & 2.115 & 2.004 & 2.184 & 2.284 & 2.176 & 612.59 & 20336.8 & 0 & 2.00/3 \\
\code{p1--GP--p2} & triang. no estr. & 0.1 & 2.301 & 2.106 & 1.992 & 2.089 & 2.038 & 2.087 & 2.222 & 2.054 & 2.002 & 2.117 & 2.186 & 2.113 & 623.45 & 21725.5 & 0 & 2.00/3 \\
\hline
\code{p1--GP--p3} & hexa & 0.0 & 1.912 & 2.146 & 1.999 & 2.196 & 1.951 & 2.195 & 2.612 & 2.168 & 2.002 & 2.177 & 2.105 & 2.173 & 102.50 & 18421.3 & 0 & 2.00/3 \\
\code{p1--GP--p3} & hexa & 0.1 & 1.971 & 2.163 & 1.999 & 2.217 & 1.975 & 2.215 & 2.613 & 2.141 & 2.001 & 2.172 & 2.080 & 2.167 & 109.22 & 19318.6 & 0 & 2.00/3 \\
\code{p1--GP--p3} & triang. no estr. & 0.0 & 2.271 & 2.126 & 1.992 & 2.082 & 2.012 & 2.081 & 2.384 & 2.115 & 2.004 & 2.219 & 2.284 & 2.210 & 633.00 & 20513.5 & 0 & 2.00/3 \\
\code{p1--GP--p3} & triang. no estr. & 0.1 & 2.301 & 2.106 & 1.992 & 2.109 & 2.038 & 2.106 & 2.222 & 2.054 & 2.002 & 2.144 & 2.186 & 2.139 & 616.82 & 21804.6 & 0 & 2.00/3 \\
\hline
\code{p2--na--NO} & hexa & 0.0 & 2.127 & 2.567 & 2.176 & 2.176 & 2.176 & 2.176 & 2.016 & 2.179 & 2.139 & 2.139 & 2.139 & 2.139 & 93.266 & 18711.6 & 0 & 2.00/3 \\
\code{p2--na--NO} & hexa & 0.1 & 2.128 & 2.607 & 2.180 & 2.180 & 2.181 & 2.181 & 2.018 & 2.209 & 2.153 & 2.153 & 2.153 & 2.153 & 97.913 & 20308.0 & 0 & 2.00/3 \\
\code{p2--na--NO} & triang. no estr. & 0.0 & 2.056 & 2.374 & 2.166 & 2.166 & 2.161 & 2.161 & 1.982 & 2.068 & 2.260 & 2.260 & 2.263 & 2.263 & 599.62 & 21949.5 & 0 & 2.00/3 \\
\code{p2--na--NO} & triang. no estr. & 0.1 & 2.043 & 2.402 & 2.162 & 2.162 & 2.158 & 2.158 & 1.980 & 2.081 & 2.286 & 2.286 & 2.285 & 2.285 & 576.14 & 24413.9 & 0 & 2.00/3 \\
\hline
\code{p2--CCp1--p1} & hexa & 0.0 & 2.149 & 2.608 & 2.201 & 2.356 & 2.203 & 2.337 & 2.019 & 2.209 & 2.154 & 2.292 & 2.154 & 2.244 & 121.17 & 19564.2 & 0 & 2.00/3 \\
\code{p2--CCp1--p1} & hexa & 0.1 & 2.151 & 2.651 & 2.208 & 2.356 & 2.211 & 2.338 & 2.021 & 2.247 & 2.171 & 2.291 & 2.172 & 2.246 & 108.40 & 20884.4 & 0 & 2.00/3 \\
\code{p2--CCp1--p1} & triang. no estr. & 0.0 & 2.056 & 2.374 & 2.166 & 2.166 & 2.161 & 2.161 & 1.982 & 2.068 & 2.260 & 2.260 & 2.263 & 2.263 & 632.48 & 22835.8 & 0 & 2.00/3 \\
\code{p2--CCp1--p1} & triang. no estr. & 0.1 & 2.043 & 2.402 & 2.162 & 2.162 & 2.158 & 2.158 & 1.980 & 2.081 & 2.286 & 2.286 & 2.285 & 2.285 & 681.24 & 25000.8 & 0 & 2.00/3 \\
\hline
\code{p2--CCp1--p2} & hexa & 0.0 & 2.212 & 2.708 & 2.272 & 2.230 & 2.283 & 2.188 & 2.026 & 2.303 & 2.199 & 2.150 & 2.206 & 2.095 & 117.83 & 19824.1 & 0 & 2.00/3 \\
\code{p2--CCp1--p2} & hexa & 0.1 & 2.223 & 2.759 & 2.291 & 2.229 & 2.304 & 2.188 & 2.030 & 2.364 & 2.229 & 2.150 & 2.238 & 2.095 & 120.32 & 21014.9 & 0 & 2.00/3 \\
\code{p2--CCp1--p2} & triang. no estr. & 0.0 & 2.118 & 2.532 & 2.273 & 2.101 & 2.279 & 2.040 & 1.977 & 2.149 & 2.394 & 2.057 & 2.419 & 2.018 & 655.10 & 23158.1 & 0 & 2.00/3 \\
\code{p2--CCp1--p2} & triang. no estr. & 0.1 & 2.104 & 2.578 & 2.283 & 2.101 & 2.292 & 2.039 & 1.972 & 2.186 & 2.455 & 2.057 & 2.479 & 2.018 & 606.77 & 25195.1 & 0 & 2.00/3 \\
\hline
\code{p2--CCp1--p3} & hexa & 0.0 & 2.212 & 2.708 & 2.272 & 2.248 & 2.283 & 2.215 & 2.026 & 2.303 & 2.199 & 2.167 & 2.206 & 2.114 & 118.26 & 19938.4 & 0 & 2.00/3 \\
\code{p2--CCp1--p3} & hexa & 0.1 & 2.223 & 2.759 & 2.291 & 2.248 & 2.304 & 2.214 & 2.030 & 2.364 & 2.229 & 2.167 & 2.238 & 2.114 & 132.47 & 21079.0 & 0 & 2.00/3 \\
\code{p2--CCp1--p3} & triang. no estr. & 0.0 & 2.118 & 2.532 & 2.273 & 2.126 & 2.279 & 2.050 & 1.977 & 2.149 & 2.394 & 2.074 & 2.419 & 2.023 & 655.96 & 23321.8 & 0 & 2.00/3 \\
\code{p2--CCp1--p3} & triang. no estr. & 0.1 & 2.104 & 2.578 & 2.283 & 2.126 & 2.292 & 2.050 & 1.972 & 2.186 & 2.455 & 2.074 & 2.479 & 2.023 & 592.54 & 25287.0 & 0 & 2.00/3 \\
\hline
\code{p2--CCp2--p1} & hexa & 0.0 & 2.149 & 2.608 & 2.203 & 3.039 & 2.203 & 2.289 & 2.019 & 2.210 & 2.155 & 2.989 & 2.155 & 2.169 & 126.07 & 20011.9 & 0 & 2.00/3 \\
\code{p2--CCp2--p1} & hexa & 0.1 & 2.152 & 2.651 & 2.211 & 3.040 & 2.211 & 2.310 & 2.021 & 2.247 & 2.173 & 2.991 & 2.172 & 2.189 & 126.42 & 20929.3 & 0 & 2.00/3 \\
\code{p2--CCp2--p1} & triang. no estr. & 0.0 & 2.056 & 2.374 & 2.166 & 2.166 & 2.161 & 2.161 & 1.982 & 2.068 & 2.260 & 2.260 & 2.263 & 2.263 & 714.55 & 23917.2 & 0 & 2.00/3 \\
\code{p2--CCp2--p1} & triang. no estr. & 0.1 & 2.043 & 2.402 & 2.162 & 2.162 & 2.158 & 2.158 & 1.980 & 2.081 & 2.286 & 2.286 & 2.285 & 2.285 & 625.61 & 25225.6 & 0 & 2.00/3 \\
\hline
\code{p2--CCp2--p2} & hexa & 0.0 & 2.212 & 2.708 & 2.284 & 3.063 & 2.283 & 2.480 & 2.026 & 2.303 & 2.206 & 3.000 & 2.206 & 2.250 & 132.79 & 20290.6 & 0 & 2.00/3 \\
\code{p2--CCp2--p2} & hexa & 0.1 & 2.224 & 2.760 & 2.305 & 3.063 & 2.305 & 2.527 & 2.030 & 2.365 & 2.239 & 3.001 & 2.239 & 2.294 & 132.29 & 21085.1 & 0 & 2.00/3 \\
\code{p2--CCp2--p2} & triang. no estr. & 0.0 & 2.118 & 2.532 & 2.286 & 3.012 & 2.279 & 2.428 & 1.977 & 2.149 & 2.413 & 2.995 & 2.420 & 2.443 & 735.63 & 24244.0 & 0 & 2.00/3 \\
\code{p2--CCp2--p2} & triang. no estr. & 0.1 & 2.104 & 2.579 & 2.300 & 3.012 & 2.293 & 2.467 & 1.972 & 2.186 & 2.481 & 2.996 & 2.482 & 2.511 & 603.12 & 25454.5 & 0 & 2.00/3 \\
\hline
\code{p2--CCp2--p3} & hexa & 0.0 & 2.212 & 2.708 & 2.284 & 3.058 & 2.283 & 2.480 & 2.026 & 2.303 & 2.206 & 2.999 & 2.206 & 2.251 & 135.22 & 20417.3 & 0 & 2.00/3 \\
\code{p2--CCp2--p3} & hexa & 0.1 & 2.224 & 2.760 & 2.305 & 3.058 & 2.305 & 2.527 & 2.030 & 2.365 & 2.239 & 3.000 & 2.239 & 2.296 & 130.50 & 21160.4 & 0 & 2.00/3 \\
\code{p2--CCp2--p3} & triang. no estr. & 0.0 & 2.118 & 2.532 & 2.286 & 3.010 & 2.279 & 2.425 & 1.977 & 2.149 & 2.413 & 2.995 & 2.420 & 2.443 & 697.92 & 24412.9 & 0 & 2.00/3 \\
\code{p2--CCp2--p3} & triang. no estr. & 0.1 & 2.104 & 2.579 & 2.300 & 3.010 & 2.293 & 2.465 & 1.972 & 2.186 & 2.481 & 2.996 & 2.482 & 2.512 & 598.37 & 25581.7 & 0 & 2.00/3 \\
\hline
\code{p2--CCp3--p1} & hexa & 0.0 & 2.149 & 2.608 & 2.203 & 3.175 & 2.203 & 2.220 & 2.019 & 2.210 & 2.155 & 2.123 & 2.155 & 2.154 & 148.20 & 21250.7 & 0 & 2.00/3 \\
\code{p2--CCp3--p1} & hexa & 0.1 & 2.152 & 2.651 & 2.211 & 3.250 & 2.211 & 2.227 & 2.021 & 2.247 & 2.173 & 2.151 & 2.172 & 2.171 & 137.68 & 21647.8 & 0 & 2.00/3 \\
\code{p2--CCp3--p1} & triang. no estr. & 0.0 & 2.056 & 2.374 & 2.166 & 2.166 & 2.161 & 2.161 & 1.982 & 2.068 & 2.260 & 2.260 & 2.263 & 2.263 & 753.15 & 26381.9 & 0 & 2.00/3 \\
\code{p2--CCp3--p1} & triang. no estr. & 0.1 & 2.043 & 2.402 & 2.162 & 2.162 & 2.158 & 2.158 & 1.980 & 2.081 & 2.286 & 2.286 & 2.285 & 2.285 & 659.73 & 27285.2 & 0 & 2.00/3 \\
\hline
\code{p2--CCp3--p2} & hexa & 0.0 & 2.212 & 2.708 & 2.284 & 3.846 & 2.283 & 2.347 & 2.026 & 2.303 & 2.206 & 2.397 & 2.206 & 2.204 & 148.27 & 21478.1 & 0 & 2.00/3 \\
\code{p2--CCp3--p2} & hexa & 0.1 & 2.224 & 2.760 & 2.305 & 3.936 & 2.304 & 2.369 & 2.030 & 2.365 & 2.239 & 2.584 & 2.238 & 2.237 & 157.69 & 21927.2 & 0 & 2.00/3 \\
\code{p2--CCp3--p2} & triang. no estr. & 0.0 & 2.118 & 2.532 & 2.286 & 3.422 & 2.279 & 2.345 & 1.977 & 2.149 & 2.413 & 2.437 & 2.420 & 2.419 & 729.00 & 26704.4 & 0 & 2.00/3 \\
\code{p2--CCp3--p2} & triang. no estr. & 0.1 & 2.104 & 2.579 & 2.300 & 3.528 & 2.293 & 2.358 & 1.972 & 2.186 & 2.481 & 2.550 & 2.482 & 2.482 & 657.26 & 27605.0 & 0 & 2.00/3 \\
\hline
\code{p2--CCp3--p3} & hexa & 0.0 & 2.212 & 2.708 & 2.284 & 3.839 & 2.283 & 2.346 & 2.026 & 2.303 & 2.206 & 2.351 & 2.206 & 2.204 & 159.02 & 21614.3 & 0 & 2.00/3 \\
\code{p2--CCp3--p3} & hexa & 0.1 & 2.224 & 2.760 & 2.305 & 3.933 & 2.304 & 2.368 & 2.030 & 2.365 & 2.239 & 2.531 & 2.238 & 2.237 & 188.88 & 22066.7 & 0 & 2.00/3 \\
\code{p2--CCp3--p3} & triang. no estr. & 0.0 & 2.118 & 2.532 & 2.286 & 3.389 & 2.279 & 2.343 & 1.977 & 2.149 & 2.413 & 2.404 & 2.420 & 2.419 & 714.03 & 26849.9 & 0 & 2.00/3 \\
\code{p2--CCp3--p3} & triang. no estr. & 0.1 & 2.104 & 2.579 & 2.300 & 3.497 & 2.293 & 2.356 & 1.972 & 2.186 & 2.481 & 2.509 & 2.482 & 2.482 & 644.62 & 27732.2 & 0 & 2.00/3 \\
\hline
\code{p2--GP--p1} & hexa & 0.0 & 2.149 & 2.608 & 2.004 & 2.716 & 2.109 & 2.720 & 2.019 & 2.210 & 2.002 & 2.416 & 2.034 & 2.422 & 99.329 & 18816.7 & 0 & 2.00/3 \\
\code{p2--GP--p1} & hexa & 0.1 & 2.152 & 2.651 & 2.004 & 2.757 & 2.107 & 2.760 & 2.021 & 2.247 & 2.002 & 2.470 & 2.033 & 2.476 & 110.50 & 20350.3 & 0 & 2.00/3 \\
\code{p2--GP--p1} & triang. no estr. & 0.0 & 2.056 & 2.374 & 2.166 & 2.166 & 2.161 & 2.161 & 1.982 & 2.068 & 2.260 & 2.260 & 2.263 & 2.263 & 581.18 & 22050.9 & 0 & 2.00/3 \\
\code{p2--GP--p1} & triang. no estr. & 0.1 & 2.043 & 2.402 & 2.162 & 2.162 & 2.158 & 2.158 & 1.980 & 2.081 & 2.286 & 2.286 & 2.285 & 2.285 & 626.10 & 24445.3 & 0 & 2.00/3 \\
\hline
\code{p2--GP--p2} & hexa & 0.0 & 2.212 & 2.708 & 2.004 & 2.942 & 2.118 & 2.941 & 2.026 & 2.303 & 2.001 & 2.734 & 2.020 & 2.738 & 116.93 & 19113.0 & 0 & 2.00/3 \\
\code{p2--GP--p2} & hexa & 0.1 & 2.224 & 2.760 & 2.004 & 2.969 & 2.116 & 2.968 & 2.030 & 2.365 & 2.001 & 2.794 & 2.019 & 2.797 & 111.94 & 20456.4 & 0 & 2.00/3 \\
\code{p2--GP--p2} & triang. no estr. & 0.0 & 2.118 & 2.532 & 1.993 & 2.860 & 2.044 & 2.862 & 1.977 & 2.149 & 2.000 & 2.751 & 2.026 & 2.757 & 618.26 & 22383.5 & 0 & 2.00/3 \\
\code{p2--GP--p2} & triang. no estr. & 0.1 & 2.104 & 2.579 & 1.993 & 2.895 & 2.042 & 2.896 & 1.972 & 2.186 & 2.000 & 2.820 & 2.025 & 2.823 & 648.71 & 24583.9 & 0 & 2.00/3 \\
\hline
\code{p2--GP--p3} & hexa & 0.0 & 2.212 & 2.708 & 2.004 & 2.940 & 2.118 & 2.940 & 2.026 & 2.303 & 2.001 & 2.739 & 2.020 & 2.743 & 129.66 & 19251.7 & 0 & 2.00/3 \\
\code{p2--GP--p3} & hexa & 0.1 & 2.224 & 2.760 & 2.004 & 2.967 & 2.117 & 2.966 & 2.030 & 2.365 & 2.001 & 2.798 & 2.019 & 2.801 & 124.41 & 20530.1 & 0 & 2.00/3 \\
\code{p2--GP--p3} & triang. no estr. & 0.0 & 2.118 & 2.532 & 1.993 & 2.863 & 2.044 & 2.865 & 1.977 & 2.149 & 2.000 & 2.758 & 2.026 & 2.764 & 637.22 & 22551.2 & 0 & 2.00/3 \\
\code{p2--GP--p3} & triang. no estr. & 0.1 & 2.104 & 2.579 & 1.993 & 2.897 & 2.041 & 2.898 & 1.972 & 2.186 & 2.000 & 2.825 & 2.025 & 2.829 & 691.48 & 24650.1 & 0 & 2.00/3 \\
\hline
\code{p3--na--NO} & hexa & 0.0 & 2.121 & 2.481 & 2.256 & 2.256 & 2.258 & 2.258 & 1.999 & 2.037 & 2.472 & 2.472 & 2.469 & 2.469 & 147.82 & 21837.0 & 0 & 2.00/3 \\
\code{p3--na--NO} & hexa & 0.1 & 2.104 & 2.476 & 2.238 & 2.238 & 2.240 & 2.240 & 1.996 & 2.034 & 2.467 & 2.467 & 2.464 & 2.464 & 152.52 & 24499.4 & 0 & 2.00/3 \\
\code{p3--na--NO} & triang. no estr. & 0.0 & 2.033 & 2.201 & 2.348 & 2.348 & 2.345 & 2.345 & 1.972 & 1.988 & 3.183 & 3.183 & 3.172 & 3.172 & 689.38 & 27545.2 & 0 & 2.00/3 \\
\code{p3--na--NO} & triang. no estr. & 0.1 & 2.021 & 2.196 & 2.333 & 2.333 & 2.331 & 2.331 & 1.972 & 1.987 & 3.182 & 3.182 & 3.171 & 3.171 & 760.32 & 31535.9 & 0 & 2.00/3 \\
\hline
\code{p3--CCp1--p1} & hexa & 0.0 & 2.180 & 2.629 & 2.347 & 2.353 & 2.360 & 2.339 & 2.003 & 2.055 & 2.608 & 2.289 & 2.617 & 2.247 & 167.38 & 23001.8 & 0 & 2.00/3 \\
\code{p3--CCp1--p1} & hexa & 0.1 & 2.155 & 2.624 & 2.323 & 2.353 & 2.335 & 2.339 & 1.998 & 2.051 & 2.600 & 2.289 & 2.608 & 2.247 & 204.47 & 25294.3 & 0 & 2.00/3 \\
\code{p3--CCp1--p1} & triang. no estr. & 0.0 & 2.033 & 2.201 & 2.348 & 2.348 & 2.345 & 2.345 & 1.972 & 1.988 & 3.183 & 3.183 & 3.172 & 3.172 & 737.71 & 28449.0 & 0 & 2.00/3 \\
\code{p3--CCp1--p1} & triang. no estr. & 0.1 & 2.021 & 2.196 & 2.333 & 2.333 & 2.331 & 2.331 & 1.972 & 1.987 & 3.182 & 3.182 & 3.171 & 3.171 & 623.17 & 31985.4 & 0 & 2.00/3 \\
\hline
\code{p3--CCp1--p2} & hexa & 0.0 & 2.672 & 3.361 & 2.672 & 2.229 & 2.497 & 2.184 & 2.142 & 2.316 & 1.405 & 2.150 & 1.266 & 2.093 & 152.92 & 23136.5 & 0 & 2.00/3 \\
\code{p3--CCp1--p2} & hexa & 0.1 & 2.615 & 3.360 & 2.623 & 2.229 & 2.442 & 2.184 & 2.103 & 2.281 & 1.306 & 2.150 & 1.209 & 2.093 & 202.04 & 25404.4 & 0 & 2.00/3 \\
\code{p3--CCp1--p2} & triang. no estr. & 0.0 & 2.439 & 2.843 & 2.197 & 2.101 & 2.046 & 2.038 & 1.999 & 2.026 & 0.132 & 2.057 & 0.463 & 2.016 & 745.72 & 28688.4 & 0 & 2.00/3 \\
\code{p3--CCp1--p2} & triang. no estr. & 0.1 & 2.369 & 2.829 & 2.121 & 2.101 & 1.976 & 2.038 & 1.995 & 2.022 & 0.125 & 2.057 & 0.459 & 2.016 & 632.27 & 32128.5 & 0 & 2.00/3 \\
\hline
\code{p3--CCp1--p3} & hexa & 0.0 & 2.672 & 3.361 & 2.672 & 2.247 & 2.497 & 2.210 & 2.142 & 2.316 & 1.405 & 2.166 & 1.266 & 2.111 & 153.33 & 23218.0 & 0 & 2.00/3 \\
\code{p3--CCp1--p3} & hexa & 0.1 & 2.615 & 3.360 & 2.623 & 2.247 & 2.442 & 2.210 & 2.103 & 2.281 & 1.306 & 2.166 & 1.209 & 2.111 & 191.23 & 25439.1 & 0 & 2.00/3 \\
\code{p3--CCp1--p3} & triang. no estr. & 0.0 & 2.439 & 2.843 & 2.196 & 2.126 & 2.045 & 2.048 & 1.999 & 2.026 & 0.132 & 2.073 & 0.463 & 2.021 & 752.75 & 28791.6 & 0 & 2.00/3 \\
\code{p3--CCp1--p3} & triang. no estr. & 0.1 & 2.369 & 2.829 & 2.121 & 2.126 & 1.976 & 2.048 & 1.995 & 2.023 & 0.125 & 2.073 & 0.459 & 2.021 & 640.42 & 32187.4 & 0 & 2.00/3 \\
\hline
\code{p3--CCp2--p1} & hexa & 0.0 & 2.181 & 2.631 & 2.362 & 3.042 & 2.362 & 2.762 & 2.003 & 2.055 & 2.629 & 2.999 & 2.619 & 2.745 & 158.82 & 23118.4 & 0 & 2.00/3 \\
\code{p3--CCp2--p1} & hexa & 0.1 & 2.156 & 2.626 & 2.337 & 3.042 & 2.337 & 2.758 & 1.998 & 2.051 & 2.620 & 2.999 & 2.611 & 2.740 & 161.93 & 25259.8 & 0 & 2.00/3 \\
\code{p3--CCp2--p1} & triang. no estr. & 0.0 & 2.033 & 2.201 & 2.348 & 2.348 & 2.345 & 2.345 & 1.972 & 1.988 & 3.183 & 3.183 & 3.172 & 3.172 & 860.40 & 28760.1 & 0 & 2.00/3 \\
\code{p3--CCp2--p1} & triang. no estr. & 0.1 & 2.021 & 2.196 & 2.333 & 2.333 & 2.331 & 2.331 & 1.972 & 1.987 & 3.182 & 3.182 & 3.171 & 3.171 & 830.09 & 32013.2 & 0 & 2.00/3 \\
\hline
\code{p3--CCp2--p2} & hexa & 0.0 & 3.520 & 3.971 & 2.698 & 3.063 & 2.707 & 3.071 & 4.171 & 4.123 & 0.669 & 3.002 & 0.641 & 2.724 & 159.00 & 23274.3 & 0 & 2.00/3 \\
\code{p3--CCp2--p2} & hexa & 0.1 & 3.547 & 4.002 & 2.648 & 3.063 & 2.657 & 3.069 & 4.104 & 4.100 & 0.514 & 3.002 & 0.488 & 2.717 & 162.70 & 25376.0 & 0 & 2.00/3 \\
\code{p3--CCp2--p2} & triang. no estr. & 0.0 & 3.748 & 3.924 & 2.081 & 3.013 & 2.060 & 2.855 & 2.167 & 2.959 & 0.011 & 2.997 & 0.007 & 1.544 & 906.22 & 29070.9 & 0 & 2.00/3 \\
\code{p3--CCp2--p2} & triang. no estr. & 0.1 & 3.692 & 3.923 & 1.983 & 3.013 & 1.961 & 2.852 & 1.804 & 2.888 & 0.004 & 2.997 & 0.001 & 1.543 & 824.68 & 32175.8 & 0 & 2.00/3 \\
\hline
\code{p3--CCp2--p3} & hexa & 0.0 & 3.520 & 3.971 & 2.698 & 3.058 & 2.707 & 3.064 & 4.171 & 4.123 & 0.669 & 3.001 & 0.641 & 2.729 & 145.69 & 23372.0 & 0 & 2.00/3 \\
\code{p3--CCp2--p3} & hexa & 0.1 & 3.547 & 4.002 & 2.648 & 3.058 & 2.657 & 3.062 & 4.104 & 4.100 & 0.514 & 3.001 & 0.488 & 2.721 & 166.94 & 25416.4 & 0 & 2.00/3 \\
\code{p3--CCp2--p3} & triang. no estr. & 0.0 & 3.748 & 3.924 & 2.081 & 3.011 & 2.060 & 2.849 & 2.167 & 2.959 & 0.011 & 2.996 & 0.007 & 1.567 & 891.74 & 29233.2 & 0 & 2.00/3 \\
\code{p3--CCp2--p3} & triang. no estr. & 0.1 & 3.692 & 3.923 & 1.982 & 3.011 & 1.961 & 2.846 & 1.805 & 2.888 & 0.004 & 2.996 & 0.001 & 1.565 & 810.88 & 32261.1 & 0 & 2.00/3 \\
\hline
\code{p3--CCp3--p1} & hexa & 0.0 & 2.181 & 2.631 & 2.362 & 4.082 & 2.362 & 2.385 & 2.003 & 2.056 & 2.629 & 3.391 & 2.619 & 2.619 & 164.32 & 23744.0 & 0 & 2.00/3 \\
\code{p3--CCp3--p1} & hexa & 0.1 & 2.156 & 2.626 & 2.337 & 4.082 & 2.337 & 2.361 & 1.998 & 2.051 & 2.620 & 3.388 & 2.611 & 2.611 & 209.86 & 25311.2 & 0 & 2.00/3 \\
\code{p3--CCp3--p1} & triang. no estr. & 0.0 & 2.033 & 2.201 & 2.348 & 2.348 & 2.345 & 2.345 & 1.972 & 1.988 & 3.183 & 3.183 & 3.172 & 3.172 & 952.56 & 31101.1 & 0 & 2.00/3 \\
\code{p3--CCp3--p1} & triang. no estr. & 0.1 & 2.021 & 2.196 & 2.333 & 2.333 & 2.331 & 2.331 & 1.972 & 1.987 & 3.182 & 3.182 & 3.171 & 3.171 & 832.89 & 32230.4 & 0 & 2.00/3 \\
\hline
\code{p3--CCp3--p2} & hexa & 0.0 & 3.520 & 3.971 & 2.705 & 4.177 & 2.709 & 2.869 & 4.171 & 4.124 & 0.671 & 3.889 & 0.642 & 0.730 & 181.97 & 24009.8 & 0 & 2.00/3 \\
\code{p3--CCp3--p2} & hexa & 0.1 & 3.547 & 4.002 & 2.656 & 4.177 & 2.658 & 2.840 & 4.105 & 4.100 & 0.516 & 3.888 & 0.490 & 0.578 & 236.84 & 25434.9 & 0 & 2.00/3 \\
\code{p3--CCp3--p2} & triang. no estr. & 0.0 & 3.749 & 3.924 & 2.085 & 3.907 & 2.062 & 2.112 & 2.167 & 2.959 & 0.011 & 2.695 & 0.007 & 0.004 & 829.89 & 31443.8 & 0 & 2.00/3 \\
\code{p3--CCp3--p2} & triang. no estr. & 0.1 & 3.693 & 3.924 & 1.987 & 3.906 & 1.963 & 2.027 & 1.804 & 2.888 & 0.004 & 2.694 & 0.001 & -0.001 & 826.54 & 32405.2 & 0 & 2.00/3 \\
\hline
\code{p3--CCp3--p3} & hexa & 0.0 & 3.520 & 3.971 & 2.705 & 4.182 & 2.709 & 2.866 & 4.171 & 4.124 & 0.671 & 3.875 & 0.642 & 0.729 & 173.66 & 24132.2 & 0 & 2.00/3 \\
\code{p3--CCp3--p3} & hexa & 0.1 & 3.547 & 4.002 & 2.656 & 4.182 & 2.658 & 2.835 & 4.105 & 4.100 & 0.516 & 3.873 & 0.490 & 0.577 & 255.78 & 25497.1 & 0 & 2.00/3 \\
\code{p3--CCp3--p3} & triang. no estr. & 0.0 & 3.749 & 3.924 & 2.085 & 3.889 & 2.061 & 2.108 & 2.167 & 2.959 & 0.011 & 2.592 & 0.007 & 0.004 & 862.24 & 31610.3 & 0 & 2.00/3 \\
\code{p3--CCp3--p3} & triang. no estr. & 0.1 & 3.693 & 3.924 & 1.986 & 3.889 & 1.963 & 2.022 & 1.804 & 2.888 & 0.004 & 2.591 & 0.001 & -0.001 & 738.79 & 32495.2 & 0 & 2.00/3 \\
\hline
\code{p3--GP--p1} & hexa & 0.0 & 2.181 & 2.631 & 2.002 & 2.762 & 1.992 & 2.757 & 2.003 & 2.056 & 2.002 & 2.661 & 2.008 & 2.650 & 148.12 & 21874.5 & 0 & 2.00/3 \\
\code{p3--GP--p1} & hexa & 0.1 & 2.156 & 2.626 & 2.002 & 2.756 & 1.991 & 2.752 & 1.998 & 2.051 & 2.002 & 2.653 & 2.008 & 2.642 & 178.17 & 24556.5 & 0 & 2.00/3 \\
\code{p3--GP--p1} & triang. no estr. & 0.0 & 2.033 & 2.201 & 2.348 & 2.348 & 2.345 & 2.345 & 1.972 & 1.988 & 3.183 & 3.183 & 3.172 & 3.172 & 693.51 & 27581.7 & 0 & 2.00/3 \\
\code{p3--GP--p1} & triang. no estr. & 0.1 & 2.021 & 2.196 & 2.333 & 2.333 & 2.331 & 2.331 & 1.972 & 1.987 & 3.182 & 3.182 & 3.171 & 3.171 & 800.77 & 31560.1 & 0 & 2.00/3 \\
\hline
\code{p3--GP--p2} & hexa & 0.0 & 3.520 & 3.971 & 2.000 & 3.744 & 1.999 & 3.748 & 4.171 & 4.124 & 2.001 & 1.458 & 2.006 & 1.396 & 170.50 & 21996.1 & 0 & 2.00/3 \\
\code{p3--GP--p2} & hexa & 0.1 & 3.547 & 4.002 & 1.999 & 3.757 & 1.998 & 3.761 & 4.105 & 4.100 & 2.001 & 1.348 & 2.006 & 1.284 & 202.72 & 24666.6 & 0 & 2.00/3 \\
\code{p3--GP--p2} & triang. no estr. & 0.0 & 3.749 & 3.924 & 1.992 & 2.844 & 2.005 & 2.843 & 2.167 & 2.959 & 2.000 & -0.092 & 2.012 & -0.097 & 805.00 & 27710.0 & 0 & 2.00/3 \\
\code{p3--GP--p2} & triang. no estr. & 0.1 & 3.693 & 3.924 & 1.992 & 2.836 & 2.005 & 2.836 & 1.804 & 2.888 & 2.000 & -0.098 & 2.012 & -0.103 & 699.63 & 31685.8 & 0 & 2.00/3 \\
\hline
\code{p3--GP--p3} & hexa & 0.0 & 3.520 & 3.971 & 2.000 & 3.738 & 1.999 & 3.742 & 4.171 & 4.124 & 2.001 & 1.420 & 2.006 & 1.355 & 174.93 & 22054.6 & 0 & 2.00/3 \\
\code{p3--GP--p3} & hexa & 0.1 & 3.547 & 4.002 & 2.000 & 3.751 & 1.999 & 3.756 & 4.105 & 4.100 & 2.001 & 1.305 & 2.006 & 1.239 & 246.90 & 24749.6 & 0 & 2.00/3 \\
\code{p3--GP--p3} & triang. no estr. & 0.0 & 3.749 & 3.924 & 1.992 & 2.812 & 2.005 & 2.811 & 2.167 & 2.959 & 2.000 & -0.108 & 2.012 & -0.114 & 882.90 & 27779.4 & 0 & 2.00/3 \\
\code{p3--GP--p3} & triang. no estr. & 0.1 & 3.693 & 3.923 & 1.992 & 2.803 & 2.005 & 2.803 & 1.804 & 2.888 & 2.000 & -0.115 & 2.012 & -0.120 & 886.85 & 31746.3 & 0 & 2.00/3 \\
\bottomrule
\end{longtable}
\normalsize

\paragraph{$\Delta t=10^{-3}$}
\tiny
\setlength{\tabcolsep}{1.4pt}
\begin{longtable}{lllrrrrrrrrrrrrrrrr}
\caption{Resultados manufacturados 2D para \code{JFNK}, \code{crankNicolson}, masa \code{consistent}, estados \code{RKF45} y $\Delta t=10^{-3}$. El primer bloque de convergencia se ajusta sobre $N=10$--$100$; el segundo sobre $N=70$--$100$. La primera columna registra la terna \code{Vm-states-Iion} efectivamente usada en la corrida. \textbf{RB} es la cantidad de pasos de tiempo que agotaron el presupuesto de iteraciones no lineales, acumulada sobre los 91 tamaños de malla de la fila (4550 pasos de tiempo en total por fila). Con \code{nonlinearAcceptUnconverged=false} un paso no convergido se revierte a los valores del inicio del paso mientras el reloj avanza igual. \textbf{$n_{\mathrm{nl}}$} reporta la mediana y el máximo de iteraciones no lineales por paso sobre esos mismos pasos; el presupuesto configurado es \code{nonlinearIterations}$=2000$.}\label{tab:es-results-2d-jfnk-1p0em03}\\
\toprule
& & & \multicolumn{6}{c}{Ajuste $N=10$--$100$} & \multicolumn{6}{c}{Ajuste $N=70$--$100$} & & & & \\
\cmidrule(lr){4-9}\cmidrule(lr){10-15}
\code{Vm-states-Iion} & Tipo de malla & $\alpha$ & $V_m$ & $V_m^G$ & $u_1$ & $u_1^G$ & $u_2$ & $u_2^G$ & $V_m$ & $V_m^G$ & $u_1$ & $u_1^G$ & $u_2$ & $u_2^G$ & $t_{\Sigma}$ [s] & RSS$_{\Sigma}$ [MB] & RB & $n_{\mathrm{nl}}$ \\
\midrule
\endfirsthead
\toprule
& & & \multicolumn{6}{c}{Ajuste $N=10$--$100$} & \multicolumn{6}{c}{Ajuste $N=70$--$100$} & & & & \\
\cmidrule(lr){4-9}\cmidrule(lr){10-15}
\code{Vm-states-Iion} & Tipo de malla & $\alpha$ & $V_m$ & $V_m^G$ & $u_1$ & $u_1^G$ & $u_2$ & $u_2^G$ & $V_m$ & $V_m^G$ & $u_1$ & $u_1^G$ & $u_2$ & $u_2^G$ & $t_{\Sigma}$ [s] & RSS$_{\Sigma}$ [MB] & RB & $n_{\mathrm{nl}}$ \\
\midrule
\endhead
\code{NO--na--NO} & hexa & 0.0 & 1.997 & 1.997 & 1.998 & 1.998 & 1.998 & 1.998 & 2.000 & 2.000 & 2.004 & 2.004 & 2.004 & 2.004 & 62.774 & 16499.9 & 0 & 2.00/3 \\
\code{NO--na--NO} & hexa & 0.1 & 1.851 & 1.851 & 1.607 & 1.607 & 1.613 & 1.613 & 2.294 & 2.294 & 2.140 & 2.140 & 2.142 & 2.142 & 104.37 & 18378.4 & 0 & 2.00/2 \\
\code{NO--na--NO} & triang. no estr. & 0.0 & 0.798 & 0.798 & 0.777 & 0.777 & 0.783 & 0.783 & 0.963 & 0.963 & 1.038 & 1.038 & 1.107 & 1.107 & 229.82 & 17289.7 & 0 & 2.00/3 \\
\code{NO--na--NO} & triang. no estr. & 0.1 & 0.785 & 0.785 & 0.771 & 0.771 & 0.776 & 0.776 & 0.953 & 0.953 & 1.027 & 1.027 & 1.097 & 1.097 & 335.59 & 20799.8 & 0 & 2.00/3 \\
\hline
\code{NO--CCp1--p1} & hexa & 0.0 & 1.997 & 1.997 & 1.999 & 2.354 & 1.998 & 2.333 & 2.000 & 2.000 & 2.004 & 2.290 & 2.004 & 2.242 & 99.950 & 17020.5 & 0 & 2.00/3 \\
\code{NO--CCp1--p1} & hexa & 0.1 & 1.851 & 1.851 & 1.608 & 2.354 & 1.613 & 2.310 & 2.294 & 2.294 & 2.140 & 2.290 & 2.142 & 2.226 & 139.82 & 18502.9 & 0 & 2.00/2 \\
\code{NO--CCp1--p1} & triang. no estr. & 0.0 & 0.798 & 0.798 & 0.777 & 0.777 & 0.783 & 0.783 & 0.963 & 0.963 & 1.038 & 1.038 & 1.107 & 1.107 & 463.94 & 18085.7 & 0 & 2.00/3 \\
\code{NO--CCp1--p1} & triang. no estr. & 0.1 & 0.785 & 0.785 & 0.771 & 0.771 & 0.776 & 0.776 & 0.953 & 0.953 & 1.027 & 1.027 & 1.097 & 1.097 & 503.50 & 20923.6 & 0 & 2.00/3 \\
\hline
\code{NO--CCp1--p2} & hexa & 0.0 & 1.997 & 1.997 & 2.000 & 2.229 & 1.999 & 2.184 & 2.000 & 2.000 & 2.004 & 2.150 & 2.004 & 2.094 & 130.30 & 17331.8 & 0 & 2.00/3 \\
\code{NO--CCp1--p2} & hexa & 0.1 & 1.851 & 1.851 & 1.608 & 2.229 & 1.614 & 2.181 & 2.295 & 2.295 & 2.140 & 2.150 & 2.142 & 2.092 & 154.64 & 18665.6 & 0 & 2.00/2 \\
\code{NO--CCp1--p2} & triang. no estr. & 0.0 & 0.798 & 0.798 & 0.777 & 2.111 & 0.783 & 1.163 & 0.963 & 0.963 & 1.038 & 2.102 & 1.107 & 1.150 & 438.67 & 18426.9 & 0 & 2.00/3 \\
\code{NO--CCp1--p2} & triang. no estr. & 0.1 & 0.785 & 0.785 & 0.771 & 2.112 & 0.776 & 1.171 & 0.953 & 0.953 & 1.027 & 2.105 & 1.097 & 1.143 & 521.01 & 21072.2 & 0 & 2.00/3 \\
\hline
\code{NO--CCp1--p3} & hexa & 0.0 & 1.997 & 1.997 & 2.000 & 2.247 & 1.999 & 2.211 & 2.000 & 2.000 & 2.004 & 2.167 & 2.004 & 2.113 & 132.84 & 17490.5 & 0 & 2.00/3 \\
\code{NO--CCp1--p3} & hexa & 0.1 & 1.851 & 1.851 & 1.608 & 2.247 & 1.614 & 2.206 & 2.295 & 2.295 & 2.140 & 2.167 & 2.142 & 2.110 & 151.61 & 18741.9 & 0 & 2.00/2 \\
\code{NO--CCp1--p3} & triang. no estr. & 0.0 & 0.798 & 0.798 & 0.777 & 2.139 & 0.783 & 1.131 & 0.963 & 0.963 & 1.038 & 2.131 & 1.107 & 1.144 & 439.70 & 18599.6 & 0 & 2.00/3 \\
\code{NO--CCp1--p3} & triang. no estr. & 0.1 & 0.785 & 0.785 & 0.771 & 2.140 & 0.776 & 1.138 & 0.953 & 0.953 & 1.027 & 2.135 & 1.097 & 1.137 & 524.48 & 21149.1 & 0 & 2.00/3 \\
\hline
\code{NO--CCp2--p1} & hexa & 0.0 & 1.997 & 1.997 & 2.000 & 3.042 & 1.997 & 2.510 & 1.999 & 1.999 & 2.004 & 2.998 & 2.003 & 2.157 & 167.05 & 17530.6 & 0 & 2.00/3 \\
\code{NO--CCp2--p1} & hexa & 0.1 & 1.851 & 1.851 & 1.603 & 3.038 & 1.613 & 1.942 & 2.294 & 2.294 & 2.140 & 2.987 & 2.142 & 2.167 & 200.77 & 18748.5 & 0 & 2.00/2 \\
\code{NO--CCp2--p1} & triang. no estr. & 0.0 & 0.798 & 0.798 & 0.777 & 0.777 & 0.783 & 0.783 & 0.963 & 0.963 & 1.038 & 1.038 & 1.107 & 1.107 & 729.76 & 19279.0 & 0 & 2.00/3 \\
\code{NO--CCp2--p1} & triang. no estr. & 0.1 & 0.785 & 0.785 & 0.771 & 0.771 & 0.776 & 0.776 & 0.953 & 0.953 & 1.027 & 1.027 & 1.097 & 1.097 & 767.38 & 21262.3 & 0 & 2.00/3 \\
\hline
\code{NO--CCp2--p2} & hexa & 0.0 & 1.995 & 1.995 & 2.003 & 3.063 & 1.997 & 2.671 & 1.999 & 1.999 & 2.004 & 3.002 & 2.003 & 2.246 & 196.75 & 17862.5 & 0 & 2.00/3 \\
\code{NO--CCp2--p2} & hexa & 0.1 & 1.853 & 1.853 & 1.597 & 3.061 & 1.619 & 2.063 & 2.289 & 2.289 & 2.140 & 2.996 & 2.141 & 2.192 & 217.72 & 19033.8 & 0 & 2.00/2 \\
\code{NO--CCp2--p2} & triang. no estr. & 0.0 & 0.799 & 0.799 & 0.777 & 1.940 & 0.783 & 0.775 & 0.963 & 0.963 & 1.038 & 1.095 & 1.107 & 1.097 & 840.38 & 19631.4 & 0 & 2.00/3 \\
\code{NO--CCp2--p2} & triang. no estr. & 0.1 & 0.786 & 0.786 & 0.771 & 1.955 & 0.776 & 0.772 & 0.953 & 0.953 & 1.027 & 1.090 & 1.097 & 1.087 & 707.36 & 21596.1 & 0 & 2.00/3 \\
\hline
\code{NO--CCp2--p3} & hexa & 0.0 & 1.995 & 1.995 & 2.003 & 3.058 & 1.997 & 2.670 & 1.999 & 1.999 & 2.004 & 3.001 & 2.003 & 2.251 & 224.42 & 17998.3 & 0 & 2.00/3 \\
\code{NO--CCp2--p3} & hexa & 0.1 & 1.853 & 1.853 & 1.597 & 3.056 & 1.619 & 2.064 & 2.289 & 2.289 & 2.140 & 2.995 & 2.141 & 2.193 & 226.37 & 19155.7 & 0 & 2.00/2 \\
\code{NO--CCp2--p3} & triang. no estr. & 0.0 & 0.799 & 0.799 & 0.777 & 1.950 & 0.783 & 0.775 & 0.963 & 0.963 & 1.038 & 1.098 & 1.107 & 1.097 & 743.74 & 19802.3 & 0 & 2.00/3 \\
\code{NO--CCp2--p3} & triang. no estr. & 0.1 & 0.786 & 0.786 & 0.771 & 1.966 & 0.776 & 0.772 & 0.953 & 0.953 & 1.027 & 1.092 & 1.097 & 1.087 & 808.05 & 21765.1 & 0 & 2.00/3 \\
\hline
\code{NO--CCp3--p1} & hexa & 0.0 & 1.997 & 1.997 & 1.998 & 3.944 & 1.997 & 2.030 & 1.999 & 1.999 & 2.004 & 2.505 & 2.004 & 2.004 & 362.53 & 18892.7 & 0 & 2.00/3 \\
\code{NO--CCp3--p1} & hexa & 0.1 & 1.851 & 1.851 & 1.600 & 3.114 & 1.613 & 1.657 & 2.294 & 2.294 & 2.140 & 2.324 & 2.142 & 2.170 & 356.97 & 20106.8 & 0 & 2.00/2 \\
\code{NO--CCp3--p1} & triang. no estr. & 0.0 & 0.798 & 0.798 & 0.777 & 0.777 & 0.783 & 0.783 & 0.963 & 0.963 & 1.038 & 1.038 & 1.107 & 1.107 & 1284.8 & 21788.5 & 0 & 2.00/3 \\
\code{NO--CCp3--p1} & triang. no estr. & 0.1 & 0.785 & 0.785 & 0.771 & 0.771 & 0.776 & 0.776 & 0.953 & 0.953 & 1.027 & 1.027 & 1.097 & 1.097 & 1109.8 & 23745.7 & 0 & 2.00/3 \\
\hline
\code{NO--CCp3--p2} & hexa & 0.0 & 1.996 & 1.996 & 1.998 & 4.244 & 1.996 & 2.092 & 1.999 & 1.999 & 2.004 & 3.862 & 2.003 & 2.002 & 388.72 & 19203.1 & 0 & 2.00/3 \\
\code{NO--CCp3--p2} & hexa & 0.1 & 1.853 & 1.853 & 1.590 & 3.420 & 1.617 & 1.718 & 2.289 & 2.289 & 2.140 & 2.581 & 2.141 & 2.203 & 371.88 & 20400.8 & 0 & 2.00/2 \\
\code{NO--CCp3--p2} & triang. no estr. & 0.0 & 0.799 & 0.799 & 0.777 & 1.070 & 0.783 & 0.795 & 0.963 & 0.963 & 1.038 & 1.038 & 1.107 & 1.110 & 1354.0 & 22118.9 & 0 & 2.00/3 \\
\code{NO--CCp3--p2} & triang. no estr. & 0.1 & 0.786 & 0.786 & 0.771 & 1.080 & 0.776 & 0.797 & 0.953 & 0.953 & 1.027 & 1.028 & 1.097 & 1.100 & 1096.4 & 24075.7 & 0 & 2.00/3 \\
\hline
\code{NO--CCp3--p3} & hexa & 0.0 & 1.996 & 1.996 & 1.998 & 4.257 & 1.996 & 2.090 & 1.999 & 1.999 & 2.004 & 3.849 & 2.003 & 2.002 & 412.25 & 19356.7 & 0 & 2.00/3 \\
\code{NO--CCp3--p3} & hexa & 0.1 & 1.853 & 1.853 & 1.590 & 3.408 & 1.617 & 1.717 & 2.289 & 2.289 & 2.140 & 2.570 & 2.141 & 2.203 & 375.60 & 20540.9 & 0 & 2.00/2 \\
\code{NO--CCp3--p3} & triang. no estr. & 0.0 & 0.799 & 0.799 & 0.777 & 1.053 & 0.783 & 0.795 & 0.963 & 0.963 & 1.038 & 1.038 & 1.107 & 1.110 & 1357.6 & 22304.1 & 0 & 2.00/3 \\
\code{NO--CCp3--p3} & triang. no estr. & 0.1 & 0.786 & 0.786 & 0.771 & 1.062 & 0.776 & 0.796 & 0.953 & 0.953 & 1.027 & 1.028 & 1.097 & 1.100 & 1205.0 & 24248.2 & 0 & 2.00/3 \\
\hline
\code{NO--GP--p1} & hexa & 0.0 & 1.997 & 1.997 & 1.999 & 2.050 & 1.998 & 2.053 & 2.000 & 2.000 & 2.000 & 2.007 & 2.000 & 2.007 & 87.918 & 16718.4 & 0 & 2.00/3 \\
\code{NO--GP--p1} & hexa & 0.1 & 1.851 & 1.851 & 1.999 & 2.042 & 1.990 & 2.046 & 2.294 & 2.294 & 2.000 & 2.015 & 2.008 & 2.016 & 112.16 & 18447.0 & 0 & 2.00/2 \\
\code{NO--GP--p1} & triang. no estr. & 0.0 & 0.798 & 0.798 & 0.777 & 0.777 & 0.783 & 0.783 & 0.963 & 0.963 & 1.038 & 1.038 & 1.107 & 1.107 & 299.64 & 17478.6 & 0 & 2.00/3 \\
\code{NO--GP--p1} & triang. no estr. & 0.1 & 0.785 & 0.785 & 0.771 & 0.771 & 0.776 & 0.776 & 0.953 & 0.953 & 1.027 & 1.027 & 1.097 & 1.097 & 455.19 & 20904.7 & 0 & 2.00/3 \\
\hline
\code{NO--GP--p2} & hexa & 0.0 & 1.997 & 1.997 & 1.999 & 2.013 & 1.998 & 2.015 & 1.999 & 1.999 & 2.000 & 2.002 & 2.000 & 2.002 & 219.26 & 17038.3 & 0 & 2.00/3 \\
\code{NO--GP--p2} & hexa & 0.1 & 1.853 & 1.853 & 1.999 & 2.013 & 1.998 & 2.014 & 2.290 & 2.290 & 2.000 & 2.003 & 2.001 & 2.003 & 209.11 & 18552.2 & 0 & 2.00/2 \\
\code{NO--GP--p2} & triang. no estr. & 0.0 & 0.799 & 0.799 & 1.878 & 1.417 & 1.101 & 1.410 & 0.963 & 0.963 & 1.677 & 1.218 & 1.112 & 1.269 & 541.47 & 17817.7 & 0 & 2.00/3 \\
\code{NO--GP--p2} & triang. no estr. & 0.1 & 0.786 & 0.786 & 1.880 & 1.426 & 1.105 & 1.419 & 0.953 & 0.953 & 1.680 & 1.219 & 1.104 & 1.270 & 539.30 & 21041.9 & 0 & 2.00/3 \\
\hline
\code{NO--GP--p3} & hexa & 0.0 & 1.997 & 1.997 & 1.999 & 2.016 & 1.999 & 2.018 & 1.999 & 1.999 & 2.000 & 2.002 & 2.000 & 2.002 & 253.15 & 17191.9 & 0 & 2.00/3 \\
\code{NO--GP--p3} & hexa & 0.1 & 1.853 & 1.853 & 1.999 & 2.016 & 1.998 & 2.018 & 2.290 & 2.290 & 2.000 & 2.004 & 2.002 & 2.004 & 265.83 & 18636.6 & 0 & 2.00/2 \\
\code{NO--GP--p3} & triang. no estr. & 0.0 & 0.799 & 0.799 & 1.878 & 1.388 & 1.101 & 1.383 & 0.963 & 0.963 & 1.677 & 1.205 & 1.112 & 1.258 & 777.03 & 18010.7 & 0 & 2.00/3 \\
\code{NO--GP--p3} & triang. no estr. & 0.1 & 0.786 & 0.786 & 1.880 & 1.397 & 1.106 & 1.391 & 0.953 & 0.953 & 1.680 & 1.206 & 1.104 & 1.258 & 770.03 & 21110.3 & 0 & 2.00/3 \\
\hline
\code{p1--na--NO} & hexa & 0.0 & 1.912 & 2.147 & 1.616 & 1.616 & 1.626 & 1.626 & 2.611 & 2.168 & 2.380 & 2.380 & 2.385 & 2.385 & 150.02 & 17887.6 & 0 & 2.00/3 \\
\code{p1--na--NO} & hexa & 0.1 & 1.971 & 2.163 & 1.677 & 1.677 & 1.687 & 1.687 & 2.615 & 2.141 & 2.414 & 2.414 & 2.419 & 2.419 & 175.68 & 19071.9 & 0 & 2.00/2 \\
\code{p1--na--NO} & triang. no estr. & 0.0 & 2.292 & 2.131 & 2.004 & 2.004 & 2.003 & 2.003 & 2.600 & 2.164 & 2.545 & 2.545 & 2.542 & 2.542 & 839.15 & 19894.4 & 0 & 2.00/3 \\
\code{p1--na--NO} & triang. no estr. & 0.1 & 2.345 & 2.112 & 2.092 & 2.092 & 2.089 & 2.089 & 2.552 & 2.095 & 2.516 & 2.516 & 2.515 & 2.515 & 966.93 & 21562.5 & 0 & 2.00/3 \\
\hline
\code{p1--CCp1--p1} & hexa & 0.0 & 1.912 & 2.147 & 1.616 & 2.350 & 1.626 & 1.955 & 2.611 & 2.168 & 2.380 & 2.295 & 2.385 & 2.369 & 210.62 & 18392.6 & 0 & 2.00/3 \\
\code{p1--CCp1--p1} & hexa & 0.1 & 1.971 & 2.163 & 1.677 & 2.351 & 1.687 & 2.045 & 2.615 & 2.141 & 2.414 & 2.294 & 2.419 & 2.382 & 241.87 & 19302.2 & 0 & 2.00/2 \\
\code{p1--CCp1--p1} & triang. no estr. & 0.0 & 2.292 & 2.131 & 2.004 & 2.004 & 2.003 & 2.003 & 2.600 & 2.164 & 2.545 & 2.545 & 2.542 & 2.542 & 886.52 & 20537.9 & 0 & 2.00/3 \\
\code{p1--CCp1--p1} & triang. no estr. & 0.1 & 2.345 & 2.112 & 2.092 & 2.092 & 2.089 & 2.089 & 2.552 & 2.095 & 2.516 & 2.516 & 2.515 & 2.515 & 1042.2 & 21806.4 & 0 & 2.00/3 \\
\hline
\code{p1--CCp1--p2} & hexa & 0.0 & 1.912 & 2.147 & 1.616 & 2.228 & 1.626 & 2.057 & 2.611 & 2.168 & 2.380 & 2.154 & 2.385 & 2.233 & 214.61 & 18677.9 & 0 & 2.00/3 \\
\code{p1--CCp1--p2} & hexa & 0.1 & 1.971 & 2.163 & 1.676 & 2.228 & 1.686 & 2.104 & 2.615 & 2.141 & 2.414 & 2.153 & 2.419 & 2.209 & 253.57 & 19440.2 & 0 & 2.00/2 \\
\code{p1--CCp1--p2} & triang. no estr. & 0.0 & 2.292 & 2.131 & 2.004 & 2.103 & 2.003 & 2.036 & 2.600 & 2.164 & 2.545 & 2.063 & 2.542 & 2.283 & 860.24 & 20876.2 & 0 & 2.00/3 \\
\code{p1--CCp1--p2} & triang. no estr. & 0.1 & 2.345 & 2.112 & 2.092 & 2.103 & 2.089 & 2.066 & 2.552 & 2.095 & 2.516 & 2.061 & 2.515 & 2.188 & 998.30 & 22029.6 & 0 & 2.00/3 \\
\hline
\code{p1--CCp1--p3} & hexa & 0.0 & 1.912 & 2.147 & 1.616 & 2.246 & 1.626 & 2.063 & 2.611 & 2.168 & 2.380 & 2.171 & 2.385 & 2.256 & 209.58 & 18812.9 & 0 & 2.00/3 \\
\code{p1--CCp1--p3} & hexa & 0.1 & 1.971 & 2.163 & 1.676 & 2.246 & 1.686 & 2.115 & 2.615 & 2.141 & 2.414 & 2.170 & 2.419 & 2.234 & 237.29 & 19520.2 & 0 & 2.00/2 \\
\code{p1--CCp1--p3} & triang. no estr. & 0.0 & 2.292 & 2.131 & 2.004 & 2.127 & 2.003 & 2.040 & 2.600 & 2.164 & 2.545 & 2.081 & 2.542 & 2.312 & 847.53 & 21045.5 & 0 & 2.00/3 \\
\code{p1--CCp1--p3} & triang. no estr. & 0.1 & 2.345 & 2.112 & 2.092 & 2.128 & 2.089 & 2.076 & 2.552 & 2.095 & 2.516 & 2.078 & 2.515 & 2.216 & 1030.2 & 22141.1 & 0 & 2.00/3 \\
\hline
\code{p1--CCp2--p1} & hexa & 0.0 & 1.912 & 2.147 & 1.616 & 2.903 & 1.626 & 1.665 & 2.611 & 2.168 & 2.380 & 2.782 & 2.385 & 2.394 & 238.64 & 18898.4 & 0 & 2.00/3 \\
\code{p1--CCp2--p1} & hexa & 0.1 & 1.971 & 2.163 & 1.676 & 2.953 & 1.687 & 1.736 & 2.615 & 2.141 & 2.414 & 2.858 & 2.419 & 2.428 & 260.96 & 19460.9 & 0 & 2.00/2 \\
\code{p1--CCp2--p1} & triang. no estr. & 0.0 & 2.292 & 2.131 & 2.004 & 2.004 & 2.003 & 2.003 & 2.600 & 2.164 & 2.545 & 2.545 & 2.542 & 2.542 & 1038.1 & 21734.5 & 0 & 2.00/3 \\
\code{p1--CCp2--p1} & triang. no estr. & 0.1 & 2.345 & 2.112 & 2.092 & 2.092 & 2.089 & 2.089 & 2.552 & 2.095 & 2.516 & 2.516 & 2.515 & 2.515 & 1177.4 & 22368.2 & 0 & 2.00/3 \\
\hline
\code{p1--CCp2--p2} & hexa & 0.0 & 1.912 & 2.147 & 1.616 & 2.990 & 1.627 & 1.689 & 2.611 & 2.168 & 2.380 & 2.882 & 2.385 & 2.400 & 247.89 & 19179.4 & 0 & 2.00/3 \\
\code{p1--CCp2--p2} & hexa & 0.1 & 1.972 & 2.163 & 1.676 & 3.018 & 1.687 & 1.765 & 2.615 & 2.141 & 2.414 & 2.931 & 2.419 & 2.435 & 276.80 & 19673.1 & 0 & 2.00/2 \\
\code{p1--CCp2--p2} & triang. no estr. & 0.0 & 2.292 & 2.131 & 2.004 & 2.676 & 2.003 & 1.978 & 2.600 & 2.164 & 2.546 & 2.665 & 2.542 & 2.522 & 1028.7 & 22067.9 & 0 & 2.00/3 \\
\code{p1--CCp2--p2} & triang. no estr. & 0.1 & 2.345 & 2.112 & 2.092 & 2.799 & 2.089 & 2.063 & 2.552 & 2.095 & 2.516 & 2.731 & 2.515 & 2.504 & 1202.9 & 22669.4 & 0 & 2.00/3 \\
\hline
\code{p1--CCp2--p3} & hexa & 0.0 & 1.912 & 2.147 & 1.616 & 2.987 & 1.627 & 1.688 & 2.611 & 2.168 & 2.380 & 2.884 & 2.385 & 2.400 & 262.69 & 19310.6 & 0 & 2.00/3 \\
\code{p1--CCp2--p3} & hexa & 0.1 & 1.972 & 2.163 & 1.676 & 3.015 & 1.687 & 1.765 & 2.615 & 2.141 & 2.414 & 2.931 & 2.419 & 2.435 & 293.94 & 19820.1 & 0 & 2.00/2 \\
\code{p1--CCp2--p3} & triang. no estr. & 0.0 & 2.292 & 2.131 & 2.004 & 2.681 & 2.003 & 1.978 & 2.600 & 2.164 & 2.546 & 2.668 & 2.542 & 2.522 & 1032.6 & 22243.1 & 0 & 2.00/3 \\
\code{p1--CCp2--p3} & triang. no estr. & 0.1 & 2.345 & 2.112 & 2.092 & 2.802 & 2.089 & 2.062 & 2.552 & 2.095 & 2.516 & 2.736 & 2.515 & 2.504 & 1219.5 & 22819.4 & 0 & 2.00/3 \\
\hline
\code{p1--CCp3--p1} & hexa & 0.0 & 1.912 & 2.147 & 1.617 & 2.034 & 1.626 & 1.630 & 2.611 & 2.168 & 2.380 & 2.378 & 2.385 & 2.406 & 388.93 & 20236.1 & 0 & 2.00/3 \\
\code{p1--CCp3--p1} & hexa & 0.1 & 1.971 & 2.163 & 1.677 & 2.164 & 1.687 & 1.692 & 2.615 & 2.141 & 2.414 & 2.408 & 2.419 & 2.441 & 405.12 & 20745.1 & 0 & 2.00/2 \\
\code{p1--CCp3--p1} & triang. no estr. & 0.0 & 2.292 & 2.131 & 2.004 & 2.004 & 2.003 & 2.003 & 2.600 & 2.164 & 2.545 & 2.545 & 2.542 & 2.542 & 1377.6 & 24229.1 & 0 & 2.00/3 \\
\code{p1--CCp3--p1} & triang. no estr. & 0.1 & 2.345 & 2.112 & 2.092 & 2.092 & 2.089 & 2.089 & 2.552 & 2.095 & 2.516 & 2.516 & 2.515 & 2.515 & 1489.4 & 24696.5 & 0 & 2.00/3 \\
\hline
\code{p1--CCp3--p2} & hexa & 0.0 & 1.913 & 2.147 & 1.617 & 2.239 & 1.627 & 1.635 & 2.611 & 2.168 & 2.380 & 2.377 & 2.385 & 2.426 & 398.56 & 20515.3 & 0 & 2.00/3 \\
\code{p1--CCp3--p2} & hexa & 0.1 & 1.972 & 2.163 & 1.677 & 2.394 & 1.688 & 1.698 & 2.615 & 2.141 & 2.414 & 2.402 & 2.419 & 2.463 & 428.87 & 21005.9 & 0 & 2.00/2 \\
\code{p1--CCp3--p2} & triang. no estr. & 0.0 & 2.292 & 2.131 & 2.004 & 2.031 & 2.003 & 1.965 & 2.600 & 2.164 & 2.546 & 2.521 & 2.542 & 2.523 & 1494.7 & 24554.1 & 0 & 2.00/3 \\
\code{p1--CCp3--p2} & triang. no estr. & 0.1 & 2.345 & 2.112 & 2.092 & 2.141 & 2.089 & 2.052 & 2.552 & 2.095 & 2.516 & 2.506 & 2.515 & 2.510 & 1457.9 & 25036.9 & 0 & 2.00/3 \\
\hline
\code{p1--CCp3--p3} & hexa & 0.0 & 1.913 & 2.147 & 1.617 & 2.224 & 1.627 & 1.634 & 2.611 & 2.168 & 2.380 & 2.376 & 2.385 & 2.426 & 407.19 & 20617.6 & 0 & 2.00/3 \\
\code{p1--CCp3--p3} & hexa & 0.1 & 1.972 & 2.163 & 1.677 & 2.379 & 1.688 & 1.698 & 2.615 & 2.141 & 2.414 & 2.401 & 2.419 & 2.462 & 430.20 & 21137.6 & 0 & 2.00/2 \\
\code{p1--CCp3--p3} & triang. no estr. & 0.0 & 2.292 & 2.131 & 2.004 & 2.028 & 2.003 & 1.965 & 2.600 & 2.164 & 2.546 & 2.522 & 2.542 & 2.523 & 1516.1 & 24723.6 & 0 & 2.00/3 \\
\code{p1--CCp3--p3} & triang. no estr. & 0.1 & 2.345 & 2.112 & 2.092 & 2.136 & 2.089 & 2.052 & 2.552 & 2.095 & 2.516 & 2.506 & 2.515 & 2.510 & 1452.3 & 25196.1 & 0 & 2.00/3 \\
\hline
\code{p1--GP--p1} & hexa & 0.0 & 1.912 & 2.147 & 1.997 & 2.152 & 1.836 & 2.158 & 2.611 & 2.168 & 2.005 & 2.309 & 2.249 & 2.308 & 183.83 & 17987.7 & 0 & 2.00/3 \\
\code{p1--GP--p1} & hexa & 0.1 & 1.971 & 2.163 & 1.999 & 2.219 & 1.895 & 2.225 & 2.615 & 2.141 & 2.003 & 2.309 & 2.222 & 2.308 & 184.59 & 19116.1 & 0 & 2.00/2 \\
\code{p1--GP--p1} & triang. no estr. & 0.0 & 2.292 & 2.131 & 2.004 & 2.004 & 2.003 & 2.003 & 2.600 & 2.164 & 2.545 & 2.545 & 2.542 & 2.542 & 799.04 & 19938.2 & 0 & 2.00/3 \\
\code{p1--GP--p1} & triang. no estr. & 0.1 & 2.345 & 2.112 & 2.092 & 2.092 & 2.089 & 2.089 & 2.552 & 2.095 & 2.516 & 2.516 & 2.515 & 2.515 & 852.76 & 21573.2 & 0 & 2.00/3 \\
\hline
\code{p1--GP--p2} & hexa & 0.0 & 1.913 & 2.147 & 1.998 & 2.181 & 1.949 & 2.179 & 2.611 & 2.168 & 2.001 & 2.155 & 2.096 & 2.151 & 300.35 & 18274.8 & 0 & 2.00/3 \\
\code{p1--GP--p2} & hexa & 0.1 & 1.972 & 2.163 & 1.999 & 2.199 & 1.972 & 2.197 & 2.615 & 2.141 & 2.001 & 2.150 & 2.071 & 2.145 & 287.40 & 19229.6 & 0 & 2.00/2 \\
\code{p1--GP--p2} & triang. no estr. & 0.0 & 2.292 & 2.131 & 1.991 & 2.067 & 2.009 & 2.066 & 2.600 & 2.164 & 2.002 & 2.176 & 2.268 & 2.168 & 920.70 & 20295.7 & 0 & 2.00/3 \\
\code{p1--GP--p2} & triang. no estr. & 0.1 & 2.345 & 2.112 & 1.992 & 2.087 & 2.034 & 2.085 & 2.552 & 2.095 & 2.000 & 2.109 & 2.167 & 2.105 & 1012.2 & 21716.7 & 0 & 2.00/3 \\
\hline
\code{p1--GP--p3} & hexa & 0.0 & 1.913 & 2.147 & 1.999 & 2.194 & 1.949 & 2.193 & 2.611 & 2.168 & 2.001 & 2.170 & 2.096 & 2.166 & 315.72 & 18429.8 & 0 & 2.00/3 \\
\code{p1--GP--p3} & hexa & 0.1 & 1.972 & 2.163 & 1.999 & 2.215 & 1.973 & 2.213 & 2.615 & 2.141 & 2.001 & 2.165 & 2.071 & 2.160 & 344.67 & 19322.5 & 0 & 2.00/2 \\
\code{p1--GP--p3} & triang. no estr. & 0.0 & 2.292 & 2.131 & 1.992 & 2.080 & 2.009 & 2.079 & 2.600 & 2.164 & 2.002 & 2.209 & 2.268 & 2.200 & 1028.1 & 20451.8 & 0 & 2.00/3 \\
\code{p1--GP--p3} & triang. no estr. & 0.1 & 2.345 & 2.112 & 1.992 & 2.106 & 2.034 & 2.104 & 2.552 & 2.095 & 2.000 & 2.134 & 2.167 & 2.129 & 1107.5 & 21798.8 & 0 & 2.00/3 \\
\hline
\code{p2--na--NO} & hexa & 0.0 & 2.128 & 2.568 & 2.142 & 2.142 & 2.143 & 2.143 & 2.020 & 2.183 & 2.025 & 2.025 & 2.026 & 2.026 & 154.77 & 18705.6 & 0 & 2.00/3 \\
\code{p2--na--NO} & hexa & 0.1 & 2.129 & 2.608 & 2.143 & 2.143 & 2.144 & 2.144 & 2.022 & 2.213 & 2.027 & 2.027 & 2.028 & 2.028 & 177.02 & 20298.4 & 0 & 2.00/2 \\
\code{p2--na--NO} & triang. no estr. & 0.0 & 2.058 & 2.376 & 2.090 & 2.090 & 2.086 & 2.086 & 1.991 & 2.077 & 1.998 & 1.998 & 2.002 & 2.002 & 838.03 & 21912.5 & 0 & 2.00/3 \\
\code{p2--na--NO} & triang. no estr. & 0.1 & 2.045 & 2.405 & 2.079 & 2.079 & 2.075 & 2.075 & 1.989 & 2.091 & 1.996 & 1.996 & 1.997 & 1.997 & 871.55 & 24391.2 & 0 & 2.00/3 \\
\hline
\code{p2--CCp1--p1} & hexa & 0.0 & 2.150 & 2.609 & 2.164 & 2.356 & 2.166 & 2.334 & 2.023 & 2.214 & 2.029 & 2.293 & 2.030 & 2.232 & 219.58 & 19535.3 & 0 & 2.00/3 \\
\code{p2--CCp1--p1} & hexa & 0.1 & 2.153 & 2.652 & 2.167 & 2.356 & 2.170 & 2.335 & 2.026 & 2.251 & 2.031 & 2.293 & 2.032 & 2.235 & 256.56 & 20889.2 & 0 & 2.00/2 \\
\code{p2--CCp1--p1} & triang. no estr. & 0.0 & 2.058 & 2.376 & 2.090 & 2.090 & 2.086 & 2.086 & 1.991 & 2.077 & 1.998 & 1.998 & 2.002 & 2.002 & 1146.2 & 22805.4 & 0 & 2.00/3 \\
\code{p2--CCp1--p1} & triang. no estr. & 0.1 & 2.045 & 2.405 & 2.079 & 2.079 & 2.075 & 2.075 & 1.989 & 2.091 & 1.996 & 1.996 & 1.997 & 1.997 & 1044.9 & 24997.5 & 0 & 2.00/3 \\
\hline
\code{p2--CCp1--p2} & hexa & 0.0 & 2.213 & 2.709 & 2.224 & 2.230 & 2.235 & 2.188 & 2.031 & 2.308 & 2.039 & 2.151 & 2.041 & 2.094 & 228.91 & 19806.9 & 0 & 2.00/3 \\
\code{p2--CCp1--p2} & hexa & 0.1 & 2.225 & 2.760 & 2.235 & 2.230 & 2.248 & 2.187 & 2.036 & 2.369 & 2.043 & 2.151 & 2.046 & 2.094 & 242.31 & 21005.3 & 0 & 2.00/2 \\
\code{p2--CCp1--p2} & triang. no estr. & 0.0 & 2.121 & 2.535 & 2.161 & 2.102 & 2.166 & 2.040 & 1.990 & 2.162 & 2.003 & 2.058 & 2.012 & 2.018 & 1105.5 & 23121.4 & 0 & 2.00/3 \\
\code{p2--CCp1--p2} & triang. no estr. & 0.1 & 2.108 & 2.582 & 2.154 & 2.101 & 2.160 & 2.039 & 1.987 & 2.199 & 1.998 & 2.058 & 2.003 & 2.018 & 1105.9 & 25192.9 & 0 & 2.00/3 \\
\hline
\code{p2--CCp1--p3} & hexa & 0.0 & 2.213 & 2.709 & 2.224 & 2.248 & 2.235 & 2.214 & 2.031 & 2.308 & 2.039 & 2.168 & 2.041 & 2.113 & 238.91 & 19939.7 & 0 & 2.00/3 \\
\code{p2--CCp1--p3} & hexa & 0.1 & 2.225 & 2.760 & 2.235 & 2.248 & 2.248 & 2.214 & 2.036 & 2.369 & 2.043 & 2.168 & 2.046 & 2.113 & 238.71 & 21069.7 & 0 & 2.00/2 \\
\code{p2--CCp1--p3} & triang. no estr. & 0.0 & 2.122 & 2.535 & 2.161 & 2.127 & 2.166 & 2.050 & 1.990 & 2.162 & 2.003 & 2.074 & 2.012 & 2.023 & 1207.3 & 23264.0 & 0 & 2.00/3 \\
\code{p2--CCp1--p3} & triang. no estr. & 0.1 & 2.108 & 2.582 & 2.154 & 2.126 & 2.160 & 2.050 & 1.987 & 2.199 & 1.998 & 2.074 & 2.003 & 2.023 & 1239.9 & 25275.6 & 0 & 2.00/3 \\
\hline
\code{p2--CCp2--p1} & hexa & 0.0 & 2.150 & 2.609 & 2.166 & 3.039 & 2.166 & 2.253 & 2.023 & 2.214 & 2.029 & 2.986 & 2.030 & 2.044 & 264.32 & 20008.4 & 0 & 2.00/3 \\
\code{p2--CCp2--p1} & hexa & 0.1 & 2.153 & 2.653 & 2.170 & 3.040 & 2.170 & 2.269 & 2.026 & 2.251 & 2.031 & 2.989 & 2.032 & 2.050 & 266.71 & 20936.1 & 0 & 2.00/2 \\
\code{p2--CCp2--p1} & triang. no estr. & 0.0 & 2.058 & 2.376 & 2.090 & 2.090 & 2.086 & 2.086 & 1.991 & 2.077 & 1.998 & 1.998 & 2.002 & 2.002 & 1235.2 & 23870.7 & 0 & 2.00/3 \\
\code{p2--CCp2--p1} & triang. no estr. & 0.1 & 2.045 & 2.405 & 2.079 & 2.079 & 2.075 & 2.075 & 1.989 & 2.091 & 1.996 & 1.996 & 1.997 & 1.997 & 1321.5 & 25229.5 & 0 & 2.00/3 \\
\hline
\code{p2--CCp2--p2} & hexa & 0.0 & 2.214 & 2.709 & 2.236 & 3.063 & 2.235 & 2.434 & 2.032 & 2.308 & 2.041 & 2.999 & 2.041 & 2.086 & 282.98 & 20284.4 & 0 & 2.00/3 \\
\code{p2--CCp2--p2} & hexa & 0.1 & 2.226 & 2.761 & 2.250 & 3.063 & 2.249 & 2.475 & 2.037 & 2.371 & 2.046 & 3.000 & 2.046 & 2.105 & 286.14 & 21071.6 & 0 & 2.00/2 \\
\code{p2--CCp2--p2} & triang. no estr. & 0.0 & 2.122 & 2.536 & 2.172 & 3.012 & 2.166 & 2.317 & 1.990 & 2.162 & 2.003 & 2.991 & 2.012 & 2.038 & 1196.2 & 24206.0 & 0 & 2.00/3 \\
\code{p2--CCp2--p2} & triang. no estr. & 0.1 & 2.108 & 2.583 & 2.166 & 3.012 & 2.160 & 2.339 & 1.987 & 2.200 & 1.999 & 2.993 & 2.002 & 2.039 & 1391.4 & 25433.3 & 0 & 2.00/3 \\
\hline
\code{p2--CCp2--p3} & hexa & 0.0 & 2.214 & 2.709 & 2.236 & 3.058 & 2.235 & 2.434 & 2.032 & 2.308 & 2.041 & 2.998 & 2.041 & 2.088 & 286.51 & 20421.3 & 0 & 2.00/3 \\
\code{p2--CCp2--p3} & hexa & 0.1 & 2.226 & 2.761 & 2.250 & 3.058 & 2.249 & 2.474 & 2.037 & 2.371 & 2.046 & 2.999 & 2.046 & 2.106 & 303.51 & 21165.0 & 0 & 2.00/2 \\
\code{p2--CCp2--p3} & triang. no estr. & 0.0 & 2.122 & 2.536 & 2.172 & 3.009 & 2.166 & 2.314 & 1.990 & 2.162 & 2.003 & 2.991 & 2.012 & 2.039 & 1195.3 & 24362.8 & 0 & 2.00/3 \\
\code{p2--CCp2--p3} & triang. no estr. & 0.1 & 2.108 & 2.583 & 2.166 & 3.010 & 2.160 & 2.336 & 1.987 & 2.200 & 1.999 & 2.993 & 2.002 & 2.040 & 1316.5 & 25555.4 & 0 & 2.00/3 \\
\hline
\code{p2--CCp3--p1} & hexa & 0.0 & 2.150 & 2.609 & 2.166 & 3.135 & 2.166 & 2.183 & 2.023 & 2.214 & 2.029 & 1.992 & 2.030 & 2.029 & 401.90 & 21240.8 & 0 & 2.00/3 \\
\code{p2--CCp3--p1} & hexa & 0.1 & 2.153 & 2.653 & 2.170 & 3.206 & 2.170 & 2.186 & 2.026 & 2.251 & 2.031 & 2.002 & 2.032 & 2.031 & 423.47 & 21652.0 & 0 & 2.00/2 \\
\code{p2--CCp3--p1} & triang. no estr. & 0.0 & 2.058 & 2.376 & 2.090 & 2.090 & 2.086 & 2.086 & 1.991 & 2.077 & 1.998 & 1.998 & 2.002 & 2.002 & 1500.9 & 26335.4 & 0 & 2.00/3 \\
\code{p2--CCp3--p1} & triang. no estr. & 0.1 & 2.045 & 2.405 & 2.079 & 2.079 & 2.075 & 2.075 & 1.989 & 2.091 & 1.996 & 1.996 & 1.997 & 1.997 & 1643.3 & 27242.5 & 0 & 2.00/3 \\
\hline
\code{p2--CCp3--p2} & hexa & 0.0 & 2.214 & 2.709 & 2.236 & 3.796 & 2.235 & 2.299 & 2.032 & 2.308 & 2.041 & 2.178 & 2.041 & 2.039 & 427.34 & 21488.3 & 0 & 2.00/3 \\
\code{p2--CCp3--p2} & hexa & 0.1 & 2.226 & 2.761 & 2.249 & 3.882 & 2.249 & 2.314 & 2.037 & 2.371 & 2.046 & 2.315 & 2.046 & 2.044 & 465.26 & 21933.4 & 0 & 2.00/2 \\
\code{p2--CCp3--p2} & triang. no estr. & 0.0 & 2.122 & 2.536 & 2.172 & 3.302 & 2.166 & 2.231 & 1.990 & 2.162 & 2.003 & 1.988 & 2.012 & 2.010 & 1582.1 & 26653.0 & 0 & 2.00/3 \\
\code{p2--CCp3--p2} & triang. no estr. & 0.1 & 2.108 & 2.583 & 2.166 & 3.390 & 2.160 & 2.225 & 1.987 & 2.200 & 1.999 & 2.010 & 2.002 & 2.001 & 1728.1 & 27559.0 & 0 & 2.00/3 \\
\hline
\code{p2--CCp3--p3} & hexa & 0.0 & 2.214 & 2.709 & 2.236 & 3.789 & 2.235 & 2.298 & 2.032 & 2.308 & 2.041 & 2.133 & 2.041 & 2.039 & 433.79 & 21626.6 & 0 & 2.00/3 \\
\code{p2--CCp3--p3} & hexa & 0.1 & 2.226 & 2.761 & 2.249 & 3.877 & 2.249 & 2.313 & 2.037 & 2.371 & 2.046 & 2.262 & 2.046 & 2.044 & 537.75 & 22086.9 & 0 & 2.00/2 \\
\code{p2--CCp3--p3} & triang. no estr. & 0.0 & 2.122 & 2.536 & 2.172 & 3.267 & 2.166 & 2.229 & 1.990 & 2.162 & 2.003 & 1.957 & 2.012 & 2.010 & 1627.4 & 26819.6 & 0 & 2.00/3 \\
\code{p2--CCp3--p3} & triang. no estr. & 0.1 & 2.108 & 2.583 & 2.166 & 3.356 & 2.160 & 2.223 & 1.987 & 2.200 & 1.999 & 1.972 & 2.002 & 2.001 & 1692.3 & 27730.4 & 0 & 2.00/3 \\
\hline
\code{p2--GP--p1} & hexa & 0.0 & 2.150 & 2.609 & 2.003 & 2.689 & 2.102 & 2.694 & 2.023 & 2.214 & 2.001 & 2.304 & 2.011 & 2.312 & 206.03 & 18811.8 & 0 & 2.00/3 \\
\code{p2--GP--p1} & hexa & 0.1 & 2.153 & 2.653 & 2.003 & 2.728 & 2.100 & 2.732 & 2.026 & 2.251 & 2.001 & 2.350 & 2.011 & 2.358 & 229.20 & 20346.2 & 0 & 2.00/2 \\
\code{p2--GP--p1} & triang. no estr. & 0.0 & 2.058 & 2.376 & 2.090 & 2.090 & 2.086 & 2.086 & 1.991 & 2.077 & 1.998 & 1.998 & 2.002 & 2.002 & 935.71 & 22002.3 & 0 & 2.00/3 \\
\code{p2--GP--p1} & triang. no estr. & 0.1 & 2.045 & 2.405 & 2.079 & 2.079 & 2.075 & 2.075 & 1.989 & 2.091 & 1.996 & 1.996 & 1.997 & 1.997 & 863.88 & 24427.9 & 0 & 2.00/3 \\
\hline
\code{p2--GP--p2} & hexa & 0.0 & 2.214 & 2.709 & 2.003 & 2.924 & 2.114 & 2.923 & 2.032 & 2.308 & 2.000 & 2.644 & 2.007 & 2.650 & 290.14 & 19087.4 & 0 & 2.00/3 \\
\code{p2--GP--p2} & hexa & 0.1 & 2.226 & 2.761 & 2.003 & 2.952 & 2.112 & 2.951 & 2.037 & 2.371 & 2.000 & 2.706 & 2.007 & 2.711 & 316.92 & 20450.9 & 0 & 2.00/2 \\
\code{p2--GP--p2} & triang. no estr. & 0.0 & 2.122 & 2.536 & 1.992 & 2.805 & 2.037 & 2.807 & 1.990 & 2.162 & 1.998 & 2.489 & 2.003 & 2.499 & 1036.4 & 22340.1 & 0 & 2.00/3 \\
\code{p2--GP--p2} & triang. no estr. & 0.1 & 2.108 & 2.583 & 1.992 & 2.840 & 2.035 & 2.842 & 1.987 & 2.200 & 1.998 & 2.553 & 2.002 & 2.561 & 987.41 & 24553.8 & 0 & 2.00/3 \\
\hline
\code{p2--GP--p3} & hexa & 0.0 & 2.214 & 2.709 & 2.004 & 2.923 & 2.114 & 2.922 & 2.032 & 2.308 & 2.000 & 2.650 & 2.007 & 2.656 & 346.39 & 19235.2 & 0 & 2.00/3 \\
\code{p2--GP--p3} & hexa & 0.1 & 2.226 & 2.761 & 2.004 & 2.950 & 2.113 & 2.950 & 2.037 & 2.371 & 2.000 & 2.711 & 2.007 & 2.716 & 366.70 & 20522.7 & 0 & 2.00/2 \\
\code{p2--GP--p3} & triang. no estr. & 0.0 & 2.122 & 2.536 & 1.993 & 2.809 & 2.037 & 2.812 & 1.990 & 2.162 & 1.998 & 2.500 & 2.003 & 2.510 & 1128.1 & 22514.3 & 0 & 2.00/3 \\
\code{p2--GP--p3} & triang. no estr. & 0.1 & 2.108 & 2.583 & 1.993 & 2.843 & 2.035 & 2.846 & 1.987 & 2.200 & 1.998 & 2.564 & 2.002 & 2.572 & 1084.0 & 24631.3 & 0 & 2.00/3 \\
\hline
\code{p3--na--NO} & hexa & 0.0 & 2.125 & 2.485 & 2.132 & 2.132 & 2.135 & 2.135 & 2.013 & 2.050 & 2.025 & 2.025 & 2.024 & 2.024 & 192.01 & 21809.6 & 0 & 2.00/3 \\
\code{p3--na--NO} & hexa & 0.1 & 2.108 & 2.481 & 2.114 & 2.114 & 2.117 & 2.117 & 2.010 & 2.048 & 2.020 & 2.020 & 2.020 & 2.020 & 320.68 & 24530.8 & 0 & 2.00/2 \\
\code{p3--na--NO} & triang. no estr. & 0.0 & 2.041 & 2.209 & 2.061 & 2.061 & 2.060 & 2.060 & 1.999 & 2.015 & 2.008 & 2.008 & 2.008 & 2.008 & 968.40 & 27545.1 & 0 & 2.00/3 \\
\code{p3--na--NO} & triang. no estr. & 0.1 & 2.029 & 2.204 & 2.046 & 2.046 & 2.046 & 2.046 & 1.999 & 2.014 & 2.008 & 2.008 & 2.007 & 2.007 & 884.29 & 31519.3 & 0 & 2.00/3 \\
\hline
\code{p3--CCp1--p1} & hexa & 0.0 & 2.185 & 2.634 & 2.188 & 2.354 & 2.201 & 2.339 & 2.020 & 2.072 & 2.038 & 2.290 & 2.040 & 2.248 & 305.03 & 23002.7 & 0 & 2.00/3 \\
\code{p3--CCp1--p1} & hexa & 0.1 & 2.160 & 2.629 & 2.164 & 2.354 & 2.176 & 2.339 & 2.016 & 2.068 & 2.031 & 2.290 & 2.032 & 2.248 & 392.23 & 25270.6 & 0 & 2.00/2 \\
\code{p3--CCp1--p1} & triang. no estr. & 0.0 & 2.041 & 2.209 & 2.061 & 2.061 & 2.060 & 2.060 & 1.999 & 2.015 & 2.008 & 2.008 & 2.008 & 2.008 & 1291.4 & 28435.5 & 0 & 2.00/3 \\
\code{p3--CCp1--p1} & triang. no estr. & 0.1 & 2.029 & 2.204 & 2.046 & 2.046 & 2.046 & 2.046 & 1.999 & 2.014 & 2.008 & 2.008 & 2.007 & 2.007 & 1154.7 & 31984.8 & 0 & 2.00/3 \\
\hline
\code{p3--CCp1--p2} & hexa & 0.0 & 2.620 & 3.309 & 2.640 & 2.229 & 2.603 & 2.184 & 2.117 & 2.267 & 2.264 & 2.150 & 2.211 & 2.094 & 324.98 & 23127.6 & 0 & 2.00/3 \\
\code{p3--CCp1--p2} & hexa & 0.1 & 2.564 & 3.307 & 2.592 & 2.229 & 2.554 & 2.184 & 2.084 & 2.237 & 2.206 & 2.150 & 2.160 & 2.094 & 428.72 & 25401.0 & 0 & 2.00/2 \\
\code{p3--CCp1--p2} & triang. no estr. & 0.0 & 2.396 & 2.791 & 2.510 & 2.101 & 2.460 & 2.038 & 2.003 & 2.026 & 2.041 & 2.058 & 2.014 & 2.017 & 1333.6 & 28660.7 & 0 & 2.00/3 \\
\code{p3--CCp1--p2} & triang. no estr. & 0.1 & 2.330 & 2.777 & 2.443 & 2.101 & 2.397 & 2.038 & 1.999 & 2.022 & 2.031 & 2.058 & 2.006 & 2.017 & 1166.1 & 32124.5 & 0 & 2.00/3 \\
\hline
\code{p3--CCp1--p3} & hexa & 0.0 & 2.620 & 3.309 & 2.640 & 2.247 & 2.603 & 2.210 & 2.117 & 2.267 & 2.264 & 2.167 & 2.211 & 2.112 & 333.14 & 23206.1 & 0 & 2.00/3 \\
\code{p3--CCp1--p3} & hexa & 0.1 & 2.564 & 3.307 & 2.592 & 2.247 & 2.554 & 2.210 & 2.084 & 2.237 & 2.206 & 2.167 & 2.160 & 2.112 & 469.38 & 25461.2 & 0 & 2.00/2 \\
\code{p3--CCp1--p3} & triang. no estr. & 0.0 & 2.396 & 2.791 & 2.509 & 2.126 & 2.460 & 2.048 & 2.003 & 2.026 & 2.041 & 2.074 & 2.014 & 2.023 & 1348.9 & 28769.0 & 0 & 2.00/3 \\
\code{p3--CCp1--p3} & triang. no estr. & 0.1 & 2.330 & 2.777 & 2.443 & 2.126 & 2.397 & 2.048 & 2.000 & 2.022 & 2.031 & 2.074 & 2.006 & 2.023 & 1202.7 & 32179.6 & 0 & 2.00/3 \\
\hline
\code{p3--CCp2--p1} & hexa & 0.0 & 2.187 & 2.637 & 2.200 & 3.042 & 2.203 & 2.642 & 2.020 & 2.073 & 2.041 & 2.999 & 2.040 & 2.257 & 391.36 & 23106.3 & 0 & 2.00/3 \\
\code{p3--CCp2--p1} & hexa & 0.1 & 2.162 & 2.631 & 2.175 & 3.042 & 2.177 & 2.638 & 2.016 & 2.069 & 2.033 & 2.999 & 2.033 & 2.252 & 521.75 & 25281.9 & 0 & 2.00/2 \\
\code{p3--CCp2--p1} & triang. no estr. & 0.0 & 2.041 & 2.209 & 2.061 & 2.061 & 2.060 & 2.060 & 1.999 & 2.015 & 2.008 & 2.008 & 2.008 & 2.008 & 1547.0 & 28727.0 & 0 & 2.00/3 \\
\code{p3--CCp2--p1} & triang. no estr. & 0.1 & 2.029 & 2.204 & 2.046 & 2.046 & 2.046 & 2.046 & 1.999 & 2.014 & 2.008 & 2.008 & 2.007 & 2.007 & 1854.1 & 32015.7 & 0 & 2.00/3 \\
\hline
\code{p3--CCp2--p2} & hexa & 0.0 & 3.539 & 3.992 & 3.284 & 3.063 & 3.303 & 3.110 & 4.324 & 4.279 & 4.079 & 3.002 & 4.086 & 3.055 & 396.06 & 23282.3 & 0 & 2.00/3 \\
\code{p3--CCp2--p2} & hexa & 0.1 & 3.570 & 4.025 & 3.311 & 3.063 & 3.331 & 3.109 & 4.299 & 4.272 & 4.076 & 3.002 & 4.084 & 3.049 & 561.58 & 25389.0 & 0 & 2.00/2 \\
\code{p3--CCp2--p2} & triang. no estr. & 0.0 & 4.057 & 4.096 & 3.842 & 3.013 & 3.861 & 3.085 & 4.476 & 4.164 & 4.251 & 2.997 & 4.258 & 3.023 & 1681.5 & 29025.2 & 0 & 2.00/3 \\
\code{p3--CCp2--p2} & triang. no estr. & 0.1 & 4.077 & 4.105 & 3.880 & 3.013 & 3.901 & 3.082 & 4.460 & 4.139 & 4.199 & 2.997 & 4.214 & 3.022 & 1659.8 & 32176.2 & 0 & 2.00/3 \\
\hline
\code{p3--CCp2--p3} & hexa & 0.0 & 3.539 & 3.992 & 3.284 & 3.058 & 3.303 & 3.102 & 4.324 & 4.279 & 4.079 & 3.001 & 4.086 & 3.050 & 424.70 & 23377.0 & 0 & 2.00/3 \\
\code{p3--CCp2--p3} & hexa & 0.1 & 3.570 & 4.025 & 3.311 & 3.058 & 3.331 & 3.101 & 4.299 & 4.272 & 4.076 & 3.001 & 4.084 & 3.044 & 548.85 & 25443.2 & 0 & 2.00/2 \\
\code{p3--CCp2--p3} & triang. no estr. & 0.0 & 4.057 & 4.096 & 3.842 & 3.011 & 3.861 & 3.072 & 4.477 & 4.165 & 4.252 & 2.997 & 4.259 & 3.019 & 1625.8 & 29189.0 & 0 & 2.00/3 \\
\code{p3--CCp2--p3} & triang. no estr. & 0.1 & 4.076 & 4.105 & 3.880 & 3.011 & 3.901 & 3.070 & 4.461 & 4.139 & 4.199 & 2.997 & 4.214 & 3.017 & 1579.9 & 32250.7 & 0 & 2.00/3 \\
\hline
\code{p3--CCp3--p1} & hexa & 0.0 & 2.186 & 2.637 & 2.200 & 3.993 & 2.202 & 2.225 & 2.020 & 2.073 & 2.041 & 2.812 & 2.040 & 2.040 & 526.58 & 23777.7 & 0 & 2.00/3 \\
\code{p3--CCp3--p1} & hexa & 0.1 & 2.161 & 2.631 & 2.174 & 3.993 & 2.177 & 2.201 & 2.016 & 2.069 & 2.033 & 2.809 & 2.033 & 2.032 & 653.85 & 25298.9 & 0 & 2.00/2 \\
\code{p3--CCp3--p1} & triang. no estr. & 0.0 & 2.041 & 2.209 & 2.061 & 2.061 & 2.060 & 2.060 & 1.999 & 2.015 & 2.008 & 2.008 & 2.008 & 2.008 & 1660.4 & 31065.9 & 0 & 2.00/3 \\
\code{p3--CCp3--p1} & triang. no estr. & 0.1 & 2.029 & 2.204 & 2.046 & 2.046 & 2.046 & 2.046 & 1.999 & 2.014 & 2.008 & 2.008 & 2.007 & 2.007 & 1871.3 & 32210.7 & 0 & 2.00/3 \\
\hline
\code{p3--CCp3--p2} & hexa & 0.0 & 3.538 & 3.993 & 3.290 & 4.231 & 3.304 & 3.441 & 4.324 & 4.279 & 4.081 & 4.284 & 4.087 & 4.101 & 561.34 & 24015.9 & 0 & 2.00/3 \\
\code{p3--CCp3--p2} & hexa & 0.1 & 3.569 & 4.026 & 3.319 & 4.231 & 3.333 & 3.485 & 4.299 & 4.272 & 4.078 & 4.283 & 4.084 & 4.097 & 661.30 & 25419.6 & 0 & 2.00/2 \\
\code{p3--CCp3--p2} & triang. no estr. & 0.0 & 4.058 & 4.096 & 3.846 & 4.109 & 3.863 & 3.934 & 4.476 & 4.164 & 4.253 & 4.049 & 4.258 & 4.241 & 1691.8 & 31398.8 & 0 & 2.00/3 \\
\code{p3--CCp3--p2} & triang. no estr. & 0.1 & 4.078 & 4.105 & 3.884 & 4.109 & 3.903 & 3.988 & 4.461 & 4.139 & 4.202 & 4.049 & 4.214 & 4.191 & 1831.7 & 32387.7 & 0 & 2.00/3 \\
\hline
\code{p3--CCp3--p3} & hexa & 0.0 & 3.538 & 3.993 & 3.290 & 4.241 & 3.304 & 3.438 & 4.324 & 4.279 & 4.081 & 4.302 & 4.087 & 4.101 & 533.17 & 24150.7 & 0 & 2.00/3 \\
\code{p3--CCp3--p3} & hexa & 0.1 & 3.569 & 4.026 & 3.319 & 4.241 & 3.333 & 3.481 & 4.299 & 4.272 & 4.078 & 4.300 & 4.084 & 4.097 & 510.03 & 25470.3 & 0 & 2.00/2 \\
\code{p3--CCp3--p3} & triang. no estr. & 0.0 & 4.057 & 4.096 & 3.846 & 4.114 & 3.862 & 3.932 & 4.477 & 4.165 & 4.254 & 4.050 & 4.259 & 4.242 & 1732.8 & 31562.0 & 0 & 2.00/3 \\
\code{p3--CCp3--p3} & triang. no estr. & 0.1 & 4.077 & 4.105 & 3.884 & 4.114 & 3.903 & 3.986 & 4.461 & 4.139 & 4.202 & 4.050 & 4.214 & 4.192 & 1720.3 & 32487.7 & 0 & 2.00/3 \\
\hline
\code{p3--GP--p1} & hexa & 0.0 & 2.187 & 2.637 & 2.001 & 2.598 & 1.989 & 2.596 & 2.020 & 2.073 & 2.000 & 2.068 & 1.999 & 2.066 & 290.60 & 21859.5 & 0 & 2.00/3 \\
\code{p3--GP--p1} & hexa & 0.1 & 2.161 & 2.631 & 2.001 & 2.592 & 1.988 & 2.589 & 2.016 & 2.069 & 2.000 & 2.060 & 1.999 & 2.059 & 432.06 & 24569.5 & 0 & 2.00/2 \\
\code{p3--GP--p1} & triang. no estr. & 0.0 & 2.041 & 2.209 & 2.061 & 2.061 & 2.060 & 2.060 & 1.999 & 2.015 & 2.008 & 2.008 & 2.008 & 2.008 & 1230.9 & 27569.1 & 0 & 2.00/3 \\
\code{p3--GP--p1} & triang. no estr. & 0.1 & 2.029 & 2.204 & 2.046 & 2.046 & 2.046 & 2.046 & 1.999 & 2.014 & 2.008 & 2.008 & 2.007 & 2.007 & 1092.4 & 31558.3 & 0 & 2.00/3 \\
\hline
\code{p3--GP--p2} & hexa & 0.0 & 3.538 & 3.992 & 1.999 & 4.003 & 1.997 & 4.010 & 4.324 & 4.279 & 2.000 & 4.194 & 2.001 & 4.196 & 381.25 & 21971.2 & 0 & 2.00/3 \\
\code{p3--GP--p2} & hexa & 0.1 & 3.569 & 4.026 & 1.999 & 4.033 & 1.997 & 4.040 & 4.299 & 4.272 & 2.000 & 4.212 & 2.000 & 4.213 & 511.96 & 24681.8 & 0 & 2.00/2 \\
\code{p3--GP--p2} & triang. no estr. & 0.0 & 4.057 & 4.096 & 1.992 & 4.077 & 2.002 & 4.083 & 4.476 & 4.164 & 1.998 & 4.204 & 2.003 & 4.204 & 1724.5 & 27714.9 & 0 & 2.00/3 \\
\code{p3--GP--p2} & triang. no estr. & 0.1 & 4.077 & 4.105 & 1.992 & 4.094 & 2.001 & 4.098 & 4.461 & 4.139 & 1.998 & 4.195 & 2.003 & 4.198 & 1145.9 & 31669.6 & 0 & 2.00/3 \\
\hline
\code{p3--GP--p3} & hexa & 0.0 & 3.538 & 3.992 & 1.999 & 4.000 & 1.998 & 4.008 & 4.324 & 4.279 & 2.000 & 4.198 & 2.001 & 4.200 & 440.53 & 22050.9 & 0 & 2.00/3 \\
\code{p3--GP--p3} & hexa & 0.1 & 3.569 & 4.026 & 1.999 & 4.031 & 1.997 & 4.039 & 4.299 & 4.272 & 2.000 & 4.217 & 2.001 & 4.219 & 472.62 & 24729.1 & 0 & 2.00/2 \\
\code{p3--GP--p3} & triang. no estr. & 0.0 & 4.057 & 4.096 & 1.992 & 4.079 & 2.002 & 4.086 & 4.477 & 4.165 & 1.998 & 4.227 & 2.003 & 4.229 & 1466.0 & 27775.6 & 0 & 2.00/3 \\
\code{p3--GP--p3} & triang. no estr. & 0.1 & 4.077 & 4.105 & 1.992 & 4.098 & 2.002 & 4.103 & 4.461 & 4.139 & 1.998 & 4.219 & 2.003 & 4.222 & 1625.2 & 31740.2 & 0 & 2.00/3 \\
\bottomrule
\end{longtable}
\normalsize


